\chapter{The theory of middle convolution}
\label{chap:2}
\setcounter{section}{-1}

\section{Transition from irreducible local systems on open sets of
\texorpdfstring{\(\mathbb{P}^1\)}{P1} to irreducible middle extension sheaves
on \texorpdfstring{\(\mathbb{A}^1\)}{A1}}
\label{sec:2.0}

\sourcepara{2.0.1}
Let \(S\) be a nonempty finite subset of \(\mathbb{P}^1(\mathbb{C})\), and
\(U:=\mathbb{P}^1-S\) the open complement,
\(j:U^{\mathrm{an}}\to(\mathbb{P}^1)^{\mathrm{an}}\) the inclusion,
\(\mathcal{F}\) an irreducible \(\mathbb{C}\)-local system on
\(U^{\mathrm{an}}\) of rank \(n\geq 1\). We know by
\hyperref[thm:1.1.2]{Theorem~1.1.2} that \(\mathcal{F}\) is physically rigid
if and only if
\[
  \chi\bigl((\mathbb{P}^1)^{\mathrm{an}},
  j_*\operatorname{End}(\mathcal{F})\bigr)=2.
\]
Motivated by this fact, we define the \emph{index of rigidity} of
\(\mathcal{F}\) on \(U\), noted \(\operatorname{rig}(\mathcal{F},U)\), to be
the integer
\[
  \operatorname{rig}(\mathcal{F},U):=
  \chi\bigl((\mathbb{P}^1)^{\mathrm{an}},
  j_*\operatorname{End}(\mathcal{F})\bigr).
\]

\setcounter{theorem}{1}
\begin{lemma}
\label{lem:2.0.2}
Let \(S\) be a nonempty finite subset of \(\mathbb{P}^1(\mathbb{C})\),
\(U:=\mathbb{P}^1-S\) the open complement,
\(j:U^{\mathrm{an}}\to(\mathbb{P}^1)^{\mathrm{an}}\) the inclusion,
\(\mathcal{F}\) an irreducible \(\mathbb{C}\)-local system on
\(U^{\mathrm{an}}\) of rank \(n\geq 1\).
\begin{enumerate}[label=(\alph*)]
\item If \(\mathcal{L}\) is any rank one local system on \(U^{\mathrm{an}}\),
we have
\[
  \operatorname{rig}(\mathcal{F}\otimes\mathcal{L},U)
  =\operatorname{rig}(\mathcal{F},U).
\]
\item If \(T\) is any finite subset of \(U^{\mathrm{an}}(\mathbb{C})\), and
\(k:U^{\mathrm{an}}-T\to U^{\mathrm{an}}\) the inclusion, we have
\[
  \operatorname{rig}(k^*\mathcal{F},U-T)
  =\operatorname{rig}(\mathcal{F},U).
\]
\end{enumerate}
\end{lemma}

\begin{proof}
(1) holds because the local system \(\operatorname{End}(\mathcal{F})\) on
\(U^{\mathrm{an}}\) does not change if we replace \(\mathcal{F}\) by
\(\mathcal{F}\otimes\mathcal{L}\). (2) holds because for any local system
\(\mathcal{G}\) on \(U^{\mathrm{an}}\), we have
\(\mathcal{G}\cong k_*k^*\mathcal{G}\). Applying this to
\(\operatorname{End}(\mathcal{F})\), we get
\[
  j_*\operatorname{End}(\mathcal{F})
  =j_*k_*k^*\operatorname{End}(\mathcal{F})
  =j_*k_*\operatorname{End}(k^*\mathcal{F}),
\]
and applying \(\chi((\mathbb{P}^1)^{\mathrm{an}},\_)\) gives the assertion.
\end{proof}

\sourcepara{2.0.3}
Recall \cite[7.3.1]{Ka-ESDE} that on a connected smooth curve
\(U/\mathbb{C}\), an algebraically constructible sheaf \(\mathcal{F}\) of
\(\mathbb{C}\)-vector spaces on \(U^{\mathrm{an}}\) is called a
\emph{middle extension sheaf} if for some (or equivalently for every) nonempty
Zariski open set \(k:V\to U\) such that \(k^*\mathcal{F}\) is a local system on
\(V^{\mathrm{an}}\), we have \(\mathcal{F}\cong k_*k^*\mathcal{F}\). A middle
extension sheaf \(\mathcal{F}\) on \(U\) is called an \emph{irreducible middle
extension} if for some (or equivalently for every) nonempty Zariski open set
\(k:V\to U\) such that \(k^*\mathcal{F}\) is a local system,
\(k^*\mathcal{F}\) is an irreducible local system.

\sourcepara{2.0.4}
We now specialize to the case when \(U\) is a Zariski open set of
\(\mathbb{P}^1\), \(j:U\to\mathbb{P}^1\) its inclusion. We will now define the
index of rigidity \(\operatorname{rig}_U(\mathcal{F})\) of an irreducible
middle extension \(\mathcal{F}\) on \(U^{\mathrm{an}}\). Given such an
\(\mathcal{F}\), pick a nonempty Zariski open set \(k:V\to U\) such that
\(k^*\mathcal{F}\) is an irreducible local system. The integer
\[
  \operatorname{rig}(k^*\mathcal{F},V):=
  \chi\bigl((\mathbb{P}^1)^{\mathrm{an}},
  j_*k_*\operatorname{End}(k^*\mathcal{F})\bigr)
\]
is independent of the auxiliary choice of \(V\), thanks to
\hyperref[lem:2.0.2]{Lemma~2.0.2} above. We call it
\(\operatorname{rig}_U(\mathcal{F})\).

\sourcepara{2.0.5}
The situation now is this. For any nonempty Zariski open set \(U\) in
\(\mathbb{P}^1\), we have the category \(\operatorname{IrrME}(U)\) of
irreducible middle extensions on \(U^{\mathrm{an}}\). Whenever \(V\subset U\),
with inclusion \(k:V\to U\), the functors
\[
  k^*:\operatorname{IrrME}(U)\to\operatorname{IrrME}(V)
  \quad\text{and}\quad
  k_*:\operatorname{IrrME}(V)\to\operatorname{IrrME}(U)
\]
are inverse equivalences. Given two nonempty Zariski open sets \(U_1\) and
\(U_2\), pick any nonempty Zariski open \(V\) in \(U_1\cap U_2\), with
inclusions \(k_i:V\to U_i\). Then we get an equivalence
\[
  \varphi_{1,2}:=(k_2)_*(k_1)^*:
  \operatorname{IrrME}(U_1)\to\operatorname{IrrME}(U_2)
\]
which is independent of the auxiliary choice of \(V\). Given three nonempty
Zariski open sets \(U_1,U_2,U_3\), we have
\[
  \varphi_{1,3}=\varphi_{2,3}\circ\varphi_{1,2}.
\]
So we can canonically identify all of these categories.

\sourcepara{2.0.6}
The \(\mathbb{Z}\)-valued functions
\(\mathcal{F}\mapsto\operatorname{rig}_U(\mathcal{F})\) on these categories
respect these identifications, so we may speak of the single function
\[
  \mathcal{F}\mapsto\operatorname{rig}(\mathcal{F}).
\]

\sourcepara{2.0.7}
At first glance, it would seem most natural to work with the single category
\(\operatorname{IrrME}(\mathbb{P}^1)\). However, it turns out to be better to
pick two points in \(\mathbb{P}^1(\mathbb{C})\), label them \(\infty\) and
\(0\), and work on the open set
\(\mathbb{A}^1:=\mathbb{P}^1-\{\infty\}\). Because we have specified the
origin \(0\), this \(\mathbb{A}^1\) has an additive group structure. By
embedding the category \(\operatorname{IrrME}(\mathbb{A}^1)\) in the slightly
larger category \(\operatorname{IrrPerv}(\mathbb{A}^1)\) of irreducible
perverse sheaves on \(\mathbb{A}^1\), we can bring to bear the whole mechanism
of additive convolution.

\section[Transition from middle extensions to perverse sheaves on
\texorpdfstring{\(\mathbb{A}^1\)}{A1}]{Transition from irreducible middle
extension sheaves on
\texorpdfstring{\(\mathbb{A}^1\)}{A1} to irreducible perverse sheaves on
\texorpdfstring{\(\mathbb{A}^1\)}{A1}}
\label{sec:2.1}

\sourcepara{2.1.1}
On any separated \(\mathbb{C}\)-scheme \(X\) of finite type, a sheaf
\(\mathcal{F}\) of \(\mathbb{C}\)-vector spaces on \(X^{\mathrm{an}}\) is said
to be algebraically constructible if there exists a finite partition
\[
  X^{\mathrm{red}}=\coprod_i Y_i
\]
as the disjoint union of smooth connected subschemes \(Y_i\), such that on each
\(Y_i\), \(\mathcal{F}|Y_i^{\mathrm{an}}\) is a local system on
\(Y_i^{\mathrm{an}}\). We denote by
\(\mathrm{D}(X^{\mathrm{an}},\mathbb{C})\) the derived category of the category
of all \(\mathbb{C}\)-sheaves on \(X^{\mathrm{an}}\), and by
\(\mathrm{D}^b_c(X^{\mathrm{an}},\mathbb{C})\) the full subcategory of
\(\mathrm{D}(X^{\mathrm{an}},\mathbb{C})\) consisting of those objects \(K\)
for which
\begin{enumerate}[label=(\alph*)]
\item each cohomology sheaf \(\mathcal{H}^i(K)\) is algebraically
constructible,
\item only finitely many of the sheaves \(\mathcal{H}^i(K)\) are nonzero.
\end{enumerate}
These \(\mathrm{D}^b_c\) support the full Grothendieck formalism of the ``six
operations''. In this formalism, the dualizing complex \(K_{X/\mathbb{C}}\) is
defined as \(f^!\mathbb{C}\), \(f:X\to\operatorname{Spec}(\mathbb{C})\)
denoting the structural morphism.

\sourcepara{2.1.2}
Recall \cite[Ch.~4]{BBD} that an object \(K\) of
\(\mathrm{D}^b_c(X^{\mathrm{an}},\mathbb{C})\) is called
\emph{semiperverse} if its cohomology sheaves \(\mathcal{H}^i(K)\) satisfy
\[
  \dim\operatorname{Supp}(\mathcal{H}^i(K))\leq -i,
  \quad\text{for every integer }i.
\]
An object \(K\) is called \emph{perverse} if both \(K\) and its Verdier dual
\[
  \mathbb{D}_{X/\mathbb{C}}K:=\mathrm{R}\mathcal{H}om(K,f^!\mathbb{C}),
  \quad f:X\to\operatorname{Spec}(\mathbb{C})
\]
denoting the structural morphism, are semiperverse. The main facts about
perversity, semiperversity and duality we will use are the following
\cite[Ch.~4]{BBD}:
\begin{enumerate}[label=(\arabic*)]
\item if \(f:X\to Y\) is an affine morphism, then \(K\mapsto \mathrm{R}f_*K\) preserves
semiperversity;
\item if \(f:X\to Y\) is a quasifinite morphism, then \(K\mapsto \mathrm{R}f_!K\)
preserves semiperversity;
\item if \(f:X\to Y\) is an arbitrary morphism whose geometric fibres all have
dimension \(\leq d\), then \(L\mapsto f^*L[d]\) preserves semiperversity;
\item duality interchanges \(\mathrm{R}f_!\) and \(\mathrm{R}f_*\);
\item duality interchanges \(f^!\) and \(f^*\);
\item if \(f:X\to Y\) is a smooth morphism everywhere of relative dimension
\(d\), then \(f^!=f^*[2d](d)\). Consequently \(f^*[d](d/2)\) is self-dual, and
\(K\mapsto f^*K[d]\) preserves perversity;
\item if \(X\) is smooth over \(\mathbb{C}\), purely of dimension \(d\), then
for any local system \(\mathcal{F}\) on \(X\), \(\mathcal{F}[d]\) is perverse,
and
\[
  \mathbb{D}_{X/k}(\mathcal{F}[d])=\mathcal{F}^{\vee}[d](d).
\]
\end{enumerate}

\sourcepara{2.1.3}
In this discussion, the field \(\mathbb{C}\) occurs in two ways, as the ground
field over which our variable scheme \(X\) is given, and as the coefficient
field. Since we speak of \(X^{\mathrm{an}}\), the field \(\mathbb{C}\) with
its classical topology is being used as the ground field. But the coefficient
field \(\mathbb{C}\) enters in a purely algebraic way, and could be replaced by
any field to which it is isomorphic, for instance by
\(\overline{\mathbb{Q}}_\ell\) (if we grant the axiom of choice). So we might
just as well work with
\(\mathrm{D}^b_c(X^{\mathrm{an}},\overline{\mathbb{Q}}_\ell)\) whenever it is
convenient.

\section{Review of
\texorpdfstring{\(\mathrm{D}^b_c(X,\overline{\mathbb{Q}}_\ell)\)}{Dbc(X,Qlbar)}}
\label{sec:2.2}

\sourcepara{2.2.1}
Let \(k\) be a perfect field of characteristic \(p\neq\ell\). For variable
separated \(k\)-schemes of finite type \(X/k\), we can speak of
\(\mathrm{D}^b_c(X,\overline{\mathbb{Q}}_\ell)\). For morphisms
\(f:X\to Y\) between separated \(k\)-schemes of finite type, one knows
(cf. \cite{De-WeilII} for the case when \(k\) is either algebraically closed or
finite, \cite{Ek}, \cite{Ka-Lau}, \cite[SGA 4, XVIII, 3]{SGA}) that these
\(\mathrm{D}^b_c\) support the full Grothendieck formalism of the ``six
operations''. In this formalism, the (relative to \(k\)) dualizing complex
\(K_X\) in \(\mathrm{D}^b_c(X,\overline{\mathbb{Q}}_\ell)\) is defined as
\(\pi^!\overline{\mathbb{Q}}_\ell\), where \(\pi\) denotes the structural
morphism \(\pi:X\to\operatorname{Spec}(k)\). In terms of \(K_X\), the Verdier
dual \(\mathbb{D}(L)\) of an object \(L\) of
\(\mathrm{D}^b_c(X,\overline{\mathbb{Q}}_\ell)\) is defined as
\[
  \mathrm{R}\mathcal{H}om(L,K_X).
\]
One knows that \(L\cong\mathbb{D}\mathbb{D}(L)\) by the natural map. The
duality theorem asserts that for \(f:X\to Y\) a morphism of finite type between
separated \(k\)-schemes of finite type, one has
\[
  \mathbb{D}(\mathrm{R}f_!L)\cong \mathrm{R}f_*\mathbb{D}(L),\qquad
  \mathbb{D}(\mathrm{R}f_*L)\cong \mathrm{R}f_!\mathbb{D}(L).
\]
If \(X/k\) is a smooth separated \(k\)-scheme of finite type and everywhere of
the same relative dimension, noted \(\dim X\), then
\[
  K_X=\overline{\mathbb{Q}}_\ell[2\dim X](\dim X),
\]
and so
\[
  \mathbb{D}(L)=
  \mathrm{R}\mathcal{H}om(L,\overline{\mathbb{Q}}_\ell)[2\dim X](\dim X).
\]

\sourcepara{2.2.2}
Given two separated \(k\)-schemes \(X/k\) and \(Y/k\) of finite type,
``external tensor product over \(\overline{\mathbb{Q}}_\ell\)'' defines a
bi-exact bilinear pairing
\[
\begin{aligned}
  \mathrm{D}^b_c(X,\overline{\mathbb{Q}}_\ell)
  \times
  \mathrm{D}^b_c(Y,\overline{\mathbb{Q}}_\ell)
  &\to
  \mathrm{D}^b_c(X\times_kY,\overline{\mathbb{Q}}_\ell),\\
  (K,L)\mapsto K\times L:=\operatorname{pr}_1^*K\otimes\operatorname{pr}_2^*L.
\end{aligned}
\]
One knows that \(\mathbb{D}(K\times L)=\mathbb{D}(K)\times\mathbb{D}(L)\).

\sourcepara{2.2.3}
If \(k\) happens to be \(\mathbb{C}\), the comparison theorem gives an exact,
fully faithful ``passage to the analytic'' functor
\[
  \mathrm{D}^b_c(X,\overline{\mathbb{Q}}_\ell)\to
  \mathrm{D}^b_c(X^{\mathrm{an}},\overline{\mathbb{Q}}_\ell)
\]
which is not, however, an equivalence of categories. Everything recalled above
about \(\mathrm{D}^b_c(X,\overline{\mathbb{Q}}_\ell)\) is true also of
\(\mathrm{D}^b_c(X^{\mathrm{an}},\overline{\mathbb{Q}}_\ell)\), and all
cohomological constructions above commute with the ``passage to the analytic''
functor. Given any object in
\(\mathrm{D}^b_c(X^{\mathrm{an}},\mathbb{C})\), for all \(\ell\gg 0\) there
exists an isomorphism
\(\mathbb{C}\cong\overline{\mathbb{Q}}_\ell\) such that the corresponding
object in
\(\mathrm{D}^b_c(X^{\mathrm{an}},\overline{\mathbb{Q}}_\ell)\) lies in (the
essential image of)
\(\mathrm{D}^b_c(X,\overline{\mathbb{Q}}_\ell)\), cf. \cite[6.1]{BBD}. In
\hyperref[prop:5.9.2]{Proposition~5.9.2}, we will give the elementary and down to earth proof
of this fact for local systems on open sets of \(\mathbb{A}^1\), which will
suffice for our purposes.

\section{Review of perverse sheaves}
\label{sec:2.3}

\sourcepara{2.3.1}
We continue to work with \(X/k\) as in \hyperref[par:2.2.1]{(2.2.1)}. An object \(K\) of
\(\mathrm{D}^b_c(X,\overline{\mathbb{Q}}_\ell)\) is called
\emph{semiperverse} if its cohomology sheaves \(\mathcal{H}^iK\) satisfy
\[
  \dim\operatorname{Supp}(\mathcal{H}^iK)\leq -i.
\]
An object \(K\) of \(\mathrm{D}^b_c(X,\overline{\mathbb{Q}}_\ell)\) is called
\emph{perverse} if both \(K\) and its dual \(\mathbb{D}(K)\) are semiperverse.
If \(f:X\to Y\) is an affine (respectively a quasifinite) morphism, then
\(\mathrm{R}f_*\) (respectively \(f_!=\mathrm{R}f_!\)) preserves
semiperversity. So if \(f\) is
both affine and quasifinite (e.g., finite, or an affine immersion), then by
duality both \(f_!=\mathrm{R}f_!\) and \(\mathrm{R}f_*\) preserve perversity.
If \(f:X\to Y\) is
a smooth morphism everywhere of relative dimension \(d\), then \(f^*[d]\)
preserves perversity. In particular, if \(K\) is perverse on \(X\), then its
inverse image on \(X\otimes_k\bar{k}\) is perverse on \(X\otimes_k\bar{k}\).
One knows that the full subcategory
\(\operatorname{Perv}(X,\overline{\mathbb{Q}}_\ell)\) of
\(\mathrm{D}^b_c(X,\overline{\mathbb{Q}}_\ell)\) consisting of perverse
objects is an abelian category in which every object is of finite length. If
\(\ell\) is fixed, we will often denote
\(\operatorname{Perv}(X,\overline{\mathbb{Q}}_\ell)\) simply
\(\operatorname{Perv}(X)\). The objects of \(\operatorname{Perv}(X)\) are
sometimes called ``perverse sheaves'' on \(X\). However, we will call them
``perverse objects'' to avoid confusion with ``honest'' sheaves.

\sourcepara{2.3.1.1}
We now recall from \cite[1.3]{BBD} the theory of the perverse truncations
\({}^p\tau_{\leq i}(K)\) and \({}^p\tau_{\geq i}(K)\) and of the perverse
cohomology sheaves \({}^p\mathcal{H}^i(K)\) attached to an object \(K\) in
\(\mathrm{D}^b_c(X,\overline{\mathbb{Q}}_\ell)\). This will be used (only) in
\hyperref[sec:2.12]{\S2.12}. Inside \(\mathrm{D}^b_c(X,\overline{\mathbb{Q}}_\ell)\), we
denote by \({}^p\mathrm{D}^{\leq 0}\) (respectively
\({}^p\mathrm{D}^{\geq 0}\)) the full subcategory consisting of those objects
\(K\) which are semiperverse (respectively, those objects \(K\) such that
\(\mathbb{D}K\) is semiperverse). For each integer \(i\), we define
\({}^p\mathrm{D}^{\leq i}\) (respectively \({}^p\mathrm{D}^{\geq i}\)) to be
the full subcategory of \(\mathrm{D}^b_c(X,\overline{\mathbb{Q}}_\ell)\)
consisting of those objects \(K\) such that \(K[i]\) lies in
\({}^p\mathrm{D}^{\leq 0}\) (respectively \({}^p\mathrm{D}^{\geq 0}\)).
Duality interchanges \({}^p\mathrm{D}^{\leq i}\) and
\({}^p\mathrm{D}^{\geq -i}\). By \cite[1.3.3]{BBD}, the inclusion of
\({}^p\mathrm{D}^{\leq i}\) (respectively of \({}^p\mathrm{D}^{\geq i}\)) into
\(\mathrm{D}^b_c(X,\overline{\mathbb{Q}}_\ell)\) admits a right
(respectively left) adjoint, noted \({}^p\tau_{\leq i}\) (respectively
\({}^p\tau_{\geq i}\)). It is tautological that
\[
  ({}^p\tau_{\leq i}(K))[j]={}^p\tau_{\leq 0}(K[j]),\qquad
  ({}^p\tau_{\geq i}(K))[j]={}^p\tau_{\geq 0}(K[j]).
\]
For any \(i\), and any \(K\) in
\(\mathrm{D}^b_c(X,\overline{\mathbb{Q}}_\ell)\), we have a distinguished
triangle
\[
  {}^p\tau_{\leq i}(K)\to K\to{}^p\tau_{\geq i+1}(K).
\]
For any two integers \(a,b\), there is a canonical isomorphism
\cite[1.3.5]{BBD} between the two composites
\[
  {}^p\tau_{\geq a}\circ{}^p\tau_{\leq b}
  \cong
  {}^p\tau_{\leq b}\circ{}^p\tau_{\geq a}.
\]
Unless \(a\leq b\), both of these composites are zero.

Now take the special case \(a=b=0\). In this case, the composite functor above
has values in \(\operatorname{Perv}(X,\overline{\mathbb{Q}}_\ell)\), the
intersection of \({}^p\mathrm{D}^{\leq 0}\) and of
\({}^p\mathrm{D}^{\geq 0}\). For \(K\) in
\(\mathrm{D}^b_c(X,\overline{\mathbb{Q}}_\ell)\), we define
\[
  {}^p\mathcal{H}^0(K):=
  {}^p\tau_{\geq 0}({}^p\tau_{\leq 0}(K))
  \cong
  {}^p\tau_{\leq 0}({}^p\tau_{\geq 0}(K)).
\]
For each integer \(i\), we define
\[
  {}^p\mathcal{H}^i(K):={}^p\mathcal{H}^0(K[i]),
\]
or equivalently,
\[
  {}^p\mathcal{H}^i(K)[-i]:=
  {}^p\tau_{\geq i}({}^p\tau_{\leq i}(K))
  \cong
  {}^p\tau_{\leq i}({}^p\tau_{\geq i}(K)).
\]
One shows \cite[1.3.6]{BBD} that \({}^p\mathcal{H}^0\) is a cohomological
functor from \(\mathrm{D}^b_c(X,\overline{\mathbb{Q}}_\ell)\) to the abelian
category \(\operatorname{Perv}(X,\overline{\mathbb{Q}}_\ell)\): a distinguished
triangle in \(\mathrm{D}^b_c(X,\overline{\mathbb{Q}}_\ell)\) gives rise to a
long exact sequence of perverse cohomology sheaves. The functors
\({}^p\tau_{\leq 0}\) and \({}^p\tau_{\geq 0}\) are interchanged by duality,
whence
\[
  {}^p\mathcal{H}^0(\mathbb{D}K):=
  \mathbb{D}({}^p\mathcal{H}^0(K)),
\]
and hence for every \(i\), we have
\[
  {}^p\mathcal{H}^i(\mathbb{D}K):=
  \mathbb{D}({}^p\mathcal{H}^{-i}(K)).
\]
Given any \(K\) in \(\mathrm{D}^b_c(X,\overline{\mathbb{Q}}_\ell)\), all but
finitely many of its perverse cohomology sheaves
\({}^p\mathcal{H}^i(K)\) vanish. (By duality, it suffices to show that
\({}^p\mathcal{H}^i(K)\) vanishes for \(i\) sufficiently large. But \(K[i]\)
is trivially semiperverse for \(i\) sufficiently large, so it suffices to show
that for \(K\) semiperverse, \({}^p\mathcal{H}^i(K)\) vanishes for \(i>0\). In
fact, one knows \cite[1.3.7]{BBD} that \(K\) in
\(\mathrm{D}^b_c(X,\overline{\mathbb{Q}}_\ell)\) is semiperverse if and only
if \({}^p\mathcal{H}^i(K)\) vanishes for \(i>0\). Therefore
\({}^p\mathcal{H}^i(K)=0\) for \(i\) outside \([a,b]\) if and only if \(K\)
lies in
\({}^p\mathrm{D}^{[a,b]}:={}^p\mathrm{D}^{\geq a}\cap{}^p\mathrm{D}^{\leq b}\).)
Moreover, any \(K\) in \(\mathrm{D}^b_c(X,\overline{\mathbb{Q}}_\ell)\) is a
successive extension of its shifted perverse cohomology sheaves
\({}^p\mathcal{H}^i(K)[-i]\), as one sees by induction on the length \(b-a\) of
the shortest interval \([a,b]\) such that \(K\) lies in
\({}^p\mathrm{D}^{[a,b]}\), using the distinguished triangles
\({}^p\tau_{\leq i}(K)\to K\to{}^p\tau_{\geq i+1}(K)\).

\sourcepara{2.3.2}
If \(X\) is smooth over \(k\), everywhere of relative dimension \(\dim X\), the
simplest example of a perverse object on \(X\) is provided by starting with a
lisse sheaf \(\mathcal{F}\) on \(X\), and taking the object
\(\mathcal{F}[\dim X]\) of
\(\mathrm{D}^b_c(X,\overline{\mathbb{Q}}_\ell)\) obtained by placing
\(\mathcal{F}\) in degree \(-\dim X\). The object
\(\mathcal{F}[\dim X]\) is trivially semiperverse, and its dual
\[
  \mathbb{D}(\mathcal{F}[\dim X])
  =(\mathcal{F}^{\vee}(\dim X))[\dim X],
\]
being of the same form, is also. If \(X\) is connected, and if \(\mathcal{F}\)
is irreducible as a lisse sheaf, i.e., as a representation of
\(\pi_1(X,x)\), then \(\mathcal{F}[\dim X]\) is a simple object of
\(\operatorname{Perv}(X)\).

\sourceheading[lem:2.3.2.1]{Lemma 2.3.2.1}
For \(X\) smooth over \(k\), everywhere of relative dimension \(\dim X\),
consider the following two (rather special) properties of an object \(K\) in
\(\mathrm{D}^b_c(X,\overline{\mathbb{Q}}_\ell)\):
\begin{enumerate}[label=(\arabic*)]
\item each of the cohomology sheaves \(\mathcal{H}^i(K)\) is lisse,
\item each of the perverse cohomology sheaves \({}^p\mathcal{H}^i(K)\) is of
the form (a lisse sheaf on \(X\))\([\dim X]\).
\end{enumerate}
These properties are equivalent, and, if they hold, then the perverse and
ordinary cohomology sheaves of \(K\) are related by
\begin{enumerate}[label=(\arabic*)]
\setcounter{enumi}{2}
\item
\[
  {}^p\mathcal{H}^i(K)=\mathcal{H}^{i-\dim X}(K)[\dim X].
\]
\end{enumerate}

\begin{proof}
Suppose first that (a) holds. For the usual truncation functors
\(\tau_{\leq n}\), we have distinguished triangles
\[
  \tau_{\leq n-1}K\to\tau_{\leq n}K\to\mathcal{H}^n(K)[-n].
\]
For \(n\) sufficiently negative, we have \(\tau_{\leq n}K=0\), and for
\(n\) sufficiently positive, we have
\(\tau_{\leq n-1}K\cong\tau_{\leq n}K\cong K\). Because
\({}^p\mathcal{H}^0\) is a cohomological functor, we get a long exact
sequence
\[
  \to{}^p\mathcal{H}^i(\tau_{\leq n-1}K)
  \to{}^p\mathcal{H}^i(\tau_{\leq n}K)
  \to{}^p\mathcal{H}^i(\mathcal{H}^n(K)[-n])
  \to{}^p\mathcal{H}^{i+1}(\tau_{\leq n-1}K)\to .
\]
Since \(\mathcal{H}^n(K)\) is lisse,
\(\mathcal{H}^n(K)[\dim X]\) is perverse, and hence
\[
  {}^p\mathcal{H}^a(\mathcal{H}^n(K)[-n])
  =
  \begin{cases}
    \mathcal{H}^n(K)[\dim X] & \text{for } a=n+\dim X,\\
    0 & \text{for } a\neq n+\dim X.
  \end{cases}
\]
Using this fact, and the long exact sequences above, one shows by induction on
\(n\) that
\begin{enumerate}[label=(\arabic*)]
\item \({}^p\mathcal{H}^a(\tau_{\leq n}K)=0\) for \(a>n+\dim X\),
\item the map
\(\tau_{\leq n}K\to\mathcal{H}^n(K)[-n]\) induces an isomorphism on
\({}^p\mathcal{H}^a\) for \(a=n+\dim X\),
\item the map \(\tau_{\leq n-1}K\to\tau_{\leq n}K\) induces an isomorphism on
\({}^p\mathcal{H}^a\) for \(a<n+\dim X\).
\end{enumerate}
Once we have these facts, then for any \(a\), we get
\[
  {}^p\mathcal{H}^{a+\dim X}(\tau_{\leq a}K)
  \cong{}^p\mathcal{H}^{a+\dim X}(\tau_{\leq a+1}K)
  \cong\cdots\cong{}^p\mathcal{H}^{a+\dim X}(K)
\]
by successive application of (3), and then by (2) we get
\[
  {}^p\mathcal{H}^{a+\dim X}(\tau_{\leq a}K)
  \cong
  {}^p\mathcal{H}^{a+\dim X}(\mathcal{H}^a(K)[-a])
  =\mathcal{H}^a(K)[\dim X].
\]
Thus we obtain
\({}^p\mathcal{H}^{a+\dim X}(K)\cong\mathcal{H}^a(K)[\dim X]\), which proves
(c), and consequently (b).

Now suppose that (b) holds. Repeat the above argument, with usual truncation
replaced by perverse truncation, with perverse cohomology sheaves replaced by
usual cohomology sheaves, and with \(\dim X\) replaced by \(-\dim X\). Again
we find (c), and consequently (a).
\end{proof}

\sourcepara{2.3.3}
Given a locally closed subscheme \(Y\) of \(X\) such that \(Y\) is affine, the
inclusion \(j:Y\to X\) is both affine and quasifinite (factor it as the open
immersion of \(Y\) into its closure \(\bar{Y}\), followed by the closed
immersion of \(\bar{Y}\) into \(X\)). So for a perverse object \(K\) on \(Y\),
both \(j_!K\) and \(\mathrm{R}j_*K\) are perverse on \(X\), and as functors
from \(\operatorname{Perv}(Y)\) to \(\operatorname{Perv}(X)\) both \(j_!\) and
\(\mathrm{R}j_*\) are exact. There is a natural ``forget supports'' map from
\(j_!K\) to \(\mathrm{R}j_*K\), and as \(\operatorname{Perv}(X)\) is an
abelian category it makes sense to form
\[
  j_{!*}(K):=\operatorname{Image}(j_!K\to \mathrm{R}j_*K)
  \in\operatorname{Perv}(X),
\]
called the ``middle extension'' from \(Y\) to \(X\) of the perverse object
\(K\). The functor \(j_{!*}\) is end-exact (i.e., it preserves both injections
and surjections, cf. the appendix to this chapter) from
\(\operatorname{Perv}(Y)\) to \(\operatorname{Perv}(X)\), it carries simple
objects to simple objects, and it commutes with duality. Despite the erroneous
assertion in \cite[8.1.4]{Ka-ESDE}, the functor \(j_{!*}\) is not exact in
general.

\sourcepara{2.3.3.1}
For \(K\) and \(L\) perverse on \(Y\) as in \hyperref[par:2.3.3]{(2.3.3)} above, the functors
\(j_{!*}\) and \(j^*\) induce natural maps of \(\operatorname{Hom}\) groups,
\[
  \operatorname{Hom}_Y(K,L)\xrightarrow{j_{!*}}
  \operatorname{Hom}_X(j_{!*}K,j_{!*}L)
\]
and
\[
  \operatorname{Hom}_X(j_{!*}K,j_{!*}L)\xrightarrow{j^*}
  \operatorname{Hom}_Y(K,L).
\]
These maps are inverse isomorphisms. To see this, we argue as follows. The
composite functor \(K\mapsto j^*j_{!*}K\) is the identity, so the composite map
\[
  \operatorname{Hom}_Y(K,L)\to
  \operatorname{Hom}_X(j_{!*}K,j_{!*}L)\to
  \operatorname{Hom}_Y(K,L)
\]
is the identity on \(\operatorname{Hom}_Y(K,L)\).

Because \(j_{!*}L\) is a subobject of \(\mathrm{R}j_*L\) in
\(\operatorname{Perv}(X)\), the restriction map
\[
  \operatorname{Hom}_X(j_{!*}K,j_{!*}L)\xrightarrow{j^*}
  \operatorname{Hom}_Y(K,L)
\]
is injective: we have
\[
\begin{aligned}
  \operatorname{Hom}_X(j_{!*}K,j_{!*}L)
  &\subset
  \operatorname{Hom}_X(j_{!*}K,\mathrm{R}j_*L)\\
  &=\operatorname{Hom}_Y(j^*j_{!*}K,L)
  =\operatorname{Hom}_Y(K,L).
\end{aligned}
\]
Now start with \(\varphi\) in \(\operatorname{Hom}_X(j_{!*}K,j_{!*}L)\). Both
\(\varphi\) and \(j_{!*}j^*(\varphi)\) have the same restriction, namely
\(j^*(\varphi)\), in \(\operatorname{Hom}_Y(K,L)\), hence
\(\varphi=j_{!*}j^*(\varphi)\) in
\(\operatorname{Hom}_X(j_{!*}K,j_{!*}L)\).

\sourcepara{2.3.4}
One knows that for any simple object \(S\) of \(\operatorname{Perv}(X)\) there
exists an affine locally closed subscheme \(j:Y\to X\) such that \(Y\) is
smooth over \(k\) and irreducible, and an irreducible lisse sheaf
\(\mathcal{F}\) on \(Y\) such that
\(S\) is \(j_{!*}(\mathcal{F}[\dim Y])\). Given the simple object \(S\), we
construct \(Y\) and \(\mathcal{F}\) as follows: the closure \(\bar{Y}\) of
\(Y\) is precisely the closure of the support of
\(\oplus_i\mathcal{H}^iS\), \(Y\) is any smooth affine open set of \(\bar{Y}\)
on which all the \(\mathcal{H}^iS\) are lisse, and \(\mathcal{F}\) is
\(\mathcal{H}^{-\dim Y}(S)|Y\).

\sourcepara{2.3.5}
An object \(S\) of \(\operatorname{Perv}(X)\) is called geometrically simple if
its inverse image on \(X\otimes_k\bar{k}\) is simple. Of course
``geometrically simple'' \(\Rightarrow\) ``simple''.

\sourcepara{2.3.6}
Consider the special case when \(X/k\) is a smooth, geometrically connected
curve. Then an object \(K\) of
\(\mathrm{D}^b_c(X,\overline{\mathbb{Q}}_\ell)\) is perverse if and only if
\[
  \mathcal{H}^iK=0\text{ for }i\neq -1,0,\qquad
  \mathcal{H}^{-1}K\text{ has no nonzero punctual sections},\qquad
  \mathcal{H}^0K\text{ is punctual}.
\]
We call a perverse object \(K\) punctual if \(K=\mathcal{H}^0(K)\), and we
call \(K\) ``nonpunctual'' if \(\mathcal{H}^0K=0\). If \(\mathcal{F}\) is a
lisse sheaf on an open nonempty open set \(j:U\to X\), then the middle
extension \(j_{!*}(\mathcal{F}[1])\) is none other than
\((j_*\mathcal{F})[1]\). It is for this reason that we adapted the terminology
``middle extension'' for sheaves of the type \(j_*\mathcal{F}\) with
\(\mathcal{F}\) lisse on \(U\). The dual
\(\mathbb{D}(j_{!*}(\mathcal{F}[1]))\) of such a middle extension is given by
\[
  \mathbb{D}(j_{!*}(\mathcal{F}[1]))
  =\mathbb{D}(j_*\mathcal{F}[1])
  =j_{!*}(\mathbb{D}(\mathcal{F}[1]))
  =j_*(\mathcal{F}^{\vee})[1](1).
\]
Any perverse sheaf \(K\) on \(X\) has a natural two step filtration, whose
associated graded pieces are (punctual, middle extension, punctual). To see
this we first filter \(K\) by its subobject \(\mathcal{H}^{-1}(K)[1]\), which
sits in the short exact sequence of perverse sheaves
\[
  0\to\mathcal{H}^{-1}(K)[1]\to K\to\mathcal{H}^0(K)\to0.
\]
Now denote by \(j:U\to X\) the inclusion of a nonempty affine open set \(U\) on
which \(\mathcal{H}^{-1}(K)\) is lisse. Since \(\mathcal{H}^{-1}(K)\) has no
nonzero punctual sections, we have a short exact sequence of usual sheaves on
\(X\)
\[
  0\to\mathcal{H}^{-1}(K)\to j_*j^*\mathcal{H}^{-1}(K)\to \operatorname{pct}'\to0.
\]
Shifting by \([1]\) and rotating the triangle, we get a distinguished triangle
\[
  \operatorname{pct}'\to\mathcal{H}^{-1}(K)[1]\to
  (j_*j^*\mathcal{H}^{-1}(K))[1],
\]
i.e., a short exact sequence of perverse sheaves
\[
  0\to\operatorname{pct}'\to\mathcal{H}^{-1}(K)[1]\to
  (j_*j^*\mathcal{H}^{-1}(K))[1]\to0.
\]
The filtration in question is
\(\operatorname{pct}'\subset\mathcal{H}^{-1}(K)[1]\subset K\).

There are two types of simple perverse object on \(X\):
\begin{enumerate}[label=(\arabic*)]
\item the punctual ones, whose \(Y\) is a single closed point \(x\) of \(X\);
the corresponding simple objects are \(x_*\mathcal{F}\), where
\(\mathcal{F}\) is an irreducible representation of
\(\operatorname{Gal}(\bar{k}/k(x))\) (so if \(k\) is algebraically closed, only
the delta sheaf \(\delta_x:=x_*\overline{\mathbb{Q}}_\ell\) supported at
\(x\));
\item the nonpunctual ones, whose \(Y\) is a nonempty open set \(j:U\to X\) of
\(X\); the corresponding simple objects are \((j_*\mathcal{F})[1]\), where
\(\mathcal{F}\) is an ``arithmetically irreducible'' lisse sheaf on \(U\),
i.e., one whose representation of \(\pi_1(U,\bar{u})\) is irreducible (so the
nonpunctual simples which are geometrically simple are precisely the
\(\mathcal{F}[1]\) where \(\mathcal{F}\) is an ``irreducible middle extension
sheaf'' in the terminology of \hyperref[par:2.0.3]{(2.0.3)}).
\end{enumerate}

\sourcepara{2.3.7}
If \(k\) is \(\mathbb{C}\), then for any \(\ell\) the exact, fully faithful
``passage to the analytic'' functor
\[
  \mathrm{D}^b_c(X,\overline{\mathbb{Q}}_\ell)\to
  \mathrm{D}^b_c(X^{\mathrm{an}},\overline{\mathbb{Q}}_\ell)
\]
induces an exact fully faithful functor
\[
  \operatorname{Perv}(X,\overline{\mathbb{Q}}_\ell)\to
  \operatorname{Perv}(X^{\mathrm{an}},\overline{\mathbb{Q}}_\ell).
\]
Everything said above about
\(\operatorname{Perv}(X,\overline{\mathbb{Q}}_\ell)\) holds also for
\(\operatorname{Perv}(X^{\mathrm{an}},\overline{\mathbb{Q}}_\ell)\), and all
cohomological constructions, including middle extension \(j_{!*}\), commute
with the ``passage to the analytic'' functor \hyperref[par:2.2.3]{(2.2.3)}.

\section{Review of Fourier Transform}
\label{sec:2.4}

\sourcepara{2.4.1}
Suppose \(X/k\) is \(\mathbb{A}^1/k\), \(k\) a perfect field with
\(\operatorname{char}(k):=p>0\). For each \(\ell\neq p\), we fix a nontrivial
additive character
\(\psi:\mathbb{F}_p\to\overline{\mathbb{Q}}_\ell^\times\), and its associated
lisse, rank one Artin--Schreier sheaf \(\mathcal{L}_\psi\) on
\(\mathbb{A}^1\). The derived category versions of Fourier Transform are
defined by
\[
\begin{aligned}
  \operatorname{FT}_{\psi,!}(K)
  &:=
  \mathrm{R}(\operatorname{pr}_2)_!
  (\operatorname{pr}_1^*K\otimes\mathcal{L}_\psi(xy))[1],\\
  \operatorname{FT}_{\psi,*}(K)
  &:=
  \mathrm{R}(\operatorname{pr}_2)_*
  (\operatorname{pr}_1^*K\otimes\mathcal{L}_\psi(xy))[1].
\end{aligned}
\]
Both are exact functors from
\(\mathrm{D}^b_c(\mathbb{A}^1,\overline{\mathbb{Q}}_\ell)\) to itself, which
are essentially interchanged by duality:
\[
  \mathbb{D}(\operatorname{FT}_{\psi,!}K)
  =\operatorname{FT}_{\psi,*}([-1]^*\mathbb{D}K)(1).
\]
It is easy to prove that \(\operatorname{FT}_{\psi,!}\) is essentially
involutive:
\[
  \operatorname{FT}_{\psi,!}\operatorname{FT}_{\psi,!}
  \cong[-1]^*(-1);
\]
by duality it follows that the same holds for \(\operatorname{FT}_{\psi,*}\).

\sourcepara{2.4.2}
The ``miracle'' of Fourier Transform is that there is really only one: the
natural ``forget supports'' map
\(\operatorname{FT}_{\psi,!}\to\operatorname{FT}_{\psi,*}\) is an isomorphism.
We denote it \(\operatorname{FT}_\psi\). As \(\operatorname{FT}_\psi\) (viewed
as \(\operatorname{FT}_{\psi,*}\)) preserves semiperversity, it follows from
the miracle that \(\operatorname{FT}_\psi\) preserves perversity, and so
defines an exact autoequivalence of
\(\operatorname{Perv}(\mathbb{A}^1)\). In particular,
\(\operatorname{FT}_\psi\) sends perverse simple objects to perverse simple
objects.

\section{Review of convolution}
\label{sec:2.5}

\sourcepara{2.5.1}
Suppose \(G\) is a smooth separated \(k\)-group scheme of finite type of
relative dimension noted \(\dim G\), \(\pi:G\times_kG\to G\) the
multiplication map, \(e:\operatorname{Spec}(k)\to G\) the identity section.
Given two objects \(K\) and \(L\) in
\(\mathrm{D}^b_c(G,\overline{\mathbb{Q}}_\ell)\), we define their ``compact''
or ``!'' convolution, denoted \(K*_!L\), by
\[
\begin{aligned}
  K*_!L&:=\mathrm{R}\pi_!(K\times L)\in
  \mathrm{D}^b_c(G,\overline{\mathbb{Q}}_\ell),\\
  K*_*L&:=\mathrm{R}\pi_*(K\times L)\in
  \mathrm{D}^b_c(G,\overline{\mathbb{Q}}_\ell).
\end{aligned}
\]
Duality interchanges the two sorts of convolution:
\[
  \mathbb{D}(K*_!L)\cong\mathbb{D}(K)*_*\mathbb{D}(L),\qquad
  \mathbb{D}(K*_*L)\cong\mathbb{D}(K)*_!\mathbb{D}(L).
\]
By the Leray spectral sequence and the Kunneth formula, we have
\(\operatorname{Gal}(\bar{k}/k)\)-equivariant isomorphisms of cohomology
algebras
\[
\begin{aligned}
  \mathrm{H}^*_c(G\otimes\bar{k},K*_!L)
  &\cong
  \mathrm{H}^*_c((G\times G)\otimes\bar{k},K\times L)\\
  &\cong
  \mathrm{H}^*_c(G\otimes\bar{k},K)\otimes
  \mathrm{H}^*_c(G\otimes\bar{k},L),\\
  \mathrm{H}^*(G\otimes\bar{k},K*_*L)
  &\cong
  \mathrm{H}^*((G\times G)\otimes\bar{k},K\times L)\\
  &\cong
  \mathrm{H}^*(G\otimes\bar{k},K)\otimes
  \mathrm{H}^*(G\otimes\bar{k},L).
\end{aligned}
\]

\sourcepara{2.5.2}
If we start with two (usual or perverse) sheaves \(\mathcal{F}\) and
\(\mathcal{G}\) on \(G\), viewed as objects of
\(\mathrm{D}^b_c(G,\overline{\mathbb{Q}}_\ell)\), their convolutions
\(\mathcal{F}*_!\mathcal{G}\) and \(\mathcal{F}*_*\mathcal{G}\) are
``really'' objects of \(\mathrm{D}^b_c(G,\overline{\mathbb{Q}}_\ell)\), and
not simply single (usual or perverse) sheaves placed in some degree. It is this
``instability'' of (usual or perverse) sheaves themselves under convolution
that makes \(\mathrm{D}^b_c(G,\overline{\mathbb{Q}}_\ell)\) the natural
setting for systematically discussing convolution.

\sourcepara{2.5.3}
For the convenience of the reader, we collect from \cite[8.1.9--10]{Ka-ESDE}
the standard facts about convolution.
\begin{enumerate}[label=(\arabic*)]
\setcounter{enumi}{-1}
\item If \(K\) and \(L\) are semiperverse (resp. perverse) objects on \(G\),
then \(K\times L\) is semiperverse (resp. perverse) on \(G\times_kG\).
Therefore if \(G\) is affine, and if \(K\) and \(L\) are both semiperverse on
\(G\), then \(K*_*L\) is semiperverse on \(G\). If \(K\) and \(L\) are both
perverse on \(G\) and if moreover the natural ``forget supports'' map is an
isomorphism \(K*_!L\cong K*_*L\), then \(K*_!L\cong K*_*L\) is perverse (its
dual being \(\mathbb{D}(K)*_*\mathbb{D}(L)\)).
\item Each sort of convolution is associative, and for each the
\(\delta\)-sheaf
\[
  \delta_e:=e_*\overline{\mathbb{Q}}_\ell
\]
supported at the identity of \(G\) is a two-sided identity object. If \(G\) is
commutative, then each sort of convolution is commutative as well.
\item[(2a)] If \(\varphi:G\to H\) is a homomorphism of smooth separated
\(k\)-group schemes of finite type, then for \(K\) and \(L\) on \(G\) we have
\[
\begin{aligned}
  \mathrm{R}\varphi_*(K*_*L)
  &\cong(\mathrm{R}\varphi_*K)*_*(\mathrm{R}\varphi_*L),\\
  \mathrm{R}\varphi_!(K*_!L)
  &\cong(\mathrm{R}\varphi_!K)*_!(\mathrm{R}\varphi_!L).
\end{aligned}
\]
\item[(2b)] If \(\varphi:G\to H\) is a homomorphism, then for \(K\) on \(G\)
and \(L\) on \(H\) we have
\[
\begin{aligned}
  \varphi^*((\mathrm{R}\varphi_!K)*_!L)
  &\cong K*_!(\varphi^*L),\\
  \varphi^!((\mathrm{R}\varphi_*K)*_*L)
  &\cong K*_*(\varphi^!L).
\end{aligned}
\]
\setcounter{enumi}{2}
\item For \(g\in G(k)\) denote by \(T_g:G\to G\) the map
\(x\mapsto gx\) ``left translation by \(g\)'', and by
\(\delta_g:=(T_g)_*(\delta_e)\) the delta sheaf supported at \(g\). Then for
\(g\in G(k)\), we have
\[
\begin{aligned}
  (T_g)_*&=\mathrm{R}(T_g)_*=(T_g)_!=\mathrm{R}(T_g)_!,\\
  (T_g)_*(K*_*L)&\cong((T_g)_*K)*_*L,\\
  (T_g)_*(K*_!L)&\cong((T_g)_*K)*_!L,\\
  (T_g)_*(L)&\cong(\delta_g)*L.
\end{aligned}
\]
Moreover, if \(G\) is commutative, then for \(g,h\) in \(G(k)\), we have
\[
\begin{aligned}
  (T_{gh})_*(K*_*L)&\cong((T_g)_*K)*_*((T_h)_*L),\\
  (T_{gh})_*(K*_!L)&\cong((T_g)_*K)*_!((T_h)_*L).
\end{aligned}
\]
\item If \(G\) is commutative, geometrically connected, and defined over a
finite subfield \(k_0\) of \(k\), then for every
\(\overline{\mathbb{Q}}_\ell\)-valued character \(\chi\) of \(G(k_0)\), the
associated lisse rank one \(\mathcal{L}_\chi\) on \(G\) obtained from pushing
out the Lang torsor by \(\chi\) satisfies
\(\pi^*\mathcal{L}_\chi\cong\mathcal{L}_\chi\times\mathcal{L}_\chi\), whence
by the projection formula
\[
\begin{aligned}
  (K*_!L)\otimes\mathcal{L}_\chi
  &\cong(K\otimes\mathcal{L}_\chi)*_!(L\otimes\mathcal{L}_\chi),\\
  (K*_*L)\otimes\mathcal{L}_\chi
  &\cong(K\otimes\mathcal{L}_\chi)*_*(L\otimes\mathcal{L}_\chi).
\end{aligned}
\]
\end{enumerate}

\sourcepara{2.5.4}
If \(k\) is \(\mathbb{C}\), then for any \(\ell\) the exact, fully faithful
``passage to the analytic'' functor
\[
  \mathrm{D}^b_c(X,\overline{\mathbb{Q}}_\ell)\to
  \mathrm{D}^b_c(X^{\mathrm{an}},\overline{\mathbb{Q}}_\ell)
\]
respects both sorts of convolution.

\section[Convolution operators: middle convolution]{Convolution operators on the category of perverse sheaves: middle
convolution}\label{sec:2.6}

\sourcepara{2.6.1}
Let \(k\) be an algebraically closed field. Suppose \(G\) is a connected smooth
affine \(k\)-group scheme of finite type of relative dimension noted
\(\dim G\), \(\pi:G\times_kG\to G\) the multiplication map. A perverse sheaf
\(K\) on \(G\) is said to have property \(\mathcal{P}_!\) (respectively
property \(\mathcal{P}_*\)) if for any perverse sheaf \(L\) on \(G\), the
\(!\) convolution \(L*_!K\) (respectively the \(*\) convolution \(L*_*K\)) is
again perverse. If \(K\) has property \(\mathcal{P}_!\) (respectively property
\(\mathcal{P}_*\)), the functor \(L\mapsto L*_!K\) (respectively
\(L\mapsto L*_*K\)) is an exact functor from \(\operatorname{Perv}(G)\) to
itself. Notice that since duality interchanges the two sorts of convolution,
\[
  \mathbb{D}(L*_!K)\cong\mathbb{D}(L)*_*\mathbb{D}(K),
\]
we have the equivalence
\[
  K\text{ has }\mathcal{P}_!\Longleftrightarrow
  \mathbb{D}K\text{ has }\mathcal{P}_*.
\]

\sourcepara{2.6.2}
A perverse sheaf \(K\) on \(G\) is said to have property \(\mathcal{P}\) if it
has both properties \(\mathcal{P}_!\) and \(\mathcal{P}_*\). If \(K\) has
property \(\mathcal{P}\), we define, for any perverse \(L\) on \(G\), the
middle convolution \(L*_{\mathrm{mid}}K\) to be the perverse sheaf on \(G\)
defined as the image, in the abelian category of perverse sheaves on \(G\), of
the natural ``forget supports'' map
\[
  L*_{\mathrm{mid}}K:=\operatorname{image}(L*_!K\to L*_*K).
\]
For fixed \(K\) in \(\mathcal{P}\), \(L\mapsto L*_{\mathrm{mid}}K\) is
end-exact (cf. \hyperref[sec:2.17]{\S2.17}).

\sourcepara{2.6.3}
By definition, the map
\[
  L*_!K\to L*_{\mathrm{mid}}K
\]
is surjective in \(\operatorname{Perv}\), and the map
\[
  L*_{\mathrm{mid}}K\to L*_*K
\]
is injective in \(\operatorname{Perv}\).

\sourcepara{2.6.4}
If \(K\) and \(L\) have \(\mathcal{P}_!\) (respectively \(\mathcal{P}_*\)), so
does \(K*_!L\) (respectively \(L*_*K\)), just by the associativity of
convolution. What about \(\mathcal{P}\) itself? We will see later that if
\(K\) and \(L\) both have \(\mathcal{P}\), and if \(G\) is either
\(\mathbb{G}_m\) or \(\mathbb{A}^1\), then \(K*_{\mathrm{mid}}L\) also has
\(\mathcal{P}\), but we do not know this for more general \(G\). Returning to
the general situation, we have

\sourceheading[lem:2.6.5]{Lemma 2.6.5}
In the situation \hyperref[par:2.6.1]{(2.6.1)}, if \(\mathcal{F}\), \(K\) and \(L\) are perverse
sheaves on \(G\) which all have \(\mathcal{P}\), then
\[
  (\mathcal{F}*_{\mathrm{mid}}K)*_{\mathrm{mid}}L
  =
  \mathcal{F}*_{\mathrm{mid}}(K*_{\mathrm{mid}}L).
\]

\begin{proof}
We will show this equality by showing that both are the image of
\(\mathcal{F}*_!K*_!L\) in \(\mathcal{F}*_*K*_*L\). Indeed, we may factor
this ``forget supports'' map as the composition of the two surjective maps
\[
\begin{aligned}
  (\mathcal{F}*_!K)*_!L
  &\to(\mathcal{F}*_{\mathrm{mid}}K)*_!L
  \quad(\text{surj. as }*_!L\text{ is exact on }\operatorname{Perv}),\\
  (\mathcal{F}*_{\mathrm{mid}}K)*_!L
  &\to(\mathcal{F}*_{\mathrm{mid}}K)*_{\mathrm{mid}}L
\end{aligned}
\]
and the two injective maps
\[
\begin{aligned}
  (\mathcal{F}*_{\mathrm{mid}}K)*_{\mathrm{mid}}L
  &\to(\mathcal{F}*_{\mathrm{mid}}K)*_*L,\\
  (\mathcal{F}*_{\mathrm{mid}}K)*_*L
  &\to(\mathcal{F}*_*K)*_*L
  \quad(\text{inj. as }*_*L\text{ is exact on }\operatorname{Perv}).
\end{aligned}
\]
Thus \((\mathcal{F}*_{\mathrm{mid}}K)*_{\mathrm{mid}}L\) is the image of
\(\mathcal{F}*_!K*_!L\) in \(\mathcal{F}*_*K*_*L\). Rearranging the
parentheses shows that \(\mathcal{F}*_{\mathrm{mid}}(K*_{\mathrm{mid}}L)\) is
also this image.
\end{proof}

\sourceheading[rem:2.6.6]{Remark 2.6.6}
If we only assume that \(K\) and \(L\) have \(\mathcal{P}\), but not
\(\mathcal{F}\), then the right hand side
\(\mathcal{F}*_{\mathrm{mid}}(K*_{\mathrm{mid}}L)\) isn't defined, since we
don't know that \(K*_{\mathrm{mid}}L\) has \(\mathcal{P}\). Later, we will know
that \(K*_{\mathrm{mid}}L\) has \(\mathcal{P}\) if both \(K\) and \(L\) do, at
least on both \(\mathbb{A}^1\) and \(\mathbb{G}_m\), but it will not be true
that for any \(\mathcal{F}\) in \(\operatorname{Perv}\), we have
\(\mathcal{F}*_{\mathrm{mid}}(K*_{\mathrm{mid}}L)=
(\mathcal{F}*_{\mathrm{mid}}K)*_{\mathrm{mid}}L\). Indeed, it can be that
\(K*_{\mathrm{mid}}L=\delta_e\), so the left side is \(\mathcal{F}\), but
that already \(\mathcal{F}*_!K=0\), so a fortiori
\(\mathcal{F}*_{\mathrm{mid}}K=0\). (Example:
\(\mathcal{F}=\overline{\mathbb{Q}}_\ell[1]\) on \(\mathbb{A}^1\),
\(K=\mathcal{L}_\chi[1]\), \(L=\mathcal{L}_{\bar{\chi}}[1]\).)

\sourceheading[lem:2.6.7]{Lemma 2.6.7}
Let \(k\) be an algebraically closed field. Suppose \(G\) is a connected smooth
affine \(k\)-group scheme of finite type. Let \(K\) and \(L\) be any perverse
sheaves on \(G\). Then \(K*_!L\) is perverse if and only if it is
semiperverse.

\begin{proof}
By definition, \(K*_!L\) is perverse if and only if both it and its dual are
semiperverse. So it suffices to show that for \(K\) and \(L\) perverse, the
dual of \(K*_!L\) is semiperverse. But this dual is
\(\mathbb{D}(K*_!L)=\mathbb{D}(K)*_*\mathbb{D}(L)\), the \(*\) convolution of
the perverse sheaves \(\mathbb{D}(K)\) and \(\mathbb{D}(L)\), hence is
semiperverse because \(G\) is affine (cf. \hyperref[par:2.5.3]{(2.5.3 (0))}).
\end{proof}

\sourceheading[lem:2.6.8]{Lemma 2.6.8}
Let \(k\) be an algebraically closed field. Suppose \(G\) is a connected smooth
affine \(k\)-group scheme of finite type. Let \(K\) be perverse on \(G\). Then
the following conditions are equivalent:
\begin{enumerate}[label=(\alph*)]
\item \(K\) has property \(\mathcal{P}_!\).
\item \(K*_!M\) is perverse for every perverse irreducible \(M\) on \(G\).
\end{enumerate}

\begin{proof}
That (1) implies (2) is trivial. Suppose now that (2) holds, and let \(L\) be
perverse on \(G\). It is known \cite[1.3.6 and 4.3.1(i)]{BBD} that the
category of perverse sheaves on \(G\) is abelian, and every object is of finite
length. We proceed by induction on the length of \(L\). If \(L\) is
irreducible, then \(K*_!L\) is perverse by (2). In general, pick a perverse
irreducible \(M\) in \(L\), and consider the short exact sequence of perverse
sheaves
\[
  0\to M\to L\to N\to0.
\]
By induction, both \(K*_!M\) and \(K*_!N\) are perverse, and hence both are
semiperverse. In the derived category
\(\mathrm{D}^b_c(G,\overline{\mathbb{Q}}_\ell)\), the above exact sequence of
perverse sheaves is a distinguished triangle. Applying the functor
\(K*_!(\_)\) to it yields a distinguished triangle
\[
  0\to K*_!M\to K*_!L\to K*_!N\to0
\]
whose end terms are semiperverse. The long exact cohomology sequence
\[
  \to\mathcal{H}^i(K*_!M)\to
  \mathcal{H}^i(K*_!L)\to
  \mathcal{H}^i(K*_!N)\to
\]
shows that \(K*_!L\) is itself semiperverse, and hence (by the previous lemma)
perverse. Thus (2) implies (1).
\end{proof}

\sourcesubheading{The special case of relative dimension one}

We now turn to the special case when \(G\) in \hyperref[par:2.6.1]{(2.6.1)} is further assumed to be of
relative dimension one. Although such a \(G\) is isomorphic to either
\(\mathbb{G}_a\) or \(\mathbb{G}_m\), we will, to the extent possible, give a
unified treatment of the two cases.

\sourceheading[lem:2.6.9]{Lemma 2.6.9}
Let \(k\) be an algebraically closed field. Suppose \(G\) is a connected smooth
affine \(k\)-group scheme of finite type, of relative dimension one. Let \(K\)
be any perverse irreducible on \(G\) whose isomorphism class is not
translation-invariant. Then \(K\) has property \(\mathcal{P}_!\).

\begin{proof}
By the above two lemmas, it suffices to show that for any perverse irreducible
\(L\) on \(G\), \(K*_!L\) is semiperverse. Since \(G\) is one-dimensional,
\(K*_!L\) is semiperverse if and only if
\[
  \mathcal{H}^0(K*_!L)\text{ is punctual, and }
  \mathcal{H}^i(K*_!L)=0\text{ for }i>0.
\]
If \(K\) (respectively \(L\)) is punctual, then \(K*_!L\) is a translate of
\(L\) (respectively \(K\)), and hence \(K*_!L\) is certainly perverse. If
neither \(K\) nor \(L\) is punctual, then there exists a dense open set \(U\)
in \(G\), and irreducible lisse sheaves \(\mathcal{F}\) and \(\mathcal{G}\) on
\(U\), such that, denoting by \(j:U\to G\) the inclusion, \(K\) and \(L\) are
\(j_*\mathcal{F}[1]\) and \(j_*\mathcal{G}[1]\) respectively. The stalk of
\(\mathcal{H}^i(K*_!L)\) at any geometric point \(g\) of \(G\) is the
cohomology group
\[
  \mathrm{H}^{i+2}_c\bigl(G,\operatorname{Trans}_g^*(j_*\mathcal{F})
  \otimes(\operatorname{inv})^*(j_*\mathcal{G})\bigr).
\]
Since \(G\) is one-dimensional, this group vanishes for \(i>0\), and hence
\(\mathcal{H}^i(K*_!L)=0\) for \(i>0\). It remains to show that
\(\mathcal{H}^0(K*_!L)\) is punctual, i.e., that for all but at most finitely
many \(g\) in \(G(k)\), we have
\[
  \mathrm{H}^2_c\bigl(G,\operatorname{Trans}_g^*(j_*\mathcal{F})
  \otimes(\operatorname{inv})^*(j_*\mathcal{G})\bigr)=0.
\]
For fixed \(g\) in \(G(k)\), let us denote by \(U_g\) the dense open set
\[
  U_g:=\operatorname{Trans}_g^*U\cap(\operatorname{inv})^*U
\]
on which both \(\operatorname{Trans}_g^*(j_*\mathcal{F})\) and
\((\operatorname{inv})^*(j_*\mathcal{G})\) are lisse. Then
\[
\begin{aligned}
  &\mathrm{H}^2_c\bigl(G,\operatorname{Trans}_g^*(j_*\mathcal{F})
  \otimes(\operatorname{inv})^*(j_*\mathcal{G})\bigr)\\
  &\quad=
  \mathrm{H}^2_c\bigl(U_g,\operatorname{Trans}_g^*(\mathcal{F})
  \otimes(\operatorname{inv})^*(\mathcal{G})\bigr)\\
  &\quad=
  \mathrm{H}^0\bigl(U_g,
  \operatorname{Trans}_{g*}(\mathcal{F}^{\vee})
  \otimes(\operatorname{inv})^*(\mathcal{G}^{\vee})\bigr)^{\vee}(1)\\
  &\quad=
  \operatorname{Hom}_{U_g}\bigl(\operatorname{Trans}_g^*(\mathcal{F}),
  (\operatorname{inv})^*(\mathcal{G}^{\vee})\bigr)^{\vee}(1).
\end{aligned}
\]
Since both \(\mathcal{F}\) and \(\mathcal{G}\) are irreducible, this group
vanishes unless
\[
  \operatorname{Trans}_g^*(\mathcal{F})
  \cong(\operatorname{inv})^*(\mathcal{G}^{\vee})
  \quad\text{as lisse sheaves on }U_g,
\]
i.e., unless
\[
  \operatorname{Trans}_g^*(j_*\mathcal{F}[1])
  \cong(\operatorname{inv})^*(j_*\mathcal{G}^{\vee}[1])
  \quad\text{as perverse sheaves on }G,
\]
i.e., unless
\[
  \operatorname{Trans}_g^*(K)
  \cong(\operatorname{inv})^*(\mathbb{D}(L))
  \quad\text{as perverse sheaves on }G,
\]
in which case the group is one-dimensional. By the constructibility of
\(\mathcal{H}^0(K*_!L)\), either \(\mathcal{H}^0(K*_!L)\) is punctual, in
which case \(K*_!L\) is semiperverse, or there is a dense open set \(V\) in
\(G\) such that
\[
  \operatorname{Trans}_g^*(K)\cong(\operatorname{inv})^*(\mathbb{D}(L))
  \quad\text{for all }g\in V.
\]
In this last case, we argue as follows. Fix one such \(g_0\); then
\[
  \operatorname{Trans}_g^*(K)\cong\operatorname{Trans}_{g_0}^*(K)
  \quad\text{for all }g\in V.
\]
Then the isomorphism class of \(\operatorname{Trans}_{g_0}^*(K)\) is invariant
by translation by all \(g\) in \(V_1:=g_0^{-1}V\). But the set of \(g\) in
\(G(k)\) such that \(\operatorname{Trans}_g^*\) fixes any particular
isomorphism class (here that of \(\operatorname{Trans}_{g_0}^*(K)\)) is a
subgroup of \(G(k)\). This subgroup, for \(\operatorname{Trans}_{g_0}^*(K)\),
contains the open dense set \(V_1\), and hence contains \(V_1V_1=G\). Thus the
isomorphism class of \(\operatorname{Trans}_{g_0}^*(K)\), and hence of \(K\)
itself, is translation invariant.
\end{proof}

\sourceheading[cor:2.6.10]{Corollary 2.6.10}
Let \(k\) be an algebraically closed field. Suppose \(G\) is a connected smooth
affine \(k\)-group scheme of finite type, of relative dimension one. Let \(K\)
be any perverse irreducible on \(G\) whose isomorphism class is not
translation-invariant. Then \(K\) has property \(\mathcal{P}\).

\begin{proof}
By \hyperref[lem:2.6.9]{Lemma~2.6.9}, \(K\) has \(\mathcal{P}_!\). But the
isomorphism class of \(\mathbb{D}K\) is not translation-invariant, so
\(\mathbb{D}K\) also has \(\mathcal{P}_!\), whence \(K\) has
\(\mathcal{P}_*\).
\end{proof}

\sourceheading[lem:2.6.11]{Lemma 2.6.11}
Let \(k\) be an algebraically closed field. Suppose \(G\) is a connected smooth
affine \(k\)-group scheme of finite type, of relative dimension one. For any
perverse objects \(K\) and \(L\) on \(G\), the not-necessarily perverse object
\(K*_!L\) on \(G\) has \(\mathcal{H}^i(K*_!L)=0\) for \(i>0\).

\begin{proof}
By an obvious dévissage, we reduce first to the case when \(K\) is
perverse irreducible, and then further to the case where \(L\) is perverse
irreducible as well. If either \(K\) or \(L\) is punctual, i.e., a delta
function, then \(K*_!L\) is a translate of either \(L\) or \(K\), so is
perverse, and hence has \(\mathcal{H}^i(K*_!L)=0\) for \(i>0\). If not, then
\(K\cong\mathcal{F}[1]\) and \(L\cong\mathcal{G}[1]\) for sheaves
\(\mathcal{F}\) and \(\mathcal{G}\) on \(G\). In this case,
\[
  \mathcal{H}^i(K*_!L)
  =
  \mathrm{R}^{i+2}\operatorname{sum}_!
  (\operatorname{pr}_1^*\mathcal{F}\otimes\operatorname{pr}_2^*\mathcal{G})
\]
vanishes for \(i\geq1\) for dimension reasons: \(G\) has relative dimension
one, hence \(\mathrm{R}^i\operatorname{sum}_!(\text{any sheaf})=0\) for
\(i\geq3\).
\end{proof}

\sourceheading[lem:2.6.12]{Lemma 2.6.12}
Let \(k\) be an algebraically closed field. Suppose \(G\) is a connected smooth
affine \(k\)-group scheme of finite type, of relative dimension one. Let \(K\)
be perverse on \(G\). Then \(K\) has property \(\mathcal{P}_!\) if and only if
\(K*_!L\) is perverse for every perverse irreducible \(L\) on \(G\) whose
isomorphism class is translation-invariant.

\begin{proof}
The ``only if'' is obvious. Suppose now that \(K*_!L\) is perverse for every
perverse irreducible \(L\) on \(G\) whose isomorphism class is
translation-invariant. We wish to show that \(K\) has \(\mathcal{P}_!\). For
this, it suffices (by \hyperref[lem:2.6.8]{Lemma~2.6.8}) to show that \(K*_!M\) is perverse for every
perverse irreducible \(M\) on \(G\). If the isomorphism class of \(M\) is
translation invariant, we are given the perversity of \(K*_!M\) by hypothesis.
If the isomorphism class of \(M\) is not translation invariant, then by \hyperref[cor:2.6.10]{Corollary~2.6.10},
\(M\) itself has \(\mathcal{P}_!\), and hence \(K*_!M=M*_!K\) is perverse.
\end{proof}

\sourceheading[lem:2.6.13]{Lemma 2.6.13}
Let \(k\) be an algebraically closed field. Suppose \(G\) is a connected smooth
affine \(k\)-group scheme of finite type, of relative dimension one. Let \(L\)
be a perverse irreducible on \(G\) whose isomorphism class is
translation-invariant. If \(G\) is \(\mathbb{G}_m\), then \(L\) is the shifted
Kummer sheaf \(\mathcal{L}_\chi[1]\) for some character \(\chi\) of
\(\pi_1^{\mathrm{tame}}(\mathbb{G}_m)\). If \(G\) is \(\mathbb{A}^1\), and
\(\operatorname{char}(k)=0\), then \(L\) is
\(\overline{\mathbb{Q}}_\ell[1]\), while if \(\operatorname{char}(k)=p>0\),
then \(L\) is the shifted Artin--Schreier sheaf
\(\mathcal{L}_{\psi(\alpha x)}[1]\) for some \(\alpha\) in \(k\).

\begin{proof}
Because \(L\) is translation invariant, it cannot be punctual, so it must be of
the form \(\mathcal{L}[1]\) for an irreducible middle extension sheaf
\(\mathcal{L}\) on \(G\). By translation invariance, \(\mathcal{L}\) must be
lisse on all of \(G\). Hence \(\mathcal{L}\) is a lisse irreducible sheaf on
\(G\).

If \(\operatorname{char}(k)=0\), then
\(\pi_1(\mathbb{G}_m)=\pi_1^{\mathrm{tame}}(\mathbb{G}_m)\) is abelian, so
\(\mathcal{L}\) is an \(\mathcal{L}_\chi\), while
\(\pi_1(\mathbb{A}^1)=0\), so \(\mathcal{L}\) is
\(\overline{\mathbb{Q}}_\ell\).

If \(\operatorname{char}(k)=p>0\), and \(G\) is \(\mathbb{G}_m\), the
translation invariance of \(\mathcal{L}\) forces it to have
\(\operatorname{swan}_\infty(\mathcal{L})=0\) and
\(\operatorname{swan}_0(\mathcal{L})=0\) (cf. \cite[1.1]{Ver},
\cite[4.1.6]{Ka-GKM}), hence to be tame. As
\(\pi_1^{\mathrm{tame}}(\mathbb{G}_m)\) is abelian, \(\mathcal{L}\) is an
\(\mathcal{L}_\chi\), as required. Suppose now that \(G\) is
\(\mathbb{A}^1\). Then \(L\) is a perverse irreducible whose isomorphism class
is translation invariant. Its Fourier transform \(\operatorname{FT}_\psi(L)\)
is a perverse irreducible, say \(M\), on \(\mathbb{A}^1\), whose isomorphism
class is invariant under
\(M\mapsto M\otimes\mathcal{L}_{\psi(\beta x)}\) for all \(\beta\) in \(k\).
Look at the \(I(\infty)\)-representation \(M(\infty)\) of such an \(M\). From
the fact that the \(\mathcal{L}_{\psi(\beta x)}\) give (one for each value of
\(\beta\) in \(k\)) an infinity of distinct characters of \(I(\infty)\), we see
(cf. the appendix \hyperref[sec:2.18]{\S2.18} to this chapter) that \(M(\infty)=0\). Therefore \(M\)
is punctual. Being perverse irreducible, \(M\) must be a single delta function
\(\delta_\alpha\) for some \(\alpha\) in \(k\), whence \(L\) is
\(\mathcal{L}_{\psi(\alpha x)}[1]\), as required.
\end{proof}

\sourceheading[lem:2.6.14]{Lemma 2.6.14}
Let \(k\) be an algebraically closed field. Suppose \(G\) is a connected smooth
affine \(k\)-group scheme of finite type, of relative dimension one. Let \(K\)
be perverse on \(G\). Then \(K\) has property \(\mathcal{P}_!\) if and only if
\(\operatorname{Hom}_G(K,L)=0\) for every perverse irreducible \(L\) on \(G\)
whose isomorphism class is translation-invariant.

\begin{proof}
We have shown above that \(K\) has \(\mathcal{P}_!\) if and only if \(K*_!L\)
is semiperverse for every perverse irreducible \(L\) on \(G\) whose isomorphism
class is translation-invariant. Fix one such \(L\). We know by \hyperref[lem:2.6.11]{Lemma~2.6.11} that
\(\mathcal{H}^i(K*_!L)=0\) for \(i>0\). Thus we must show that
\(\mathcal{H}^0(K*_!L)\) is punctual if and only if
\(\operatorname{Hom}_G(K,L)=0\). Because the isomorphism class of \(L\) is
translation invariant, at any point \(g\) in \(G(k)\), the stalk of
\(\mathcal{H}^0(K*_!L)\) is the same cohomology group
\[
  \mathrm{H}^0_c(G,\operatorname{inv}^*L\otimes K).
\]
Thus
\[
  \mathcal{H}^0(K*_!L)\text{ is punctual}
  \Longleftrightarrow
  \mathcal{H}^0(K*_!L)=0
  \Longleftrightarrow
  \mathrm{H}^0_c(G,\operatorname{inv}^*L\otimes K)=0.
\]
By \hyperref[lem:2.6.13]{Lemma~2.6.13}, \(L=\mathcal{L}[1]\) with \(\mathcal{L}\) lisse of rank one on
\(G\), and \(\mathcal{L}^{\vee}\cong\operatorname{inv}^*\mathcal{L}\). Thus,
ignoring Tate twists we find
\[
\begin{aligned}
  \mathrm{H}^0_c(G,\operatorname{inv}^*L\otimes K)
  &=
  \mathrm{H}^1_c(G,\operatorname{inv}^*\mathcal{L}\otimes K)
  \quad\text{is dual to}\\
  \mathrm{H}^{-1}(G,\operatorname{inv}^*\mathcal{L}^{\vee}\otimes\mathbb{D}(K))
  &=
  \mathrm{H}^0(G,\operatorname{inv}^*\mathcal{L}^{\vee}[-1]\otimes\mathbb{D}(K))\\
  &=
  \operatorname{Hom}(\overline{\mathbb{Q}}_\ell,
  \operatorname{inv}^*\mathcal{L}^{\vee}[-1]\otimes\mathbb{D}(K))\\
  &=
  \operatorname{Hom}(\operatorname{inv}^*\mathcal{L}[1],\mathbb{D}(K))
  \quad(\text{because }\mathcal{L}\text{ is lisse of rank one})\\
  &=
  \operatorname{Hom}(\mathbb{D}(L),\mathbb{D}(K))
  \quad(\text{because }\operatorname{inv}^*\mathcal{L}\cong\mathcal{L}^{\vee})\\
  &\cong \operatorname{Hom}(K,L).
\end{aligned}
\]
Thus we find that \(\mathcal{H}^0(K*_!L)\) is punctual if and only if
\(\operatorname{Hom}(K,L)=0\).
\end{proof}

\sourceheading[cor:2.6.15]{Corollary 2.6.15}
Let \(k\) be an algebraically closed field. Suppose \(G\) is a connected smooth
affine \(k\)-group scheme of finite type, of relative dimension one. Let \(K\)
be perverse on \(G\). Then \(K\) has property \(\mathcal{P}_*\) if and only if
\(\operatorname{Hom}_G(L,K)=0\) for every perverse irreducible \(L\) on \(G\)
whose isomorphism class is translation-invariant.

\begin{proof}
Indeed, \(K\) has \(\mathcal{P}_*\) if and only if \(\mathbb{D}K\) has
\(\mathcal{P}_!\), if and only if
\(\operatorname{Hom}(\mathbb{D}K,L)=0\) for every perverse irreducible \(L\) on
\(G\) whose isomorphism class is translation-invariant, if and only if
\(\operatorname{Hom}(\mathbb{D}L,K)=0\) for every such \(L\). But the class of
such \(L\) is stable by duality.
\end{proof}

\sourceheading[cor:2.6.16]{Corollary 2.6.16}
Let \(k\) be an algebraically closed field. Suppose \(G\) is a connected smooth
affine \(k\)-group scheme of finite type, of relative dimension one. Let \(K\)
be perverse on \(G\). If \(K\) has \(\mathcal{P}_!\), then every quotient of
\(K\) (as perverse sheaf) has \(\mathcal{P}_!\). If \(K\) has
\(\mathcal{P}_*\), then every subobject of \(K\) (as perverse sheaf) has
\(\mathcal{P}_*\).

\begin{proof}
Obvious from the \(\operatorname{Hom}(K,L)=0\) and
\(\operatorname{Hom}(L,K)=0\) criteria \hyperref[lem:2.6.14]{Lemma~2.6.14} and \hyperref[cor:2.6.15]{Corollary~2.6.15} for \(K\) to have
\(\mathcal{P}_!\) and \(\mathcal{P}_*\) respectively.
\end{proof}

\sourceheading[cor:2.6.16.1]{Corollary 2.6.16.1}
Let \(k\) be an algebraically closed field. Suppose \(G\) is a connected smooth
affine \(k\)-group scheme of finite type, of relative dimension one. Suppose
that
\[
  0\to K_1\to K\to K_2\to0
\]
is a short exact sequence of perverse sheaves on \(G\). If both \(K_1\) and
\(K_2\) have \(\mathcal{P}_!\) (respectively \(\mathcal{P}_*\), respectively
\(\mathcal{P}\)), then so does \(K\).

\begin{proof}
Again obvious from the \(\operatorname{Hom}(K,L)=0\) and
\(\operatorname{Hom}(L,K)=0\) criteria for \(K\) to have \(\mathcal{P}_!\) and
\(\mathcal{P}_*\) respectively.
\end{proof}

\sourceheading[cor:2.6.16.2]{Corollary 2.6.16.2}
Let \(k\) be an algebraically closed field. Suppose \(G\) is a connected smooth
affine \(k\)-group scheme of finite type, of relative dimension one. If \(K\)
is perverse on \(G\) and has \(\mathcal{P}_!\) (respectively
\(\mathcal{P}_*\)) then any perverse irreducible quotient (respectively
subobject) \(M\) of \(K\) as perverse sheaf has \(\mathcal{P}\).

\begin{proof}
Indeed by \hyperref[cor:2.6.16]{Corollary~2.6.16}, \(M\) has \(\mathcal{P}_!\) (respectively
\(\mathcal{P}_*\)). But \(M\) is itself perverse irreducible, hence by \hyperref[lem:2.6.14]{Lemma~2.6.14}
(respectively by \hyperref[cor:2.6.15]{Corollary~2.6.15}), its isomorphism class cannot be translation
invariant. So by \hyperref[cor:2.6.10]{Corollary~2.6.10}, \(M\) has \(\mathcal{P}\).
\end{proof}

\sourceheading[cor:2.6.16.3]{Corollary 2.6.16.3}
Let \(k\) be an algebraically closed field. Suppose \(G\) is a connected smooth
affine \(k\)-group scheme of finite type, of relative dimension one. If \(K\)
is perverse on \(G\) and has \(\mathcal{P}\), then every perverse irreducible
quotient and every perverse irreducible subobject of \(K\) as perverse sheaf
has \(\mathcal{P}\).

\begin{proof}
Immediate from \hyperref[cor:2.6.16.2]{Corollary~2.6.16.2}.
\end{proof}

\sourceheading[rem:2.6.16.4]{Remark 2.6.16.4}
It is not necessarily the case that if \(K\) is perverse and has
\(\mathcal{P}\), then every subobject and every quotient of \(K\) as perverse
sheaf again has \(\mathcal{P}\). To give examples, we work over \(\mathbb{C}\)
(see \hyperref[rem:2.10.4]{Remark~2.10.4} for examples in characteristic \(p>0\)), and use
Riemann--Hilbert to work with holonomic RS \(\mathcal{D}\)-modules instead of
perverse sheaves.

On \(\mathbb{G}_{a,\mathbb{C}}:=\operatorname{Spec}(\mathbb{C}[x])\), we
write \(\mathcal{D}\) for \(\mathbb{C}[x,\delta]\), where
\(\delta:=d/dx\). Consider the left \(\mathcal{D}\)-module
\(M:=\mathcal{D}/\mathcal{D}x\delta_x\). Because \(x\delta_x\) is minus its
own adjoint, \(M\) is self dual. The only irreducible holonomic
\(\mathcal{D}\)-modules whose isomorphism classes are translation invariant
are \(\mathcal{D}/\mathcal{D}(\delta-\alpha)\), for \(\alpha\) in
\(\mathbb{C}\). Only for \(\alpha=0\) is
\(\mathcal{D}/\mathcal{D}(\delta-\alpha)\) RS. Therefore \(M\) has
\(\mathcal{P}\) if and only if
\[
  \operatorname{Hom}_{\text{left-}\mathcal{D}\text{-modules}}
  (M,\mathcal{D}/\mathcal{D}\delta)=0.
\]
We will show that for any \(\alpha\) in \(\mathbb{C}\), we have
\[
  \operatorname{Hom}_{\text{left-}\mathcal{D}\text{-modules}}
  (M,\mathcal{D}/\mathcal{D}(\delta-\alpha))=0.
\]
This \(\operatorname{Hom}\) is, via the image of the class of \(1\) in
\(\mathcal{D}/\mathcal{D}x\delta_x\), just the set of elements in
\(\mathcal{D}/\mathcal{D}(\delta-\alpha)\) annihilated by \(x\delta_x\). But
\(\mathcal{D}/\mathcal{D}(\delta-\alpha)\cong e^{\alpha x}\mathbb{C}[x]\), via
\(1\mapsto e^{\alpha x}\), and
\(e^{\alpha x}\mathbb{C}[x]\subset\mathbb{C}[[x]]\). But \(x\delta_x\) is
obviously (look at lowest degree terms) injective on \(\mathbb{C}[[x]]\), and
hence \(M\) has \(\mathcal{P}\). But \(M\) admits
\(\mathcal{D}/\mathcal{D}\delta_x\) as quotient, which does not have
\(\mathcal{P}\), because
\(\mathcal{D}/\mathcal{D}\delta\subset\mathcal{D}/\mathcal{D}\delta_x\) via
\(\operatorname{Right}(x)\). Similarly, \(M\) admits, via
\(\operatorname{Right}(x)\), \(\mathcal{D}/\mathcal{D}x\delta\) as a subobject,
which again does not have \(\mathcal{P}\), because
\(\mathcal{D}/\mathcal{D}x\delta\) admits \(\mathcal{D}/\mathcal{D}\delta\) as
a quotient.

On \(\mathbb{G}_{m,\mathbb{C}}:=\operatorname{Spec}(\mathbb{C}[x,x^{-1}])\),
we write \(\mathcal{D}\) for \(\mathbb{C}[x,x^{-1},\delta]\), \(\delta\) again
given by \(\delta:=d/dx\). Consider the left \(\mathcal{D}\)-module
\[
  N:=\mathcal{D}/\mathcal{D}(x-1)\delta(x-1).
\]
Because \((x-1)\delta(x-1)\) is minus its own adjoint, \(N\) is self dual. The
irreducible holonomic \(\mathcal{D}\)-modules whose isomorphism classes are
translation invariant are those of the form
\(\mathcal{D}/\mathcal{D}(x\delta-\alpha)\), for \(\alpha\) in
\(\mathbb{C}\). Thus \(N\) has \(\mathcal{P}\) if and only if, for every
\(\alpha\) in \(\mathbb{C}\), we have
\[
  \operatorname{Hom}_{\text{left-}\mathcal{D}\text{-modules}}
  (N,\mathcal{D}/\mathcal{D}(x\delta-\alpha))=0.
\]
Just as above, this \(\operatorname{Hom}\) is the set of elements in
\(\mathcal{D}/\mathcal{D}(x\delta-\alpha)\cong x^\alpha\mathbb{C}[x,x^{-1}]\)
which are annihilated by \((x-1)\delta(x-1)\). In terms of the parameter
\(t:=x-1\), \(x^\alpha\mathbb{C}[x,x^{-1}]\subset\mathbb{C}[[t]]\), and
\((x-1)\delta(x-1)\) becomes the operator \(t(d/dt)t\), which, just as above,
is injective on \(\mathbb{C}[[t]]\). Therefore \(N\) has \(\mathcal{P}\). But
\(N\) admits \(\mathcal{D}/\mathcal{D}\delta(x-1)\) as a quotient, which does
not have \(\mathcal{P}\); indeed, since \(x\) is invertible, we have
\(\mathcal{D}/\mathcal{D}x\delta=\mathcal{D}/\mathcal{D}\delta\), and
\(\mathcal{D}/\mathcal{D}\delta\subset\mathcal{D}/\mathcal{D}\delta(x-1)\) via
\(\operatorname{Right}(x-1)\). Similarly, \(N\) admits, via
\(\operatorname{Right}(x-1)\), \(\mathcal{D}/\mathcal{D}(x-1)\delta\) as a
subobject, which does not have \(\mathcal{P}\), since it in turn admits
\(\mathcal{D}/\mathcal{D}\delta=\mathcal{D}/\mathcal{D}x\delta\) as a quotient.

\sourceheading[cor:2.6.17]{(Key) Corollary 2.6.17}
Let \(k\) be an algebraically closed field. Suppose \(G\) is a connected smooth
affine \(k\)-group scheme of finite type, of relative dimension one. Suppose
that \(K\) and \(L\) are perverse on \(G\), and each has \(\mathcal{P}\). Then
their middle convolution \(K*_{\mathrm{mid}}L\) has \(\mathcal{P}\).

\begin{proof}
We have already noted that \(K*_!L\) has \(\mathcal{P}_!\), and that
\(K*_*L\) has \(\mathcal{P}_*\), in both cases simply by the associativity of
\(!\) and \(*\) convolution. But \(K*_{\mathrm{mid}}L\) is a quotient of
\(K*_!L\), and hence has \(\mathcal{P}_!\). Similarly,
\(K*_{\mathrm{mid}}L\) is a subobject of \(K*_*L\), and hence has
\(\mathcal{P}_*\).
\end{proof}

\sourcepara{2.6.18}
By the above corollary, when \(G\) is either \(\mathbb{G}_m\) or
\(\mathbb{A}^1\), the full subcategory \((\mathcal{P})\) of the category
\(\operatorname{Perv}\) of all perverse sheaves on \(G\) consisting of those
with property \(\mathcal{P}\) is stable by middle convolution
\(*_{\mathrm{mid}}\). This category for \(\mathbb{A}^1\), with its middle
convolution, will be the essential player in all that follows.

\section{Interlude: middle direct images (relative dimension one)}\label{sec:2.7}

\sourcepara{2.7.1}
Here is the general setup for ``middle direct image''.

\sourceheading[prop:2.7.2]{Proposition 2.7.2}
Over an arbitrary base \(S\) which is itself separated and of finite type over
a field \(k\) of characteristic \(p\neq\ell\), consider a diagram
\[
\begin{tikzcd}[column sep=large,row sep=large]
U \arrow[r,"j"] \arrow[dr,"f"'] &
X \arrow[d,"\bar{f}"] &
D:=X-U \arrow[l,"i"'] \arrow[dl,"\bar{f}|D"] \\
& S &
\end{tikzcd}
\]
in which \(j:U\to X\) is an affine open immersion, \(\bar{f}\) is proper, and
\(\bar{f}|D:D\to S\) is affine (hence finite, since it is also proper).
Suppose \(K\) in \(\mathrm{D}^b_c(U,\overline{\mathbb{Q}}_\ell)\) is perverse,
and that both \(\mathrm{R}f_*K\) and \(\mathrm{R}f_!K\) are perverse on \(S\).
Then \(\mathrm{R}\bar{f}_*j_{!*}K\) is perverse on \(S\), and
\[
  \mathrm{R}\bar{f}_*j_{!*}K
  =
  \operatorname{Image}(\mathrm{R}f_!K\to \mathrm{R}f_*K)
\]
in \(\operatorname{Perv}(S)\).

\begin{proof}
On \(X\), we have two short exact sequences of perverse sheaves:
\[
  \begin{aligned}
    0 &\to \operatorname{Ker}\to j_!K\to j_{!*}K\to 0,\\
    0 &\to j_{!*}K\to \mathrm{R}j_*K\to \operatorname{Coker}\to 0.
  \end{aligned}
\]
Notice that the objects \(\operatorname{Ker}\) and \(\operatorname{Coker}\)
are supported on \(D\), hence are perverse on \(D\).

In the derived category, these are distinguished triangles, so applying the
exact functor \(\mathrm{R}\bar{f}_*\) gives two distinguished triangles on
\(S\),
\[
  \begin{aligned}
    &\to \mathrm{R}\bar{f}_*\operatorname{Ker}
      \to \mathrm{R}f_!K
      \to \mathrm{R}\bar{f}_*j_{!*}K\to,\\
    &\to \mathrm{R}\bar{f}_*j_{!*}K
      \to \mathrm{R}f_*K
      \to \mathrm{R}\bar{f}_*\operatorname{Coker}\to.
  \end{aligned}
\]
The objects \(\mathrm{R}\bar{f}_*\operatorname{Ker}\) and
\(\mathrm{R}\bar{f}_*\operatorname{Coker}\) are perverse on \(S\), because
they are both of the form \(\mathrm{R}(\bar{f}|D)_*(\text{perverse on }D)\),
and \(\bar{f}|D\), being finite, preserves perversity.

We first show that \(\mathrm{R}\bar{f}_*j_{!*}K\) is perverse. The first
distinguished triangle above shows that \(\mathrm{R}\bar{f}_*j_{!*}K\) is
semiperverse: its \(\mathcal{H}^{-i}\) is caught between the
\(\mathcal{H}^{-i}\) of one perverse sheaf and the \(\mathcal{H}^{-(i-1)}\)
of another, so certainly has dimension of support \(\leq i\). The hypothesis
that both \(\mathrm{R}f_*K\) and \(\mathrm{R}f_!K\) are perverse gives, by
duality, that also both \(\mathrm{R}f_*\mathbb{D}K\) and
\(\mathrm{R}f_!\mathbb{D}K\) are perverse. So repeating the argument we find
that
\[
  \mathrm{R}\bar{f}_*j_{!*}\mathbb{D}K
  =
  \mathbb{D}(\mathrm{R}\bar{f}_*j_{!*}K)
\]
is also semiperverse, whence \(\mathrm{R}\bar{f}_*j_{!*}K\) is perverse.

Once we know that \(\mathrm{R}\bar{f}_*j_{!*}K\) is perverse, the two
distinguished triangles above are actually short exact sequences in
\(\operatorname{Perv}(S)\)
\[
  \begin{aligned}
    0 &\to \mathrm{R}\bar{f}_*\operatorname{Ker}
      \to \mathrm{R}f_!K
      \to \mathrm{R}\bar{f}_*j_{!*}K\to 0,\\
    0 &\to \mathrm{R}\bar{f}_*j_{!*}K
      \to \mathrm{R}f_*K
      \to \mathrm{R}\bar{f}_*\operatorname{Coker}\to 0,
  \end{aligned}
\]
which together show that
\[
  \mathrm{R}\bar{f}_*j_{!*}K
  =
  \operatorname{Image}(\mathrm{R}f_!K\to \mathrm{R}f_*K).
\]
\end{proof}

\section{Middle additive convolution via middle direct image}
\label{sec:2.8}

\sourcepara{2.8.1}
We return to \(\mathbb{A}^1\) over an algebraically closed field \(k\) of
characteristic \(\neq\ell\). We will apply the idea of ``middle direct image''
to computing the middle additive convolution with a perverse object \(L\) on
\(\mathbb{A}^1\) which has property \(\mathcal{P}\).

\sourcepara{2.8.2}
Given a perverse sheaf \(K\) on \(\mathbb{A}^1\), we have the perverse sheaf
\(K_x\otimes L_{t-x}\) on \(\mathbb{A}^2\), with coordinates \(x,t\). In these
coordinates, the original ``sum'' map is the projection onto the \(t\)-line. By
the very definition of additive convolution, we have
\[
  \begin{aligned}
    K*_!+L &= \mathrm{R}(\operatorname{pr}_2)_!(K_x\otimes L_{t-x}),\\
    K*_*+L &= \mathrm{R}(\operatorname{pr}_2)_*(K_x\otimes L_{t-x}).
  \end{aligned}
\]
Since \(L\) lies in \((\mathcal{P})\), both of these objects are perverse.

\sourcepara{2.8.3}
We compactify the map \(\operatorname{pr}_2\) by compactifying the affine
\(x\)-line into the projective \(x\)-line, all over the \(\mathbb{A}^1\) of
\(t\)'s:
\[
  j:\mathbb{A}^1_x\times\mathbb{A}^1_t
  \hookrightarrow
  \mathbb{P}^1_x\times\mathbb{A}^1_t.
\]
Then we are in the set-up to which the above middle direct image proposition
applies. Thus we find

\sourceheading[prop:2.8.4]{Proposition 2.8.4}
Hypotheses and notations as in \hyperref[par:2.8.1]{(2.8.1)}, for \(K\) perverse on
\(\mathbb{A}^1\), and \(L\) perverse on \(\mathbb{A}^1\) with
\(\mathcal{P}\), the middle additive convolution \(K*_{\mathrm{mid}}+L\) is
\[
  \mathrm{R}(\operatorname{pr}_2)_*
  \bigl(\text{the perverse sheaf } j_{!*}(K_x\otimes L_{t-x})
  \text{ on } \mathbb{P}^1\times\mathbb{A}^1\bigr).
\]

Using this, we get a fibre-by-fibre recipe for \(K*_{\mathrm{mid}}+L\) over a
dense open set.

\sourceheading[cor:2.8.5]{Corollary 2.8.5}
Hypotheses as in \hyperref[prop:2.8.4]{Proposition~2.8.4} above, suppose in addition that
\(K=\mathcal{F}[1]\) and \(L=\mathcal{G}[1]\) for sheaves \(\mathcal{F}\) and
\(\mathcal{G}\) on \(\mathbb{A}^1\). Then
\begin{enumerate}[label=(\arabic*)]
\item there exists a dense open set \(U\subset\mathbb{A}^1\) over which
\begin{enumerate}[label=(\alph*)]
\item the formation of
\(\mathrm{R}j_*(\mathcal{F}_x\otimes\mathcal{G}_{t-x})\mid
\mathbb{P}^1\times U\) commutes with arbitrary change of base on \(U\),
\item \(\mathrm{R}j_*(\mathcal{F}_x\otimes\mathcal{G}_{t-x})\mid
\infty\times U\) is lisse on \(U\), of formation compatible with arbitrary
change of base on \(U\).
\end{enumerate}
\item Over \(\mathbb{P}^1\times U\), we have
\begin{enumerate}[label=(\alph*)]
\item an isomorphism
\[
  j_{!*}(K_x\otimes L_{t-x})\mid\mathbb{P}^1\times U
  \cong
  j_*(\mathcal{F}_x\otimes\mathcal{G}_{t-x})[2]
  \mid\mathbb{P}^1\times U,
\]
\item a short exact sequence of perverse sheaves
\[
  \begin{aligned}
    0 &\to
    j_*(\mathcal{F}_x\otimes\mathcal{G}_{t-x})[1]\mid\infty\times U\\
    &\to
    j_!(\mathcal{F}_x\otimes\mathcal{G}_{t-x})[2]
    \mid\mathbb{P}^1\times U\\
    &\to
    j_*(\mathcal{F}_x\otimes\mathcal{G}_{t-x})[2]
    \mid\mathbb{P}^1\times U
    \to 0,
  \end{aligned}
\]
which, as a distinguished triangle, is of formation compatible with arbitrary
change of base on \(U\).
\end{enumerate}
\item For \(E/k\) any separably closed extension field of \(k\), and any point
\(t\) in \(U(E)\), we have
\[
  (K*_{\mathrm{mid}}L)_t
  =
  \mathrm{R}\Gamma\bigl(\mathbb{P}^1\otimes_k E,
  (j:\mathbb{A}^1\to\mathbb{P}^1)_*
  (\text{the sheaf }x\mapsto\mathcal{F}_x\otimes\mathcal{G}_{t-x})\bigr)[2],
\]
and a ``short exact sequence'' (really a distinguished triangle)
\[
  \begin{aligned}
    0 &\to
    (\text{the }I(\infty)\text{-invariants in the sheaf }
    x\mapsto \mathcal{F}_x\otimes\mathcal{G}_{t-x})[1]\\
    &\to (K*_!L)_t
    \to (K*_{\mathrm{mid}}L)_t
    \to 0,
  \end{aligned}
\]
which is the fibre at \(t\) of a short exact sequence of perverse sheaves on
\(U\)
\[
  0\to
  (\text{the lisse sheaf }j_*(\mathcal{F}_x\otimes\mathcal{G}_{t-x})
  \mid\infty\times U)[1]
  \to K*_!+L
  \to K*_{\mathrm{mid}}+L
  \to 0.
\]
\item If both \(\mathcal{F}\) and \(\mathcal{G}\) are tame at \(\infty\), then
we may take \(U\) to be \(\mathbb{A}^1\), and the sheaf
\(j_*(\mathcal{F}_x\otimes\mathcal{G}_{t-x})\mid
\infty\times\mathbb{A}^1\) is lisse, of formation compatible with arbitrary
change of base on \(\mathbb{A}^1\).
\item If \(\operatorname{char}(k)=0\), then the sheaf
\(j_*(\mathcal{F}_x\otimes\mathcal{G}_{t-x})\mid
\infty\times\mathbb{A}^1\) is geometrically constant, of formation compatible
with arbitrary change of base on \(\mathbb{A}^1\).
\end{enumerate}

\begin{proof}
(1) The base-changing statements over a dense open \(U_1\) are a special case
of Deligne's generic base change theorem \cite[Cor.~2.9]{De-ThFin}.
The lisseness in (1b) results from the constructibility of
\(\mathrm{R}j_*(\mathcal{F}_x\otimes\mathcal{G}_{t-x})\), but may require
passing to a smaller dense open \(U\) in \(U_1\). Once we have (1), (2a)
follows from the ``successive partial truncation'' description
\cite[2.2.4]{BBD}, and (2b) follows from (1). Once we
have (2), (3) results from (proper base change and) the fact that, by (1), the
formation of \(j_*(\mathcal{F}_x\otimes\mathcal{G}_{t-x})\mid
\mathbb{P}^1\times U\) commutes with arbitrary change of base on \(U\). To
prove (4), notice that the sheaf \((\mathcal{F}_x\otimes\mathcal{G}_{t-x})\)
on \(\mathbb{P}^1\times\mathbb{A}^1-\infty\times\mathbb{A}^1\) is lisse in a
Zariski open neighborhood of \(\infty\times\mathbb{A}^1\): indeed, if \(f(x)\)
and \(g(x)\) are monic polynomials such that \(\mathcal{F}\) (resp.
\(\mathcal{G}\)) is lisse where \(f\) (resp. \(g\)) is invertible, then
\(\mathcal{F}_x\otimes\mathcal{G}_{t-x}\) is lisse where the
\((\pm)\)monic-in-\(x\) polynomial \(f(x)g(t-x)\) is invertible. Moreover,
this sheaf is fibre-by-fibre tamely ramified along
\(\infty\times\mathbb{A}^1\). The result (4) then follows from the relative
Abhyankar lemma, cf. \cite[4.7.2]{Ka-SE}. If
\(\operatorname{char}(k)=0\), then \(\mathcal{F}\) and \(\mathcal{G}\) are
automatically tame at \(\infty\), so by (4) the sheaf
\(j_*(\mathcal{F}_x\otimes\mathcal{G}_{t-x})\mid\infty\times\mathbb{A}^1\) is
lisse; because we are in characteristic zero, any lisse sheaf on
\(\mathbb{A}^1\) is geometrically constant.
\end{proof}

\section{Middle additive convolution with Kummer sheaves}\label{sec:2.9}

\sourcepara{2.9.1}
We continue to work over an algebraically closed field of characteristic
\(\neq\ell\). We now apply \hyperref[prop:2.8.4]{Proposition~2.8.4} to the special case of computing the middle
additive convolution with \(j_!\mathcal{L}_\chi[1]=j_*\mathcal{L}_\chi[1]\) on
\(\mathbb{A}^1\), for \(\mathcal{L}_\chi\) a nontrivial Kummer sheaf on
\(\mathbb{G}_m\), and \(j:\mathbb{G}_m\to\mathbb{A}^1\) the inclusion. We know
that \(j_*\mathcal{L}_\chi[1]\) on \(\mathbb{A}^1\) is perverse irreducible,
and its isomorphism class is not translation invariant (it has a unique point
where it is not lisse). So \(j_*\mathcal{L}_\chi[1]\) on \(\mathbb{A}^1\) has
property \(\mathcal{P}\).

Exactly as in \hyperref[prop:2.8.4]{Proposition~2.8.4} we find

\sourceheading[prop:2.9.2]{Proposition 2.9.2}
Hypotheses and notations as in \hyperref[par:2.9.1]{(2.9.1)}, for \(K\) perverse on
\(\mathbb{A}^1\), and any nontrivial Kummer sheaf \(\mathcal{L}_\chi\), the
middle additive convolution \(K*_{\mathrm{mid}}+j_*\mathcal{L}_\chi[1]\) is
\[
  \mathrm{R}(\operatorname{pr}_2)_*
  \bigl(\text{the perverse sheaf }
  j_{!*}(K_x\otimes\mathcal{L}_\chi(t-x)[1])
  \text{ on }\mathbb{P}^1\times\mathbb{A}^1\bigr).
\]

\sourcepara{2.9.3}
In order to go further, we need to see what the perverse sheaf
\(j_{!*}(K_x\otimes\mathcal{L}_\chi(t-x)[1])\) on
\(\mathbb{P}^1\times\mathbb{A}^1\) looks like. To lighten the notational
burden in the following lemma, we will sometimes write \(\mathcal{L}_\chi\) on
\(\mathbb{A}^1\) to mean \(j_*\mathcal{L}_\chi\) on \(\mathbb{A}^1\).

\sourceheading[lem:2.9.4]{Lemma 2.9.4}
Hypotheses and notations as in \hyperref[prop:2.9.2]{Proposition~2.9.2} above:
\begin{enumerate}[label=(\arabic*)]
\item For any \(t_0\) in \(\mathbb{A}^1(k)\), the restriction of
\(j_{!*}(K_x\otimes\mathcal{L}_\chi(t-x)[1])\) to
\(\mathbb{P}^1\times\{t_0\}\) is
\[
  (j_{t_0}:\mathbb{A}^1_x\times\{t_0\}\to
  \mathbb{P}^1_x\times\{t_0\})_{!*}
  (K_x\otimes\mathcal{L}_\chi(t_0-x))[1].
\]
\item Denoting by \(i:\infty\times\mathbb{A}^1\to
\mathbb{P}^1\times\mathbb{A}^1\) the inclusion, the perverse sheaf
\(j_{!*}(K_x\otimes\mathcal{L}_\chi(t-x)[1])\) on
\(\mathbb{P}^1\times\mathbb{A}^1\) sits in a short exact sequence of perverse
sheaves on \(\mathbb{P}^1\times\mathbb{A}^1\)
\[
  \begin{aligned}
    0 &\to
    i_*\bigl(\text{the constant sheaf }
    (\mathcal{H}^{-1}(K)\otimes\mathcal{L}_\chi)^{I(\infty)}
    \text{ on }\infty\times\mathbb{A}^1_t\bigr)[1]\\
    &\to j_!(K_x\otimes\mathcal{L}_\chi(t-x)[1])
    \to j_{!*}(K_x\otimes\mathcal{L}_\chi(t-x)[1])
    \to 0.
  \end{aligned}
\]
\item The middle and \(!\) additive convolutions of \(K\) with
\(j_*\mathcal{L}_\chi[1]\) sit in an exact sequence of perverse sheaves on
\(\mathbb{A}^1\)
\[
  \begin{aligned}
    0 &\to
    \bigl(\text{the constant sheaf }
    (\mathcal{H}^{-1}(K)\otimes\mathcal{L}_\chi)^{I(\infty)}\bigr)[1]
    \text{ on }\mathbb{A}^1_t\\
    &\to K*_!j_*\mathcal{L}_\chi[1]
    \to K*_{\mathrm{mid}}+j_*\mathcal{L}_\chi[1]
    \to 0.
  \end{aligned}
\]
\end{enumerate}

\begin{proof}
An affine open neighborhood of \(\infty\times\mathbb{A}^1\) is the
\(\mathbb{A}^2\) with coordinates \(z\) and \(t\), \(z=1/x\). In these
coordinates, \(\infty\times\mathbb{A}^1\) is defined by \(z=0\), and so the
function \(1-zt\) takes the value \(1\) along \(\infty\times\mathbb{A}^1\). Let
us denote by \(U\) the open in \(\mathbb{A}^2\) where \(1-tz\) is invertible.
On this open set, we have the lisse, rank one sheaf
\(\mathcal{L}_\chi(1-tz)\). Thanks to the identity
\[
  \begin{aligned}
    t-x &= (-1/z)(1-zt),\\
    \mathcal{L}_\chi(t-x)
      &= \mathcal{L}_\chi(1-tz)\otimes\mathcal{L}_\chi(-1/z).
  \end{aligned}
\]
In terms of the inclusion
\[
  j:U-\infty\times\mathbb{A}^1\to U,
\]
we have
\[
  \begin{aligned}
    &K_x\otimes\mathcal{L}_\chi(t-x)[1]
    \mid U-\infty\times\mathbb{A}^1\\
    &\quad =
    (j^*\mathcal{L}_\chi(1-tz))
    \otimes(K_{1/z}\otimes\mathcal{L}_\chi(-1/z))[1]
    \text{ on }\mathbb{G}_{m,z}\times\mathbb{A}^1_t,
    \mid U-\infty\times\mathbb{A}^1.
  \end{aligned}
\]
The key point here is that \(\mathcal{L}_\chi(1-tz)\) is lisse on all of
\(U\). So taking middle extension across \(\infty\times\mathbb{A}^1\), which
we may compute Zariski locally along \(\infty\times\mathbb{A}^1\), we see that
on \(U\) we will get
\[
  \mathcal{L}_\chi(1-tz)\otimes
  \bigl(\text{M.E. of }(K_{1/z}\otimes\mathcal{L}_\chi(-1/z))[1]
  \text{ on }\mathbb{G}_{m,z}\times\mathbb{A}^1_t,
  \text{ across }z=0\bigr).
\]
Since formation of middle extension is compatible with external products
(this being true separately for duality, for \(j_!\) and (so) for its dual
\(\mathrm{R}j_*\)), we see that, on \(U\), our middle extension is
the restriction to \(U\) of the external tensor product sheaf
\[
  \begin{aligned}
    &\mathcal{L}_\chi(1-tz)\otimes
    \bigl(\text{M.E. of }(K_{1/z}\otimes\mathcal{L}_\chi(-1/z))
    \mid\mathbb{G}_{m,z}\\
    &\qquad\text{ across }z=0\bigr)
    \otimes(\overline{\mathbb{Q}}_\ell[1]\text{ on }\mathbb{A}^1_t).
  \end{aligned}
\]
on the \(\mathbb{A}^2\) with coordinates \(z,t\). In particular, this shows
that (1) holds. If we go back to the original \(x,t\) coordinates, we find
that the middle extension across \(\infty\times\mathbb{A}^1\) is, along
\(\infty\times\mathbb{A}^1\), the constant perverse sheaf \([1]\) on
\(\mathbb{A}^1_t\) which is
\[
  (K\otimes\mathcal{L}_\chi)^{I(\infty)}[1]
  =
  (\mathcal{H}^{-1}(K)\otimes\mathcal{L}_\chi)^{I(\infty)}[2].
\]
So on \(\mathbb{P}^1_x\times\mathbb{A}^1_t\), we have, denoting by
\(i:\infty\times\mathbb{A}^1\to\mathbb{P}^1\times\mathbb{A}^1\) the
inclusion, the short exact sequence of perverse sheaves
\[
  \begin{aligned}
    0 &\to
    i_*\bigl(\text{the constant sheaf }
    (\mathcal{H}^{-1}(K)\otimes\mathcal{L}_\chi)^{I(\infty)}
    \text{ on }\infty\times\mathbb{A}^1_t\bigr)[1]\\
    &\to j_!(K_x\otimes\mathcal{L}_\chi(t-x)[1])
    \to j_{!*}(K_x\otimes\mathcal{L}_\chi(t-x)[1])
    \to 0
  \end{aligned}
\]
asserted in (2). Taking the total direct image onto the \(t\)-line, we get the
exact sequence
\[
  \begin{aligned}
    0 &\to
    \bigl(\text{the constant sheaf }
    (\mathcal{H}^{-1}(K)\otimes\mathcal{L}_\chi)^{I(\infty)}\bigr)[1]
    \text{ on }\mathbb{A}^1_t\\
    &\to K*_!j_*\mathcal{L}_\chi[1]
    \to K*_{\mathrm{mid}}+j_*\mathcal{L}_\chi[1]
    \to 0.
  \end{aligned}
\]
\end{proof}

\sourceheading[cor:2.9.5]{Corollary 2.9.5}
For \(K\) perverse on \(\mathbb{A}^1\), and any nontrivial Kummer sheaf
\(\mathcal{L}_\chi\), the stalk at any point \(t\) in \(\mathbb{A}^1(k)\) of
\(K*_{\mathrm{mid}}+j_*\mathcal{L}_\chi[1]\) is
\[
  \mathrm{R}\Gamma\bigl(\mathbb{P}^1,
  j_{!*}(\text{the perverse sheaf }
  x\mapsto K_x\otimes\mathcal{L}_\chi(t-x)
  \text{ on }\mathbb{A}^1)[1]\bigr).
\]

\begin{proof}
This results from the preceding \hyperref[prop:2.9.2]{Proposition~2.9.2}, via proper base change and
part (1) of the above \hyperref[lem:2.9.4]{Lemma~2.9.4}.
\end{proof}

\sourceheading[prop:2.9.6]{Proposition 2.9.6}
Let \(k\) be an algebraically closed field. Given any two nontrivial Kummer
sheaves \(\mathcal{L}_\rho\) and \(\mathcal{L}_\chi\) on
\(\mathbb{G}_{m,k}\), their middle additive convolution
\(j_*\mathcal{L}_\rho[1]*_{\mathrm{mid}}+j_*\mathcal{L}_\chi[1]\) on
\(\mathbb{A}^1\) is given (geometrically)
\[
  \begin{aligned}
    j_*\mathcal{L}_\rho[1]*_{\mathrm{mid}}+j_*\mathcal{L}_\chi[1]
      &\cong j_*\mathcal{L}_{\rho\chi}[1],
      &&\text{if }\rho\chi\neq 1,\\
      &\cong \delta_0,
      &&\text{if }\chi\rho=1.
  \end{aligned}
\]

\begin{proof}
Suppose first that \(\chi\rho\neq 1\). Then by part (3) of \hyperref[lem:2.9.4]{Lemma~2.9.4} above,
applied with \(K:=j_*\mathcal{L}_\rho[1]\), we see that
\[
  (j_*\mathcal{L}_\rho[1])*_!+(j_*\mathcal{L}_\chi[1])
  \cong
  j_*\mathcal{L}_\rho[1]*_{\mathrm{mid}}+j_*\mathcal{L}_\chi[1],
\]
because
\[
  (\mathcal{H}^{-1}(K)\otimes\mathcal{L}_\chi)^{I(\infty)}
  =
  (\mathcal{L}_{\rho\chi})^{I(\infty)}
  =
  0.
\]
Thus our middle convolution is given by
\[
  \mathrm{R}(\operatorname{pr}_2)_!
  \bigl(\text{the perverse sheaf }
  \mathcal{L}_\rho(x)\otimes\mathcal{L}_\chi(t-x)[2]
  \text{ on }\mathbb{A}^1\times\mathbb{A}^1\bigr).
\]
Over \(t=0\), the (geometric) stalk is
\(\mathrm{R}\Gamma_c(\mathbb{G}_m,\mathcal{L}_{\rho\chi}[2])=0\). Therefore
our middle convolution on \(\mathbb{A}^1\) is the extension by zero of its
restriction to \(\mathbb{G}_m\). So it suffices to show that over
\(\mathbb{G}_m\), our middle convolution is (geometrically) isomorphic to
\(\mathcal{L}_{\rho\chi}[1]\).

Over the open set \(\mathbb{G}_m\) where \(t\) is invertible, we make the
change of variable
\[
  (x,t)\mapsto (tx,t),
\]
and we find
\[
  \begin{aligned}
    &j_*\mathcal{L}_\rho[1]*_{\mathrm{mid}}+j_*\mathcal{L}_\chi[1]
      \mid \mathbb{G}_m\\
    &\quad\cong
    \mathrm{R}(\operatorname{pr}_2)_!
    \bigl(\mathcal{L}_\rho(tx)\otimes\mathcal{L}_\chi(t-tx)[2]
    \text{ on }\mathbb{A}^1\times\mathbb{G}_m\bigr)\\
    &\quad=
    \mathrm{R}(\operatorname{pr}_2)_!
    \bigl(\mathcal{L}_\rho(x)\otimes\mathcal{L}_\chi(1-x)[1]
    \otimes\mathcal{L}_{\rho\chi}(t)[1]
    \text{ on }\mathbb{A}^1\times\mathbb{G}_m\bigr)\\
    &\quad=
    \mathcal{L}_{\rho\chi}[1]\otimes
    \mathrm{R}\Gamma_c(\mathbb{A}^1-\{1,0\},
    \mathcal{L}_\rho(x)\otimes\mathcal{L}_\chi(1-x)[1]).
  \end{aligned}
\]
Because \(\rho\chi\neq 1\), we have
\(\mathrm{H}^2_c(\mathbb{A}^1-\{1,0\},
\mathcal{L}_\rho(x)\otimes\mathcal{L}_\chi(1-x))=0\), by considering the
local monodromy at \(\infty\) of
\(\mathcal{L}_\rho(x)\otimes\mathcal{L}_\chi(1-x)\); because
\(\mathcal{L}_\rho(x)\otimes\mathcal{L}_\chi(1-x)\) is tame, and lisse of
rank one, we have
\[
  \chi_c(\mathbb{A}^1-\{1,0\},
  \mathcal{L}_\rho(x)\otimes\mathcal{L}_\chi(1-x))=-1,
\]
and the only possibly nonvanishing groups are the \(\mathrm{H}^i_c\) for
\(i=1,2\). Thus
\[
  \mathrm{R}\Gamma_c(\mathbb{A}^1-\{1,0\},
  \mathcal{L}_\rho(x)\otimes\mathcal{L}_\chi(1-x)[1])
\]
is a one-dimensional \(\overline{\mathbb{Q}}_\ell\)-vector space, placed in
degree zero, as required.

Suppose now that \(\chi\rho=1\). In this case, \hyperref[cor:2.9.5]{Corollary~2.9.5} above will
give the required assertion. For the stalk of our middle convolution at \(t\)
in \(\mathbb{A}^1(k)\) is given by
\[
  \mathrm{R}\Gamma\bigl(\mathbb{P}^1,
  j_{!*}(\text{the perverse sheaf }
  x\mapsto\mathcal{L}_\chi(x)\otimes\mathcal{L}_\rho(t-x)[1]
  \text{ on }\mathbb{A}^1)[1]\bigr).
\]
By definition, the sheaf
\(\mathcal{L}_\chi(x)\otimes\mathcal{L}_\rho(t-x)[1]\) on \(\mathbb{A}^1\) is
lisse outside of \(\{t,0\}\), and is extended by zero across \(\{t,0\}\). So
if we denote by
\[
  k_{t,0}:\mathbb{A}^1-\{t,0\}\to\mathbb{A}^1
\]
the inclusion, this stalk is
\[
  \mathrm{R}\Gamma\bigl(\mathbb{P}^1,
  j_{!*}(k_{t,0})_!(\mathcal{L}_\chi(x)\otimes
  \mathcal{L}_\rho(t-x)[1]\text{ on }\mathbb{A}^1_x-\{t,0\})[1]\bigr).
\]
Since \(\chi\rho=1\), this stalk is
\[
  \mathrm{R}\Gamma\bigl(\mathbb{P}^1,
  j_{!*}(k_{t,0})_!(\mathcal{L}_\chi(x/(t-x))[1]
  \text{ on }\mathbb{A}^1_x-\{t,0\})[1]\bigr).
\]
For any \(t\), this coefficient sheaf is lisse of rank one at \(\infty\), so
it is lisse of rank one on \(\mathbb{P}^1_x-\{t,0\}\) and everywhere tame.

For \(t\neq 0\), the sheaf \(\mathcal{L}_\chi(x/(t-x))\) is not geometrically
constant (it has nontrivial monodromy at both \(0\) and \(t\)), and has no
nonzero punctual sections, so
\[
  \mathrm{H}^i\bigl(\mathbb{P}^1,
  j_{!*}(k_{t,0})_!(\mathcal{L}_\chi(x/(t-x))[1]
  \text{ on }\mathbb{A}^1_x-\{t,0\})[1]\bigr)=0
\]
if \(i=0,-2\). But the Euler characteristic also vanishes, so we find
\[
  \mathrm{R}\Gamma\bigl(\mathbb{P}^1,
  j_{!*}(k_{t,0})_!(\mathcal{L}_\chi(x/(t-x))[1]
  \text{ on }\mathbb{A}^1_x-\{t,0\})[1]\bigr)=0
\]
for \(t\neq 0\), and hence
\[
  j_*\mathcal{L}_\rho[1]*_{\mathrm{mid}}+j_*\mathcal{L}_\chi[1]
\]
is supported at \(t=0\) if \(\rho\chi=1\). So it remains only to compute the
stalk at zero. This stalk is
\[
  \mathrm{R}\Gamma\bigl(\mathbb{P}^1,
  j_{!*}(k_{0,0})_!(\mathcal{L}_\chi(x/(-x))[1]
  \text{ on }\mathbb{A}^1_x-\{0\})[1]\bigr).
\]
But the perverse sheaf
\(j_{!*}(k_{0,0})_!(\mathcal{L}_\chi(x/(-x))[1])\) is (geometrically) just
the constant sheaf \(\overline{\mathbb{Q}}_\ell[1]\) on
\(\mathbb{P}^1-\{0\}\), extended by zero. So our stalk at zero is
\[
  \mathrm{R}\Gamma_c(\mathbb{P}^1-\{0\},\overline{\mathbb{Q}}_\ell[2])
  =
  \mathrm{R}\Gamma_c(\mathbb{A}^1,\overline{\mathbb{Q}}_\ell[2]),
\]
the one-dimensional vector space
\(\mathrm{H}^2_c(\mathbb{A}^1,\overline{\mathbb{Q}}_\ell)\) placed in degree
zero.
\end{proof}

\sourceheading[thm:2.9.7]{Theorem 2.9.7}
Let \(k\) be an algebraically closed field, \(\mathcal{L}_\chi\) a nontrivial
Kummer sheaf on \(\mathbb{G}_m\), \(j:\mathbb{G}_m\to\mathbb{A}^1\) the
inclusion.
\begin{enumerate}[label=(\arabic*)]
\item On \(\mathbb{A}^1\), the operators
\(K\mapsto K*_{\mathrm{mid}}+j_*\mathcal{L}_\chi[1]\) and
\(K\mapsto K*_{\mathrm{mid}}+j_*\mathcal{L}_{\bar{\chi}}[1]\) on
\(\mathcal{P}\) are automorphisms of \(\mathcal{P}\), which are inverses of
each other (where we write \(\bar{\chi}\) for the character \(\chi^{-1}\)).
\item On \(\mathbb{A}^1\), the operator
\(K\mapsto K*_{\mathrm{mid}}+j_*\mathcal{L}_\chi[1]\) induces an automorphism
of the perverse irreducible objects in \(\mathcal{P}\), with inverse the
operator \(K\mapsto K*_{\mathrm{mid}}+j_*\mathcal{L}_{\bar{\chi}}[1]\).
\end{enumerate}

\begin{proof}
We know that \(j_*\mathcal{L}_\chi[1]\) is perverse irreducible on
\(\mathbb{A}^1\), and that it lies in \(\mathcal{P}\) on \(\mathbb{A}^1\)
(because it is not translation invariant). In view of the preceding \hyperref[prop:2.9.6]{Proposition~2.9.6}, \hyperref[thm:2.9.7]{Theorem~2.9.7} is a special case of the following result.
\end{proof}

\sourceheading[thm:2.9.8]{Theorem 2.9.8}
On \(\mathbb{A}^1\) over an algebraically closed field \(k\), let \(K\) and
\(L\) be perverse objects in \(\mathcal{P}\) with
\(K*_{\mathrm{mid}}+L=\delta_0=L*_{\mathrm{mid}}+K\).
\begin{enumerate}[label=(\arabic*)]
\item On \(\mathbb{A}^1\), the operators \(X\mapsto X*_{\mathrm{mid}}+K\) and
\(X\mapsto X*_{\mathrm{mid}}+L\) on \(\mathcal{P}\) are automorphisms of
\(\mathcal{P}\), which are inverses of each other.
\item The operator \(X\mapsto X*_{\mathrm{mid}}+K\) induces an automorphism of
the perverse irreducible objects in \(\mathcal{P}\), with inverse the operator
\(X\mapsto X*_{\mathrm{mid}}+L\).
\end{enumerate}

\begin{proof}
For (1), just use the fact that middle convolution is associative, and that
\(\delta_0\) is the identity for middle convolution.

For (2), let \(X\) be perverse irreducible in \(\mathcal{P}\). We must show
that the perverse sheaf \(X*_{\mathrm{mid}}+K\) is perverse irreducible. By
(1), it is nonzero, so it contains a subobject \(Y\) which is perverse
irreducible. Since \(X*_{\mathrm{mid}}+K\) lies in \(\mathcal{P}\), \(Y\)
itself lies in \(\mathcal{P}\) (by \hyperref[cor:2.6.16.3]{Corollary~2.6.16.3}). Since
\(W\mapsto W*_{\mathrm{mid}}+L\) is end-exact, it preserves injections, so
\(Y*_{\mathrm{mid}}+L\) is a subobject of
\(X*_{\mathrm{mid}}+K*_{\mathrm{mid}}+L=X\). This subobject is nonzero by
(1), so by the irreducibility of \(X\), we have
\[
  X=Y*_{\mathrm{mid}}+L.
\]
Applying \(W\mapsto W*_{\mathrm{mid}}+K\), we find
\[
  X*_{\mathrm{mid}}+K
  =
  Y*_{\mathrm{mid}}+L*_{\mathrm{mid}}+K
  =
  Y
\]
is irreducible, as required.
\end{proof}

\noindent
For later use, we record the following corollary.

\sourceheading[cor:2.9.9]{Corollary 2.9.9}
On \(\mathbb{A}^1\) over an algebraically closed field \(k\), let \(K\) and
\(L\) be perverse objects in \(\mathcal{P}\) with
\(K*_{\mathrm{mid}}+L=\delta_0=L*_{\mathrm{mid}}+K\). Then \(L\) and \(K\)
are perverse irreducible.

\begin{proof}
Indeed, by (2), the functors \(X\mapsto X*_{\mathrm{mid}}+K\) and
\(X\mapsto X*_{\mathrm{mid}}+L\) carry irreducibles in \(\mathcal{P}\) to
irreducibles in \(\mathcal{P}\). Take \(X\) to be \(\delta_0\).
\end{proof}

\section[Middle additive convolution via Fourier transform]{Interpretation of middle additive convolution via Fourier Transform}
\label{sec:2.10}

\sourcepara{2.10.1}
In this section, we work on \(\mathbb{A}^1\) over an algebraically closed field
\(k\) of characteristic \(p>0\), \(p\neq\ell\). Recall \cite[Cor.~9.6]{Br}
that Fourier Transform interchanges \(!\) convolution and tensor product:
more precisely, for \(K\) and \(L\) in
\(\mathrm{D}^b_c(\mathbb{A}^1,\overline{\mathbb{Q}}_\ell)\), we have
\[
  \begin{aligned}
    \operatorname{FT}_{\psi,!}(K*_!L)
      &=\operatorname{FT}_{\psi,!}(K)\otimes
        \operatorname{FT}_{\psi,!}(L)[-1],\\
    \operatorname{FT}_{\psi,!}(K\otimes L)[-1]
      &=\operatorname{FT}_{\psi,!}(K)*_!
        \operatorname{FT}_{\psi,!}(L)(-1).
  \end{aligned}
\]

\sourceheading[lem:2.10.2]{Lemma 2.10.2}
In the situation of \hyperref[par:2.10.1]{(2.10.1)}, let \(K\) be perverse on \(\mathbb{A}^1\). The
following conditions are equivalent.
\begin{enumerate}[label=(\arabic*)]
\item \(K\) has \(\mathcal{P}_!\).
\item \(K\) has no quotient \(\mathcal{L}_\psi(\alpha x)[1]\) for any
\(\alpha\) in \(k\).
\item \(\operatorname{FT}(K)\) has no nonzero punctual quotient.
\item \(\mathcal{H}^0(\operatorname{FT}(K))=0\).
\item \(\operatorname{FT}(K)\) is of the form \(\mathcal{F}[1]\) for some
sheaf \(\mathcal{F}\) on \(\mathbb{A}^1\) which has no nonzero punctual
sections.
\end{enumerate}

\begin{proof}
(1) \(\Leftrightarrow\) (2) by \hyperref[lem:2.6.13]{Lemma~2.6.13} and \hyperref[lem:2.6.14]{Lemma~2.6.14}.

(2) \(\Leftrightarrow\) (3) since \(\operatorname{FT}\) carries
\(\mathcal{L}_\psi(\alpha x)[1]\) to \(\delta_{-\alpha}\).

(3) \(\Rightarrow\) (4): For any perverse \(N\) on \(\mathbb{A}^1\), the
short exact sequence of perverse sheaves
\[
  0\to\mathcal{H}^{-1}(N)[1]\to N\to\mathcal{H}^0(N)\to 0
\]
expresses \(\mathcal{H}^0(N)\) as a punctual quotient of \(N\). Applying this
to \(\operatorname{FT}(K)\) gives
\(\mathcal{H}^0(\operatorname{FT}(K))=0\).

(4) \(\Leftrightarrow\) (5) by the concrete description of perverse sheaves in
dimension one.

(5) \(\Rightarrow\) (1): If \(\operatorname{FT}(K)\) is \(\mathcal{F}[1]\) for
some sheaf \(\mathcal{F}\) on \(\mathbb{A}^1\), then
\[
  L\mapsto \operatorname{FT}(K)[-1]\otimes L=\mathcal{F}\otimes L
\]
preserves semiperversity. By Fourier inversion,
\(L\mapsto K*_!L\) preserves semiperversity. By \hyperref[lem:2.6.7]{Lemma~2.6.7}, \(K\) has
\(\mathcal{P}_!\).
\end{proof}

\sourceheading[lem:2.10.3]{Lemma 2.10.3}
In the situation \hyperref[par:2.10.1]{(2.10.1)}, for \(K\) perverse on \(\mathbb{A}^1\), \(K\) has
\(\mathcal{P}\) if and only if \(\operatorname{FT}(K):=N\) is a middle
extension.

\begin{proof}
If \(K\) has \(\mathcal{P}\), then both \(K\) and \(\mathbb{D}K\) have
\(\mathcal{P}_!\). By the above lemma, both
\[
  N:=\operatorname{FT}_{\psi,!}K
  \quad\text{and}\quad
  \mathbb{D}N=\operatorname{FT}_{\bar{\psi},!}(\mathbb{D}K)
\]
have no nonzero punctual quotient. Thus (by duality) \(N\) has no nonzero
punctual subobject. By the same lemma, we know \(N\) is \(\mathcal{F}[1]\) for
some sheaf \(\mathcal{F}\) on \(\mathbb{A}^1\) which has no nonzero punctual
sections. This forces \(N\) to be a middle extension. For if
\(N\) is \(\mathcal{F}[1]\), with
\(\mathcal{F}\subset j_*j^*\mathcal{F}\) for \(j:U\to\mathbb{A}^1\) an open
where \(\mathcal{F}\) is lisse, then if
\(\mathcal{F}\neq j_*j^*\mathcal{F}\), the short exact sequence of perverse
sheaves
\[
  0\to (j_*j^*\mathcal{F})/\mathcal{F}
  \to \mathcal{F}[1]=N
  \to j_*j^*\mathcal{F}[1]\to 0
\]
exhibits a nonzero punctual subobject of \(N\).

Conversely, suppose that \(N:=\operatorname{FT}_{\psi,!}(K)\) is a middle
extension. Then so is
\(\mathbb{D}N=\operatorname{FT}_{\bar{\psi},!}(\mathbb{D}K)\). By the previous
proposition, both \(K\) and \(\mathbb{D}K\) have \(\mathcal{P}_!\), so \(K\)
has \(\mathcal{P}\).
\end{proof}

\sourceheading[rem:2.10.4]{Remarks 2.10.4}
\begin{enumerate}[label=(\arabic*)]
\item Here is a slightly variant proof of the above result. We know that a
perverse \(K\) lies in \(\mathcal{P}\) if and only if as a perverse sheaf it
has no subobject and no quotient of the form \(\mathcal{L}_\psi(\alpha x)[1]\)
for any \(\alpha\) in \(k\). By Fourier Transform, this becomes the condition
that \(\operatorname{FT}(K)\) as perverse sheaf have no subobject and no
quotient which is punctual. But this last condition is equivalent to being a
middle extension (compare \cite[2.9.1]{Ka-ESDE}), as one sees using the
natural filtration (cf. \hyperref[par:2.3.6]{(2.3.6)}) of the perverse sheaf \(\operatorname{FT}(K)\)
with associated graded (punctual, middle extension, punctual).
\item Using \hyperref[lem:2.10.3]{Lemma~2.10.3}, we can give a characteristic \(p>0\) example (compare
\hyperref[rem:2.6.16.4]{Remark~2.6.16.4}) of a perverse \(K\) on \(\mathbb{A}^1\) which has \(\mathcal{P}\),
such that \(K\) admits a quotient as perverse sheaf which does not have
\(\mathcal{P}\). Thanks to \hyperref[lem:2.10.3]{Lemma~2.10.3}, it is the same to give an example of a
perverse \(N\) on \(\mathbb{A}^1\) which is a middle extension, but which
admits a quotient as perverse sheaf which is not a middle extension. On
\(\mathbb{G}_m\) over our algebraically closed field \(k\) of characteristic
\(p\neq\ell\), there exists a lisse rank two
\(\overline{\mathbb{Q}}_\ell\)-sheaf \(\mathcal{F}\) which is a nontrivial
extension of the constant sheaf \(\overline{\mathbb{Q}}_\ell\) by itself,
because
\[
  \mathrm{H}^1(\mathbb{G}_m,\overline{\mathbb{Q}}_\ell)
  =
  \overline{\mathbb{Q}}_\ell.
\]
This sheaf is automatically tame, its global monodromy being unipotent, and
hence its local monodromy at zero is a single Jordan block, of dimension two.
Denote by \(j:\mathbb{G}_m\to\mathbb{A}^1\) the inclusion. Applying \(j_*\) to
the tautological short exact sequence of sheaves on \(\mathbb{G}_m\)
\[
  0\to\overline{\mathbb{Q}}_\ell\to\mathcal{F}\to
  \overline{\mathbb{Q}}_\ell\to 0,
\]
we get an exact sequence of sheaves on \(\mathbb{A}^1\),
\[
  0\to\overline{\mathbb{Q}}_{\ell,\mathbb{A}^1}
  \to j_*\mathcal{F}
  \to\overline{\mathbb{Q}}_{\ell,\mathbb{A}^1}.
\]
The last arrow is not surjective on the stalk at zero, since \(j_*\mathcal{F}\)
has only a one-dimensional stalk at zero, so we have a short exact sequence of
sheaves on \(\mathbb{A}^1\)
\[
  0\to\overline{\mathbb{Q}}_{\ell,\mathbb{A}^1}
  \to j_*\mathcal{F}
  \to j_!\overline{\mathbb{Q}}_{\ell,\mathbb{G}_m}
  \to 0.
\]
Shifting by \([1]\), we get a short exact sequence of perverse sheaves on
\(\mathbb{A}^1\) which exhibits
\(j_!\overline{\mathbb{Q}}_{\ell,\mathbb{G}_m}[1]\), which is visibly not a
middle extension, as a quotient of the middle extension
\((j_*\mathcal{F})[1]=j_{!*}(\mathcal{F}[1])\).
\end{enumerate}

\sourceheading[prop:2.10.5]{Proposition 2.10.5}
Let \(K\) be perverse on \(\mathbb{A}^1\), and suppose \(K\) has
\(\mathcal{P}\). Let \(\mathcal{L}_\chi\) be a nontrivial Kummer sheaf on
\(\mathbb{G}_m\), \(j:\mathbb{G}_m\to\mathbb{A}^1\) the inclusion. Write
\(\operatorname{FT}(K)=\mathcal{F}[1]\), with \(\mathcal{F}\) a middle
extension sheaf on \(\mathbb{A}^1\). Then
\[
  \operatorname{FT}(K*_{\mathrm{mid}}+j_*\mathcal{L}_\chi[1])
  =
  j_*(j^*\operatorname{FT}(K)\otimes\mathcal{L}_{\bar{\chi}})
  =
  j_*(j^*\mathcal{F}\otimes\mathcal{L}_{\bar{\chi}})[1].
\]

\begin{proof}
We have already seen in \hyperref[lem:2.9.4]{Lemma~2.9.4} that we have a short exact sequence of perverse
sheaves on \(\mathbb{A}^1\)
\[
  0\to(\text{constant sheaf})[1]\to
  K*_!j_*\mathcal{L}_\chi[1]
  \to K*_{\mathrm{mid}}+j_*\mathcal{L}_\chi[1]\to 0.
\]
Under Fourier Transform, this gives a short exact sequence
\[
  0\to(\text{pct. sheaf, supp. at }0)
  \to\operatorname{FT}(K*_!j_*\mathcal{L}_\chi[1])
  \to\operatorname{FT}(K*_{\mathrm{mid}}+j_*\mathcal{L}_\chi[1])\to 0.
\]
Using the general identity
\[
  \operatorname{FT}(K*_!L)=\operatorname{FT}(K)\otimes\operatorname{FT}(L)[-1],
\]
together with the standard geometric isomorphism
\[
  \operatorname{FT}(j_*\mathcal{L}_\chi[1])
  \cong
  j_*\mathcal{L}_{\bar{\chi}}[1],
\]
we rewrite this short exact sequence as
\[
  0\to(\text{pct. sheaf, supp. at }0)
  \to \operatorname{FT}(K)\otimes j_*\mathcal{L}_{\bar{\chi}}
  \to \operatorname{FT}(K*_{\mathrm{mid}}+j_*\mathcal{L}_\chi[1])\to 0.
\]
Restricting to \(\mathbb{G}_m\), the above exact sequence gives an
isomorphism:
\[
  j^*\operatorname{FT}(K)\otimes\mathcal{L}_{\bar{\chi}}
  \cong
  j^*\operatorname{FT}(K*_{\mathrm{mid}}+j_*\mathcal{L}_\chi[1]).
\]
Since we know a priori that
\(\operatorname{FT}(K*_{\mathrm{mid}}+j_*\mathcal{L}_\chi[1])\) is a middle
extension, being \(\operatorname{FT}\) of an object in \((\mathcal{P})\), we
have
\[
  \begin{aligned}
    \operatorname{FT}(K*_{\mathrm{mid}}+j_*\mathcal{L}_\chi[1])
      &\cong
      j_*j^*\operatorname{FT}(K*_{\mathrm{mid}}+j_*\mathcal{L}_\chi[1])\\
      &=
      j_*(j^*\operatorname{FT}(K)\otimes\mathcal{L}_{\bar{\chi}}),
  \end{aligned}
\]
as required.
\end{proof}

\sourceheading[rem:2.10.6]{Remark 2.10.6}
This proposition gives us, in positive characteristic, a second proof of
\hyperref[prop:2.9.6]{Proposition~2.9.6} (take \(K=j_*\mathcal{L}_\rho[1]\)).

\sourcepara{2.10.7}
In fact, more is true. Under Fourier Transform, the operation
\(*_{\mathrm{mid}}+\) corresponds to the obvious tensor product operation on
middle extensions.

\sourceheading[thm:2.10.8]{Theorem 2.10.8}
In the situation \hyperref[par:2.10.1]{(2.10.1)}, let \(K\) and \(L\) be
perverse on \(\mathbb{A}^1\), both \(K\) and \(L\) having \(\mathcal{P}\). Pick
a common open set \(j:U\to\mathbb{A}^1\) where both \(N:=\operatorname{FT}(K)\)
and \(M:=\operatorname{FT}(L)\) are lisse, and write
\(N=j_*\mathcal{F}[1]\), \(M=j_*\mathcal{G}[1]\), with \(\mathcal{F}\) and
\(\mathcal{G}\) lisse sheaves on \(U\). Then
\[
  \operatorname{FT}(K*_{\mathrm{mid}}+L)=j_*(\mathcal{F}\otimes\mathcal{G})[1].
\]

\begin{proof}
The key point is to show that there exists a dense open set
\(j:U\to\mathbb{A}^1\) such that when we apply \(\operatorname{FT}\) to the
``forget supports'' map \(K*_!+L\to K*_*+L\), the map we obtain,
\[
  \operatorname{FT}(K*_!+L)\to\operatorname{FT}(K*_*+L),
\]
is an isomorphism on \(U\):
\[
  j^*\operatorname{FT}(K*_!+L)\cong j^*\operatorname{FT}(K*_*+L).
\]
If this is so, then from
\[
  \operatorname{FT}(K*_{\mathrm{mid}}+L)
    =\operatorname{Image}\bigl(\operatorname{FT}(K*_!+L)\to
      \operatorname{FT}(K*_*+L)\bigr),
\]
we get
\[
  j^*\operatorname{FT}(K*_{\mathrm{mid}}+L)
    \cong j^*\operatorname{FT}(K*_!+L)
    \cong j^*\operatorname{FT}(K*_*+L).
\]
Since we know a priori that \(\operatorname{FT}(K*_{\mathrm{mid}}+L)\) is a
middle extension, we know that
\[
  \operatorname{FT}(K*_{\mathrm{mid}}+L)
    \cong j_*j^*\operatorname{FT}(K*_{\mathrm{mid}}+L)
    \cong j_*j^*\operatorname{FT}(K*_!+L).
\]
Using
\(\operatorname{FT}(K*_!+L)=\operatorname{FT}(K)\otimes\operatorname{FT}(L)[-1]
=\mathcal{F}\otimes\mathcal{G}[1]\) gives the assertion.

It remains to prove that the map
\[
  \operatorname{FT}(K*_!+L)\to\operatorname{FT}(K*_*+L)
\]
is an isomorphism on a dense open set of \(\mathbb{A}^1\).
\end{proof}

\sourceheading[lem:2.10.9]{Lemma 2.10.9}
In the situation \hyperref[par:2.10.1]{(2.10.1)}, for any two objects \(K\) and
\(L\) in \(\mathrm{D}^b_c(\mathbb{A}^1,\overline{\mathbb{Q}}_\ell)\), the
natural map
\[
  \operatorname{FT}(\text{``forget supports''}):
  \operatorname{FT}(K*_!+L)\to\operatorname{FT}(K*_*+L)
\]
is an isomorphism on a dense open set.

\begin{proof}
The idea of the proof is to exploit the basic miracle of Fourier Transform,
that \(\operatorname{FT}_{!,\psi}\cong\operatorname{FT}_{*,\psi}\). By
Deligne's generic base change theorem \cite[Cor.~2.9]{De-ThFin}, there exists a
dense open set \(U\) in \(\mathbb{A}^1\) over which the formation of each of
\(\operatorname{FT}_{*,\psi}(K)\), \(\operatorname{FT}_{*,\psi}(L)\), and
\(\operatorname{FT}_{*,\psi}(K*_*+L)\) commutes with arbitrary change of base.
By proper base change, the formation of each
\(\operatorname{FT}_{!,\psi}(K)\), \(\operatorname{FT}_{!,\psi}(L)\), and
\(\operatorname{FT}_{!,\psi}(K*_!+L)\) commutes with arbitrary change of base
over all of \(\mathbb{A}^1\).

Choose any point \(\alpha\) in \(\mathbb{A}^1(k)\). The stalk at \(\alpha\) of
\(\operatorname{FT}(K*_!+L)\), viewed as an \(\operatorname{FT}_{!,\psi}\), is
\[
  \mathrm{R}\Gamma_c(\mathbb{A}^1,(K*_!+L)\otimes\mathcal{L}_\psi(\alpha x))[1].
\]
If \(\alpha\) lies in \(U\), the stalk at \(\alpha\) of
\(\operatorname{FT}(K*_*+L)\), viewed as an \(\operatorname{FT}_{*,\psi}\), is
\[
  \mathrm{R}\Gamma(\mathbb{A}^1,(K*_*+L)\otimes\mathcal{L}_\psi(\alpha x))[1],
\]
and the map between them is induced by the ``forget supports'' map
\(K*_!+L\to K*_*+L\) on coefficients, followed by the ``forget supports'' map
\(\mathrm{R}\Gamma_c\to\mathrm{R}\Gamma\).

The source is
\[
  \begin{aligned}
  &\mathrm{R}\Gamma_c(\mathbb{A}^1,(K*_!+L)\otimes\mathcal{L}_\psi(\alpha x))[1]\\
  &\quad=
  \mathrm{R}\Gamma_c(\mathbb{A}^1,
    \mathrm{R}(\operatorname{sum})_!(\operatorname{pr}_1^*K\otimes
    \operatorname{pr}_2^*L)\otimes\mathcal{L}_\psi(\alpha x))[1]\\
  &\quad=
  \mathrm{R}\Gamma_c(\mathbb{A}^1,
    \mathrm{R}(\operatorname{sum})_!(\operatorname{pr}_1^*(K\otimes
    \mathcal{L}_\psi(\alpha x))\otimes
    \operatorname{pr}_2^*(L\otimes\mathcal{L}_\psi(\alpha x))))[1]\\
  &\qquad\text{(by the projection formula, and the additivity of \(\mathcal{L}_\psi\))}\\
  &\quad=
  \mathrm{R}\Gamma_c(\mathbb{A}^2,
    \operatorname{pr}_1^*(K\otimes\mathcal{L}_\psi(\alpha x))\otimes
    \operatorname{pr}_2^*(L\otimes\mathcal{L}_\psi(\alpha x)))[1]\quad
    \text{(by Leray)}\\
  &\quad=
  \mathrm{R}\Gamma_c(\mathbb{A}^1,K\otimes\mathcal{L}_\psi(\alpha x))
  \otimes
  \mathrm{R}\Gamma_c(\mathbb{A}^1,L\otimes\mathcal{L}_\psi(\alpha x))[1]
  \quad\text{(by Kunneth).}
  \end{aligned}
\]
In completely analogous fashion, the target, for \(\alpha\) in \(U\), is
\[
  \mathrm{R}\Gamma(\mathbb{A}^1,K\otimes\mathcal{L}_\psi(\alpha x))\otimes
  \mathrm{R}\Gamma(\mathbb{A}^1,L\otimes\mathcal{L}_\psi(\alpha x))[1],
\]
and, with these identifications, the map between them is
\((\text{``forget supports''})\otimes(\text{``forget supports''})\).
Because \(\alpha\) lies in \(U\), over which the formation of both
\(\operatorname{FT}_{*,\psi}(K)\) and \(\operatorname{FT}_{*,\psi}(L)\)
commutes with base change, these ``forget supports'' maps
\[
  \begin{aligned}
  \mathrm{R}\Gamma_c(\mathbb{A}^1,K\otimes\mathcal{L}_\psi(\alpha x))
    &\to \mathrm{R}\Gamma(\mathbb{A}^1,K\otimes\mathcal{L}_\psi(\alpha x)),\\
  \mathrm{R}\Gamma_c(\mathbb{A}^1,L\otimes\mathcal{L}_\psi(\alpha x))
    &\to \mathrm{R}\Gamma(\mathbb{A}^1,L\otimes\mathcal{L}_\psi(\alpha x)),
  \end{aligned}
\]
are just the identity maps of \(\operatorname{FT}(K)_\alpha[-1]\) and of
\(\operatorname{FT}(L)_\alpha[-1]\), respectively.

Thus the map
\(\operatorname{FT}(K*_!+L)\to\operatorname{FT}(K*_*+L)\) is an isomorphism
over \(U\), as required.
\end{proof}

\sourceheading[cor:2.10.10]{Corollary 2.10.10}
In the situation \hyperref[par:2.10.1]{(2.10.1)}, for any two objects \(K\) and
\(L\) in \(\mathrm{D}^b_c(\mathbb{A}^1,\overline{\mathbb{Q}}_\ell)\), the
mapping cone object \([K*_!+L\to K*_*+L]\) in
\(\mathrm{D}^b_c(\mathbb{A}^1,\overline{\mathbb{Q}}_\ell)\) is a direct sum of
shifted \(\mathcal{L}_\psi(\alpha x)\) sheaves, for various \(\alpha\) in \(k\).

\begin{proof}
This is the Fourier Transform of the statement that the mapping cone
\([\operatorname{FT}(K*_!+L)\to\operatorname{FT}(K*_*+L)]\) is punctual.
\end{proof}

\sourceheading[rem:2.10.11]{Remark 2.10.11}
In terms of a relative compactification of \(\operatorname{pr}_2\), say
\[
  \begin{tikzcd}
  \mathbb{A}^1\times\mathbb{A}^1
    \arrow[r,"j"] \arrow[dr,"\operatorname{pr}_2"']
    & \mathbb{P}^1\times\mathbb{A}^1
    & \infty\times\mathbb{A}^1 \arrow[l,"i"'] \arrow[dl,"\operatorname{pr}_2|_\infty"]\\
  & \mathbb{A}^1,
  \end{tikzcd}
\]
the object \([K*_!+L\to K*_*+L]\) is
\(i^*\mathrm{R}j_*(K_x\otimes L_{t-x})\). So the above
\hyperref[cor:2.10.10]{Corollary~2.10.10} says that, over an algebraically
closed field of characteristic \(p>0\),
\(i^*\mathrm{R}j_*(K_x\otimes L_{t-x})\) has all of its cohomology sheaves
direct sums of \(\mathcal{L}_\psi(\alpha x)\). This does not seem obvious a
priori, and we do not know a ``direct'' proof of it.

\sourcepara{2.10.12}
Here is a slight reformulation of the results of this section. Denote by
\(\operatorname{GalRep}(\mathbb{A}^1/k,\overline{\mathbb{Q}}_\ell)\) the
category of those continuous finite-dimensional
\(\overline{\mathbb{Q}}_\ell\)-representations \(V\) of
\(\operatorname{Gal}(k(x)^{\mathrm{sep}}/k(x))\), the Galois group of the
function field of \(\mathbb{A}^1/k\), with the following two properties:
\begin{enumerate}[label=(\arabic*)]
\item \(V\) is definable over a finite extension of \(\mathbb{Q}_\ell\);
\item \(V\) is unramified outside a finite set of places of \(k(x)\).
\end{enumerate}
If we denote by \(\eta\) the generic point of \(\mathbb{A}^1\), the functor
\[
  \begin{aligned}
  (\text{perverse middle extension sheaves on }\mathbb{A}^1)
    &\to \operatorname{GalRep}(\mathbb{A}^1/k,\overline{\mathbb{Q}}_\ell)\\
  N=j_*\mathcal{F}[1]&\mapsto \mathcal{F}_\eta=(N[-1])_\eta
  \end{aligned}
\]
is an equivalence of categories. Indeed, for each nonempty open \(U\) in
\(\mathbb{A}^1\), this functor induces an equivalence between the full
(by \hyperref[par:2.3.3.1]{(2.3.3.1)}) subcategory of its source,
\[
  (\text{perverse middle extensions on }\mathbb{A}^1\text{ which are lisse on }U),
\]
and the full subcategory of its target
\[
  (\text{objects in }\operatorname{GalRep}(\mathbb{A}^1/k,\overline{\mathbb{Q}}_\ell)
  \text{ which are unramified on }U).
\]
Using this fact, \hyperref[lem:2.10.3]{Lemma~2.10.3} and
\hyperref[thm:2.10.8]{Theorem~2.10.8} say that the composite functor
\[
  \begin{aligned}
  (\text{perverse sheaves on }\mathbb{A}^1\text{ with }\mathcal{P})
    &\to \operatorname{GalRep}(\mathbb{A}^1/k,\overline{\mathbb{Q}}_\ell)\\
  K&\mapsto (\operatorname{FT}(K)[-1])_\eta
  \end{aligned}
\]
is an equivalence of categories, under which middle additive convolution goes
over to tensor product, and under which the functor
\[
  K\mapsto \mathbb{D}_{-}K:=[x\mapsto -x]^*\mathbb{D}K
  =\mathbb{D}([x\mapsto -x]^*K)
\]
goes over into the functor \(V\mapsto V^\vee\), the contragredient
representation to \(V\).

\sourcepara{2.10.13}
Using this description, we can easily analyze the invertible objects. We say
that a perverse sheaf \(K\) on \(\mathbb{A}^1\) with \(\mathcal{P}\) is
invertible for the operation of middle convolution if there exists a perverse
\(L\) with \(\mathcal{P}\) such that \(K*_{\mathrm{mid}}+L\cong\delta_0\). If
such an \(L\) exists, it is unique, by the associativity of middle convolution.

\sourceheading[thm:2.10.14]{Theorem 2.10.14}
On \(\mathbb{A}^1\) over an algebraically closed field of characteristic
\(p>0\), a perverse sheaf \(K\) on \(\mathbb{A}^1\) with \(\mathcal{P}\) is
invertible for the operation of middle convolution if and only if
\(\operatorname{FT}(K)[-1]\) has generic rank one. In this case,
\(K*_{\mathrm{mid}}+\mathbb{D}_{-}K=\delta_0\).

\begin{proof}
Obvious by Fourier Transform, where it becomes the question of looking for
\(\otimes\)-invertible objects in \(\operatorname{GalRep}\): these are
obviously those of dimension one, and for these, the contragredient is the
\(\otimes\)-inverse.
\end{proof}

\sourceheading[cor:2.10.15]{Corollary 2.10.15}
On \(\mathbb{A}^1\) over an algebraically closed field \(k\) of characteristic
\(p>0\), the perverse sheaves \(K\) on \(\mathbb{A}^1\) with \(\mathcal{P}\)
which are both invertible for middle convolution and tame at \(\infty\) are
precisely the translated \(\delta\)-functions \(\delta_\alpha\) for \(\alpha\)
in \(k\), and the translated nontrivial Kummer sheaves
\(j_*\mathcal{L}_\chi(x-\alpha)[1]\) for \(\alpha\) in \(k\).

\begin{proof}
The objects listed are obviously invertible, and tame at \(\infty\). We must
show there are no more. First of all, any invertible object \(K\) is irreducible
as a perverse sheaf: indeed, from the Fourier Transform description, such a
\(K\) has no proper subobject in \(\mathcal{P}\), while, as already noted in
the proof of \hyperref[thm:2.9.8]{Theorem~2.9.8}, any irreducible subobject of
a perverse \(K\) with \(\mathcal{P}\) itself has \(\mathcal{P}\).

Among perverse irreducible \(K\)'s, the punctual ones are already on our list.
So we must look for irreducible middle extensions \(j_*\mathcal{F}[1]\) with
\(\mathcal{F}\) tame at \(\infty\), for which
\(\operatorname{FT}(j_*\mathcal{F})\) has generic rank one. For
\(\mathcal{F}\) tame at \(\infty\), and \(\alpha\neq 0\) in \(k\), the stalk
\(\operatorname{FT}(j_*\mathcal{F})_\alpha\) is
\[
  \mathrm{H}^1_c(\mathbb{A}^1,j_*\mathcal{F}\otimes\mathcal{L}_\psi(\alpha x));
\]
both the \(\mathrm{H}^2_c\) and \(\mathrm{H}^0\) vanish
(\(j_*\mathcal{F}\otimes\mathcal{L}_\psi(\alpha x)\) is totally wild at
\(\infty\), and has no nonzero punctual sections). So for \(\alpha\neq 0\), the
rank of the stalk \(\operatorname{FT}(j_*\mathcal{F})_\alpha\) is
\[
  \begin{aligned}
  -\chi_c(\mathbb{A}^1,j_*\mathcal{F}\otimes\mathcal{L}_\psi(\alpha x))
    &=
  -\operatorname{generic\ rank}(j_*\mathcal{F}\otimes\mathcal{L}_\psi(\alpha x))\\
  &\quad+\operatorname{Swan}_\infty(j_*\mathcal{F}\otimes\mathcal{L}_\psi(\alpha x))\\
  &\quad+\sum_{\text{finite sing }s}
    [\operatorname{Swan}_s(j_*\mathcal{F}\otimes\mathcal{L}_\psi(\alpha x))
      +\operatorname{drop}_s(j_*\mathcal{F}\otimes\mathcal{L}_\psi(\alpha x))].
  \end{aligned}
\]
Since \(\mathcal{F}\) is tame at \(\infty\), and
\(\mathcal{L}_\psi(\alpha x)\) has slope one at \(\infty\), the first two terms
cancel. Since \(\mathcal{L}_\psi(\alpha x)\) is lisse on \(\mathbb{A}^1\), at
each finite singularity \(s\) of \(\mathcal{F}\), the
\(\mathcal{L}_\psi(\alpha x)\) might as well be absent: we find
\[
  \operatorname{rank}\operatorname{FT}(j_*\mathcal{F})_\alpha
  =
  \sum_{\text{finite sing }s}
  [\operatorname{Swan}_s(j_*\mathcal{F})+\operatorname{drop}_s(j_*\mathcal{F})].
\]
Each of the terms
\([\operatorname{Swan}_s(j_*\mathcal{F})+\operatorname{drop}_s(j_*\mathcal{F})]\)
is strictly positive (the drop is nonzero), and each term where \(\mathcal{F}\)
is not tame is at least two. So if this rank is to be one, there is at
precisely one finite singularity, \(\mathcal{F}\) is tame there, and has a drop
of one. Translating the singularity to the origin, we get an \(\mathcal{F}\)
which is lisse and tame on \(\mathbb{G}_m\), irreducible, and nontrivially
ramified at zero. Since it is lisse, tame and irreducible, it must be rank one
(\(\pi_1^{\mathrm{tame}}(\mathbb{G}_m)\) is abelian), so a Kummer sheaf
\(\mathcal{L}_\chi\) on \(\mathbb{G}_m\). Since it is ramified at zero,
\(\chi\) is nontrivial.
\end{proof}

\section[Invertible objects on the affine line]{Invertible objects on
\texorpdfstring{\(\mathbb{A}^1\)}{A1} in characteristic zero}\label{sec:2.11}

\sourceheading[thm:2.11.1]{Theorem 2.11.1}
On \(\mathbb{A}^1\) over an algebraically closed field \(k\) of characteristic
zero, the perverse sheaves \(K\) on \(\mathbb{A}^1\) with \(\mathcal{P}\) which
are invertible for middle convolution are precisely the translated
\(\delta\)-functions \(\delta_\alpha\) for \(\alpha\) in \(k\), and the
translated nontrivial Kummer sheaves \(j_*\mathcal{L}_\chi(x-\alpha)[1]\) for
\(\alpha\) in \(k\).

\begin{proof}
The listed objects are visibly invertible. We must show there are no more. By
\hyperref[cor:2.9.9]{Corollary~2.9.9}, any invertible object \(K\) is perverse
irreducible. If \(K\) is punctual, it is already on our list. So suppose that
\(K\) is an irreducible middle extension \(j_*\mathcal{F}[1]\), with
\(\mathcal{F}\) nonconstant. Such a \(K\) lies in \(\mathcal{P}\).

Step 1. We show that for any irreducible middle extension
\(K=j_*\mathcal{F}[1]\), with \(K\) in \(\mathcal{P}\), there exists a
surjective map of perverse sheaves \(K*_{\mathrm{mid}}+\mathbb{D}_{-}(K)\twoheadrightarrow
\delta_0\). This is obvious in characteristic \(p\), by looking on the Fourier
Transform side. For this, we argue as follows. For any perverse sheaf \(N\) on
\(\mathbb{A}^1\), we have a short exact sequence of perverse sheaves
\[
  0\to\mathcal{H}^{-1}(N)[1]\to N\to\mathcal{H}^0(N)\to 0,
\]
in which the sheaf \(\mathcal{H}^0(N)\) is punctual. So for any \(\alpha\) in
\(\mathbb{A}^1\), we have a surjective map
\(\mathcal{H}^0(N)\twoheadrightarrow\mathcal{H}^0(N)_\alpha\otimes\delta_\alpha\),
so all in all \(N\twoheadrightarrow\mathcal{H}^0(N)_\alpha\otimes\delta_\alpha\).
Applying this to \(N=K*_{\mathrm{mid}}+\mathbb{D}_{-}(K)\), it suffices to show
that
\[
  \mathcal{H}^0(K*_{\mathrm{mid}}+\mathbb{D}_{-}(K))_0
\]
is one-dimensional. Because we are in characteristic zero, we can compute every
fibre of \(K*_{\mathrm{mid}}+\mathbb{D}_{-}(K)\) (by
\hyperref[cor:2.8.5]{Corollary~2.8.5}, (3) and (4)). Recall that
\(K=j_*\mathcal{F}[1]\), \(j:U\to\mathbb{A}^1\) a dense open set on which
\(\mathcal{F}\) is lisse, irreducible, and nonconstant. In terms of
\(k:\mathbb{A}^1\to\mathbb{P}^1\) the inclusion, we have
\[
  (K*_{\mathrm{mid}}+\mathbb{D}_{-}(K))_0
  =
  \mathrm{R}\Gamma(\mathbb{P}^1,k_*(j_*\mathcal{F}\otimes j_*(\mathcal{F}^{\vee})))[2].
\]
Thus
\[
  \begin{aligned}
  \mathcal{H}^0(K*_{\mathrm{mid}}+\mathbb{D}_{-}(K))_0
    &=
  \mathrm{H}^2(\mathbb{P}^1,k_*(j_*\mathcal{F}\otimes j_*(\mathcal{F}^{\vee})))\\
    &=
  \mathrm{H}^2_c(U,\mathcal{F}\otimes\mathcal{F}^{\vee})
  =\overline{\mathbb{Q}}_\ell(-1),
  \end{aligned}
\]
the last equality because \(\mathcal{F}\) is irreducible. In fact, for such
\(K\), one can show that
\(\mathcal{H}^0(K*_{\mathrm{mid}}+\mathbb{D}_{-}(K))_\alpha=0\) for
\(\alpha\neq 0\), provided we are in characteristic zero. In other words,
\(\mathcal{H}^0(K*_{\mathrm{mid}}+\mathbb{D}_{-}(K))\) is the required
\(\delta_0\) quotient.

Step 2. We show that if \(K*_{\mathrm{mid}}+L=\delta_0=L*_{\mathrm{mid}}+K\),
then \(L\) is \(\mathbb{D}_{-}(K)\). Indeed, applying the end-exact functor
\(X\mapsto L*_{\mathrm{mid}}+X\) to the surjection
\(K*_{\mathrm{mid}}+\mathbb{D}_{-}(K)\twoheadrightarrow\delta_0\), produced in
Step 1, we get
\[
  L*_{\mathrm{mid}}+K*_{\mathrm{mid}}+\mathbb{D}_{-}(K)
  \to L*_{\mathrm{mid}}+\delta_0,
  \quad\text{i.e.,}\quad
  \mathbb{D}_{-}(K)\to L.
\]
Since \(K\) is irreducible, so is \(\mathbb{D}_{-}(K)\), and hence
\(\mathbb{D}_{-}(K)=L\).

Step 3. We show that if
\(K*_{\mathrm{mid}}+\mathbb{D}_{-}(K)=\delta_0\) with \(K\) an irreducible
middle extension \(j_*\mathcal{F}[1]\), then \(\mathcal{F}\) is a translate of
a nontrivial Kummer sheaf \(\mathcal{L}_\chi\). Because we are in characteristic
zero, we have, by \hyperref[cor:2.8.5]{Corollary~2.8.5}, (3) and (4), a
short exact sequence of perverse sheaves on \(\mathbb{A}^1\)
\[
  0\to(\text{constant sheaf})[1]\to K*_!+\mathbb{D}_{-}(K)
  \to K*_{\mathrm{mid}}+\mathbb{D}_{-}(K)\to 0,
\]
which, if \(K*_{\mathrm{mid}}+\mathbb{D}_{-}(K)=\delta_0\), reads
\[
  0\to(\text{constant sheaf})[1]\to K*_!+\mathbb{D}_{-}(K)\to\delta_0\to 0.
\]
We exploit this last exact sequence by computing the difference in the virtual
ranks of stalks: for any \(t\neq 0\), we have
\[
  1=\operatorname{rank}(K*_!+\mathbb{D}_{-}(K))_0
    -\operatorname{rank}(K*_!+\mathbb{D}_{-}(K))_t,
\]
i.e.,
\[
  1=
  \chi_c(\mathbb{A}^1,j_*(\mathcal{F})\otimes j_*(\mathcal{F}^{\vee}))
  -
  \chi_c(\mathbb{A}^1,j_*(\mathcal{F})\otimes[x\mapsto x+t]^*j_*(\mathcal{F}^{\vee})).
\]
Denote by \(n\) the rank of \(\mathcal{F}\) (also that of
\(\mathcal{F}^{\vee}\)), and denote by \(S\) the set of finite singularities of
\(\mathcal{F}\) (also that of \(\mathcal{F}^{\vee}\)). At any \(s\) in \(S\),
both \(j_*(\mathcal{F})\) and \(j_*(\mathcal{F}^{\vee})\) have the same
dimensional fibre, say of dimension \(r_s\) (the same because \(r_s\) is the
common number of unipotent Jordan blocks in the local monodromies of
\(\mathcal{F}\) and of \(\mathcal{F}^{\vee}\) at the point \(s\), cf.
\hyperref[sec:3.1]{\S3.1}). Since we are in characteristic zero, the
Euler--Poincaré formula gives
\[
  \chi_c(\mathbb{A}^1,j_*(\mathcal{F})\otimes j_*(\mathcal{F}^{\vee}))
  =
  n^2-\sum_{s\in S}(n^2-r_s^2).
\]
For \(t\) sufficiently general, the two finite sets \(S\) and \(S+t\) are
disjoint. So for such \(t\),
\[
  \begin{aligned}
  &\chi_c(\mathbb{A}^1,j_*(\mathcal{F})\otimes
    [x\mapsto x+t]^*j_*(\mathcal{F}^{\vee}))\\
  &\quad=
  n^2-\sum_{s\in S}(n^2-r_sn)-\sum_{s\in S+t}(n^2-nr_s)\\
  &\quad=
  n^2-2\sum_{s\in S}(n^2-r_sn)\\
  &\quad=
  n^2-\sum_{s\in S}(2n)(n-r_s)\\
  &\quad=
  n^2-\sum_{s\in S}((n+r_s)+(n-r_s))(n-r_s)\\
  &\quad=
  n^2-\sum_{s\in S}(n^2-r_s^2)-\sum_{s\in S}(n-r_s)^2.
  \end{aligned}
\]
Subtracting, we find
\[
  1=\sum_{s\in S}(n-r_s)^2.
\]
Since each term \((n-r_s)^2\) is a strictly positive integer, we conclude that
\(S\) consists of a single point. Translating that point to the origin,
\(\mathcal{F}\) becomes a lisse sheaf on \(\mathbb{G}_m\) which is irreducible
and nonconstant, so necessarily a nontrivial Kummer sheaf
\(\mathcal{L}_\chi\).
\end{proof}

\section[Musings on middle-invertible objects]{Musings on
\texorpdfstring{\(*_{\mathrm{mid}}\)}{*mid}-invertible objects in
\texorpdfstring{\(\mathcal{P}\)}{P} in the
\texorpdfstring{\(\mathbb{G}_m\)}{Gm} case}\label{sec:2.12}

\sourcepara{2.12.1}
In this section, we work on \(\mathbb{G}_m\) over an algebraically closed field
\(k\) of characteristic \(\ne\ell\).

\sourceheading[lem:2.12.2]{Lemma 2.12.2}
On \(\mathbb{G}_m\) over an algebraically closed field \(k\) of characteristic
\(\ne\ell\), let \(K\) in
\(\mathrm{D}^b_c(\mathbb{G}_m,\overline{\mathbb{Q}}_\ell)\). The following
conditions are equivalent:
\begin{enumerate}[label=(\arabic*)]
\item \(\mathrm{R}\Gamma_c(\mathbb{G}_m,K\otimes\mathcal{L}_\chi)=0\) for
some Kummer sheaf \(\mathcal{L}_\chi\).
\item \(\mathrm{R}\Gamma(\mathbb{G}_m,K\otimes\mathcal{L}_\chi)=0\) for some
Kummer sheaf \(\mathcal{L}_\chi\).
\item[(3a)] each of the cohomology sheaves \(\mathcal{H}^i(K)\) is a
successive extension of Kummer sheaves \(\mathcal{L}_\chi\).
\item[(3b)] each of the perverse cohomology sheaves
\({}^p\mathcal{H}^i(K)\) is a successive extension of Kummer sheaves
\(\mathcal{L}_\chi[1]\).
\item \(\mathrm{R}\Gamma_c(\mathbb{G}_m,K\otimes\mathcal{L}_\chi)=0\) for all
but at most finitely many Kummer sheaves \(\mathcal{L}_\chi\).
\item \(\mathrm{R}\Gamma(\mathbb{G}_m,K\otimes\mathcal{L}_\chi)=0\) for all
but at most finitely many Kummer sheaves \(\mathcal{L}_\chi\).
\end{enumerate}
\begin{proof}
By \hyperref[lem:2.3.2.1]{Lemma~2.3.2.1}, (3a) \(\Leftrightarrow\) (3b). The
implications (3a) \(\Rightarrow\) (4) and (3a) \(\Rightarrow\) (5) are
obvious, because
\(\mathrm{R}\Gamma_c(\mathbb{G}_m,\mathcal{L}_\chi)
=\mathrm{R}\Gamma(\mathbb{G}_m,\mathcal{L}_\chi)=0\) for any nontrivial
Kummer sheaf. The implications (4) \(\Rightarrow\) (1) and (5)
\(\Rightarrow\) (2) are trivial. We will prove below that (2) \(\Rightarrow\)
(3a). Admitting this, we have (2) \(\Leftrightarrow\) (3a) \(\Leftrightarrow\)
(3b) \(\Leftrightarrow\) (5). As noted in
\hyperref[par:2.3.1.1]{(2.3.1.1)}, we have
\[
  {}^p\mathcal{H}^i(\mathbb{D}K):=
  \mathbb{D}({}^p\mathcal{H}^{-i}(K)).
\]
Therefore (3b) holds for \(K\) if and only if (3b) holds for \(\mathbb{D}K\).
Therefore the equivalences (2) \(\Leftrightarrow\) (3b) \(\Leftrightarrow\)
(5) hold for \(\mathbb{D}K\), which by duality means that the equivalences
(1) \(\Leftrightarrow\) (3b) \(\Leftrightarrow\) (4) hold for \(K\).

It remains to show that (2) \(\Rightarrow\) (3a). For this, we argue as
follows, cf. \cite[2.5.3]{Ka-ACT}. Replacing \(K\) by
\(K\otimes\mathcal{L}_\chi\), we may assume that
\(\mathrm{R}\Gamma(\mathbb{G}_m,K)=0\). Consider the spectral sequence
\[
  E_2^{p,q}=\mathrm{H}^p(\mathbb{G}_m,\mathcal{H}^q(K))
  \Rightarrow \mathrm{H}^{p+q}(\mathbb{G}_m,K).
\]
It has \(E_2^{p,q}=0\) unless \(p\) is \(0\) or \(1\) (cohomological dimension
of an affine curve), so degenerates at \(E_2\). Therefore for each \(\chi\),
we have
\[
  \mathrm{R}\Gamma(\mathbb{G}_m,K)=0
  \Leftrightarrow
  \mathrm{R}\Gamma(\mathbb{G}_m,\mathcal{H}^i(K))=0
  \text{ for all }i.
\]
So to prove (2) \(\Rightarrow\) (3), it suffices to do so in the case when
\(K\) is a single sheaf \(\mathcal{F}\). In this case, we are reduced to the
following sublemma.
\end{proof}

\sourceheading[sublem:2.12.3]{Sublemma 2.12.3}
On \(\mathbb{G}_m\) over an algebraically closed field \(k\) of characteristic
\(\ne\ell\), let \(\mathcal{F}\) be a constructible
\(\overline{\mathbb{Q}}_\ell\)-sheaf such that
\(\mathrm{H}^*(\mathbb{G}_m,\mathcal{F})=0\). Then \(\mathcal{F}\) is lisse on
\(\mathbb{G}_m\) and everywhere tame, hence a successive extension of Kummer
sheaves \(\mathcal{L}_\chi\).
\begin{proof}
Let \(j:U\to\mathbb{G}_m\) be a dense open set where \(\mathcal{F}\) is lisse.
Consider the canonical map \(\mathcal{F}\to j_*j^*\mathcal{F}\), whose
(punctual) kernel we denote \(\mathcal{F}_{\mathrm{pct}}\). The inclusion
\(\mathcal{F}_{\mathrm{pct}}\subset\mathcal{F}\) induces an inclusion
\[
  \mathrm{H}^0(\mathbb{G}_m,\mathcal{F}_{\mathrm{pct}})
  \subset \mathrm{H}^0(\mathbb{G}_m,\mathcal{F})=0,
\]
whence \(\mathrm{H}^0(\mathbb{G}_m,\mathcal{F}_{\mathrm{pct}})=0\), and
consequently \(\mathcal{F}_{\mathrm{pct}}=0\). Thus \(\mathcal{F}\) has no
nonzero punctual sections. For such an \(\mathcal{F}\), the Euler--Poincaré
formula
\[
\begin{aligned}
  \chi(\mathbb{G}_m,\mathcal{F})
  &=-\operatorname{Swan}_0(\mathcal{F})
    -\operatorname{Swan}_{\infty}(\mathcal{F})\\
  &\quad-\sum_{x\text{ in }\mathbb{G}_m}
    [\operatorname{drop}_x(\mathcal{F})
    +\operatorname{Swan}_x(\mathcal{F})]
\end{aligned}
\]
is a sum of terms which are each nonpositive. Since
\(\chi(\mathbb{G}_m,\mathcal{F})=0\), we see that \(\mathcal{F}\) is lisse on
\(\mathbb{G}_m\), and everywhere tame. Because
\(\pi_1^{\mathrm{tame}}(\mathbb{G}_m)\) is abelian, any such \(\mathcal{F}\)
is a successive extension of Kummer sheaves \(\mathcal{L}_\chi\), these being
the characters of \(\pi_1^{\mathrm{tame}}(\mathbb{G}_m)\).
\end{proof}

For ease of later reference, we record here the well-known
\sourceheading[lem:2.12.3.1]{Lemma 2.12.3.1}
On \(\mathbb{G}_m\) over an algebraically closed field \(k\) of characteristic
\(\ne\ell\), let \(K\) in
\(\mathrm{D}^b_c(\mathbb{G}_m,\overline{\mathbb{Q}}_\ell)\) be perverse. Then
\begin{enumerate}[label=(\arabic*)]
\item \(\chi(\mathbb{G}_m,K)\geq0\),
\item \(\chi(\mathbb{G}_m,K)=0\) if and only if \(K\) is a successive
extension of Kummer sheaves \(\mathcal{L}_\chi[1]\).
\end{enumerate}
\begin{proof}
For (1), use the short exact sequence
\[
  0\to\mathcal{H}^{-1}(K)[1]\to K\to\mathcal{H}^0(K)\to0
\]
to reduce to the case when \(K\) is either punctual, in which case the
assertion is obvious, or is of the form \(\mathcal{F}[1]\), with
\(\mathcal{F}\) a sheaf on \(\mathbb{G}_m\) with no nonzero punctual sections.
In this case the Euler--Poincaré formula shows, as above, that
\(\chi(\mathbb{G}_m,\mathcal{F})\leq0\), i.e.,
\(\chi(\mathbb{G}_m,\mathcal{F}[1])\geq0\).

For (2), the vanishing of
\[
  \chi(\mathbb{G}_m,K)
  =
  \chi(\mathbb{G}_m,\mathcal{H}^{-1}(K)[1])
  +\chi(\mathbb{G}_m,\mathcal{H}^0(K))
\]
together with the non-negativity of each summand shows that
\(\mathcal{H}^0(K)=0\), and that
\(\chi(\mathbb{G}_m,\mathcal{H}^{-1}(K)[1])=0\). But
\(\mathcal{H}^{-1}(K)\) has no nonzero punctual sections, so again as in the
proof of \hyperref[sublem:2.12.3]{Sublemma~2.12.3} the Euler--Poincaré formula
shows that \(\mathcal{H}^{-1}(K)\) is a successive extension of Kummer sheaves
\(\mathcal{L}_\chi\).
\end{proof}

\sourceheading[lem:2.12.4]{Lemma 2.12.4}
On \(\mathbb{G}_m\) over an algebraically closed field \(k\) of characteristic
\(\ne\ell\), let \(K\) in
\(\mathrm{D}^b_c(\mathbb{G}_m,\overline{\mathbb{Q}}_\ell)\). For all but at
most finitely many Kummer sheaves \(\mathcal{L}_\chi\), the ``forget
supports'' map
\[
  \mathrm{R}\Gamma_c(\mathbb{G}_m,K\otimes\mathcal{L}_\chi)
  \to
  \mathrm{R}\Gamma(\mathbb{G}_m,K\otimes\mathcal{L}_\chi)
\]
is an isomorphism.
\begin{proof}
By the spectral sequences
\[
\begin{aligned}
  E_2^{p,q}
  &=\mathrm{H}^p_c(\mathbb{G}_m,\mathcal{H}^q(K)\otimes\mathcal{L}_\chi)
  \Rightarrow
  \mathrm{H}^{p+q}_c(\mathbb{G}_m,K\otimes\mathcal{L}_\chi),\\
  E_2^{p,q}
  &=\mathrm{H}^p(\mathbb{G}_m,\mathcal{H}^q(K)\otimes\mathcal{L}_\chi)
  \Rightarrow
  \mathrm{H}^{p+q}(\mathbb{G}_m,K\otimes\mathcal{L}_\chi),
\end{aligned}
\]
we reduce immediately to the case where \(\mathcal{F}\) is a single sheaf. In
that case, we need only avoid the \(\chi\) whose inverses occur either in
\(\mathcal{F}(0)\) or \(\mathcal{F}(\infty)\), the \(I(0)\)- and
\(I(\infty)\)-representations attached \(\mathcal{F}\).
\end{proof}

\sourceheading[lem:2.12.5]{Lemma 2.12.5 (Gabber--Loeser)}
On \(\mathbb{G}_m\) over an algebraically closed field \(k\) of characteristic
\(\ne\ell\), let \(K\) and \(L\) be two objects in
\(\mathrm{D}^b_c(\mathbb{G}_m,\overline{\mathbb{Q}}_\ell)\). Then the mapping
cone object formed from the ``forget supports'' map between \(!\) and \(*\)
multiplicative convolutions,
\[
  K*_{!\times}L\to K*_{*\times}L,
\]
has all of its ordinary and perverse cohomology sheaves successive extensions
of Kummer sheaves.
\begin{proof}
(Compare the proof of \cite[page~28]{Ga-Loe}.) Choose a Kummer sheaf
\(\mathcal{L}_\chi\) such that the ``forget supports'' map is an isomorphism
\[
  \mathrm{R}\Gamma_c(\mathbb{G}_m,N\otimes\mathcal{L}_\chi)
  \cong
  \mathrm{R}\Gamma(\mathbb{G}_m,N\otimes\mathcal{L}_\chi)
\]
for \(N\) any of the three objects \(K\), \(L\), \(K*_{!\times}L\). This is
possible by the preceding lemma. For this \(\chi\), we claim the ``forget
supports'' map
\[
  \mathrm{R}\Gamma(\mathbb{G}_m,(K*_{!\times}L)\otimes\mathcal{L}_\chi)
  \to
  \mathrm{R}\Gamma(\mathbb{G}_m,(K*_{*\times}L)\otimes\mathcal{L}_\chi)
\]
is an isomorphism. Since \(\chi\) has been chosen so that
\[
  \mathrm{R}\Gamma_c(\mathbb{G}_m,(K*_{!\times}L)\otimes\mathcal{L}_\chi)
  \cong
  \mathrm{R}\Gamma(\mathbb{G}_m,(K*_{!\times}L)\otimes\mathcal{L}_\chi),
\]
it suffices to show that the ``forget supports twice'' map
\[
  \mathrm{R}\Gamma_c(\mathbb{G}_m,(K*_{!\times}L)\otimes\mathcal{L}_\chi)
  \to
  \mathrm{R}\Gamma(\mathbb{G}_m,(K*_{*\times}L)\otimes\mathcal{L}_\chi)
\]
is an isomorphism. Using the definition of multiplicative convolution, the
multiplicativity of \(\mathcal{L}_\chi\), the Leray spectral sequence for both
\(!\) and \(*\) direct image, and the Kunneth formula for both
\(\mathrm{R}\Gamma_c\) and \(\mathrm{R}\Gamma\), this map becomes the map
``forget supports'' \(\otimes\) ``forget supports''
\[
  \mathrm{R}\Gamma_c(\mathbb{G}_m,K\otimes\mathcal{L}_\chi)\otimes
  \mathrm{R}\Gamma_c(\mathbb{G}_m,L\otimes\mathcal{L}_\chi)
  \to
  \mathrm{R}\Gamma(\mathbb{G}_m,K\otimes\mathcal{L}_\chi)\otimes
  \mathrm{R}\Gamma(\mathbb{G}_m,L\otimes\mathcal{L}_\chi),
\]
which is an isomorphism by our choice of \(\chi\). The lemma now follows from
\hyperref[lem:2.12.2]{Lemma~2.12.2}, applied to \(K:=\) the mapping cone of
\(K*_{!\times}L\to K*_{*\times}L\).
\end{proof}

\sourceheading[cor:2.12.6]{Corollary 2.12.6}
On \(\mathbb{G}_m\) over an algebraically closed field \(k\) of characteristic
\(\ne\ell\), let \(K\) and \(L\) be two objects in
\(\mathrm{D}^b_c(\mathbb{G}_m,\overline{\mathbb{Q}}_\ell)\) which are
perverse. Suppose that \(K\) satisfies \(\mathcal{P}\) for multiplicative
convolution. Then the kernel and cokernel of the ``forget supports'' map
between the perverse sheaves
\[
  K*_{!\times}L\to K*_{*\times}L
\]
are perverse sheaves which are successive extensions of Kummer sheaves
\(\mathcal{L}_\chi[1]\).
\begin{proof}
The kernel and cokernel are the perverse cohomology sheaves of the mapping
cone.
\end{proof}

\sourceheading[cor:2.12.7]{Corollary 2.12.7}
On \(\mathbb{G}_m\) over an algebraically closed field \(k\) of characteristic
\(\ne\ell\), let \(K\) and \(L\) be two objects in
\(\mathrm{D}^b_c(\mathbb{G}_m,\overline{\mathbb{Q}}_\ell)\) which are
perverse. Suppose that \(K\) satisfies \(\mathcal{P}\) for multiplicative
convolution. Then we have the product formula for middle multiplicative
convolution
\[
  \chi(\mathbb{G}_m,K*_{\mathrm{mid}\times}L)
  =
  \chi(\mathbb{G}_m,K)\chi(\mathbb{G}_m,L).
\]
\begin{proof}
We have a short exact sequence of perverse sheaves
\[
  0\to\ker\to K*_{\mathrm{mid}\times}L\to K*_{!\times}L\to0
\]
on \(\mathbb{G}_m\). We know that \(\chi(\mathbb{G}_m,\ker)=0\) because
\(\ker\) is a successive extension of \(\mathcal{L}_\chi[1]\)'s. Thus we find
\[
  \chi(\mathbb{G}_m,K*_{\mathrm{mid}\times}L)
  =
  \chi(\mathbb{G}_m,K*_{!\times}L).
\]
By the \(!\) Kunneth formula we know that
\[
  \chi_c(\mathbb{G}_m,K*_{!\times}L)
  =
  \chi_c(\mathbb{G}_m,K)\chi_c(\mathbb{G}_m,L).
\]
Finally, we know that \(\chi_c(\mathbb{G}_m,N)=\chi(\mathbb{G}_m,N)\) for any
derived category object \(N\).
\end{proof}

\sourceheading[cor:2.12.8]{Corollary 2.12.8}
On \(\mathbb{G}_m\) over an algebraically closed field \(k\) of characteristic
\(\ne\ell\), let \(K\) and \(L\) be two objects in
\(\mathrm{D}^b_c(\mathbb{G}_m,\overline{\mathbb{Q}}_\ell)\) which are
perverse and which both satisfy \(\mathcal{P}\) for multiplicative
convolution. Suppose that \(K*_{\mathrm{mid}\times}L=\delta_1\), i.e.,
suppose that \(K\) in \(\mathcal{P}\) is invertible for multiplicative middle
convolution, with inverse \(L\). Then
\[
  \chi(\mathbb{G}_m,K)=\chi(\mathbb{G}_m,L)=1.
\]
\begin{proof}
Indeed, \(\chi(\mathbb{G}_m,\delta_1)=1\), and
\(\chi(\mathbb{G}_m,K)\) is a nonnegative integer for any perverse \(K\) on
\(\mathbb{G}_m\), by \hyperref[lem:2.12.3.1]{Lemma~2.12.3.1}.
\end{proof}

\sourceheading[thm:2.12.9]{Theorem 2.12.9}
On \(\mathbb{G}_m\) over an algebraically closed field \(k\) of characteristic
\(\ne\ell\), let \(K\) and \(L\) in
\(\mathrm{D}^b_c(\mathbb{G}_m,\overline{\mathbb{Q}}_\ell)\) be perverse and
satisfy \(\mathcal{P}\) for multiplicative convolution. Suppose that
\(K*_{\mathrm{mid}\times}L=\delta_1\).
\begin{enumerate}[label=(\arabic*)]
\item On \(\mathbb{G}_m\), the operators \(X\mapsto X*_{\mathrm{mid}\times}K\)
and \(X\mapsto X*_{\mathrm{mid}\times}L\) on \(\mathcal{P}\) are automorphisms
of \(\mathcal{P}\), which are inverses of each other.
\item The operator \(X\mapsto X*_{\mathrm{mid}\times}K\) induces an
automorphism of the perverse irreducible objects in \(\mathcal{P}\), with
inverse the operator \(X\mapsto X*_{\mathrm{mid}\times}L\).
\item \(L\) and \(K\) are perverse irreducible on \(\mathbb{G}_m\).
\end{enumerate}
\begin{proof}
This is entirely analogous to the additive case, cf.
\hyperref[thm:2.9.7]{Theorem~2.9.7} and
\hyperref[cor:2.9.9]{Corollary~2.9.9}.
\end{proof}

\sourceheading[thm:2.12.10]{Theorem 2.12.10}
On \(\mathbb{G}_m\) over an algebraically closed field \(k\) of characteristic
\(\ne\ell\), let \(K\) in
\(\mathrm{D}^b_c(\mathbb{G}_m,\overline{\mathbb{Q}}_\ell)\) be perverse. The
following conditions are equivalent.
\begin{enumerate}[label=(\arabic*)]
\item \(K\) lies in \(\mathcal{P}\) and \(\chi(\mathbb{G}_m,K)=1\).
\item \(K\) is perverse irreducible and \(\chi(\mathbb{G}_m,K)=1\).
\item \(K\) is an irreducible hypergeometric (cf.
\cite[3.5.4 and 8.5.3]{Ka-ESDE}).
\item \(K\) lies in \(\mathcal{P}\), and there exists a geometric isomorphism
\[
  K*_{\mathrm{mid}\times}\mathbb{D}(\operatorname{inv}^*K)\cong\delta_1.
\]
\item \(K\) lies in \(\mathcal{P}\), and is invertible for multiplicative
middle convolution.
\end{enumerate}
\begin{proof}
For (1) \(\Rightarrow\) (2), let \(L\subset K\) be nonzero perverse
irreducible. If \(L=K\), stop. If \(K/L\) is nonzero, we get a contradiction
as follows. Since \(K\) lies in \(\mathcal{P}\), \(L\) is not any
\(\mathcal{L}_\chi[1]\), so \(\chi(\mathbb{G}_m,L)>0\), by
\hyperref[lem:2.12.3.1]{Lemma~2.12.3.1}, (2). But by
\hyperref[lem:2.12.3.1]{Lemma~2.12.3.1}, (1),
\[
  \chi(\mathbb{G}_m,\text{any perverse})\geq0,
\]
so from the exact sequence
\[
  0\to L\to K\to K/L\to0
\]
we get \(\chi(\mathbb{G}_m,L)=1\), and \(\chi(\mathbb{G}_m,K/L)=0\). Thus
\(K/L\) is a successive extension of \(\mathcal{L}_\chi[1]\)'s, and so \(K\)
has an \(\mathcal{L}_\chi[1]\) quotient. But this is impossible, because \(K\)
lies in \(\mathcal{P}\).

That (2) \(\Rightarrow\) (3) is proven in \cite[3.5.4 and 8.5.3]{Ka-ESDE}.
That (4) \(\Rightarrow\) (5) is trivial, and (5) \(\Rightarrow\) (1) is
proven in \hyperref[cor:2.12.8]{Corollary~2.12.8} above.

It remains only to prove (3) \(\Rightarrow\) (1). Suppose first that we are in
characteristic \(p>0\). By the structure theorem \cite[8.5.3]{Ka-ESDE} for
hypergeometrics, we know that every irreducible hypergeometric is either a
\(\delta\) or is \(\mathcal{F}[1]\) for \(\mathcal{F}\) a hypergeometric sheaf
\(\mathcal{H}_{\lambda}(!,\psi,\chi\)'s, \(\rho\)'s) with disjoint \(\chi\)'s
and \(\rho\)'s. Such a hypergeometric is a successive
\(*_{\mathrm{mid}\times}\) convolution (because with disjoint \(\chi\)'s and
\(\rho\)'s, each \(*_{!\times}\cong *_{*\times}\) in the constructive
definition of hypergeometric objects cf. \cite[8.4.2 (5)]{Ka-ESDE}) of
sheaves \(\mathcal{L}_\chi\otimes\mathcal{L}_\psi[1]\), their multiplicative
translates (this only changes the \(\psi\)) and their multiplicative inverses.
So we are reduced to showing that
\[
  K*_{\mathrm{mid}\times}\mathbb{D}(\operatorname{inv}^*K)\cong\delta_1
\]
for \(K=\mathcal{L}_\chi\otimes\mathcal{L}_\psi[1]\). Now for all the
convolutions on \(\mathbb{G}_m\), we have
\[
  (\mathcal{L}_\chi\otimes A)*_{\mathrm{any},\times}
  (\mathcal{L}_\chi\otimes B)
  =
  \mathcal{L}_\chi\otimes(A*_{\mathrm{any},\times}B).
\]
So it remains to show that
\(K*_{\mathrm{mid}\times}\mathbb{D}(\operatorname{inv}^*K)\cong\delta_1\) for
\(K=\mathcal{L}_\psi[1]\). This will be proven in
\hyperref[cor:2.13.4]{Corollary~2.13.4} below.

If we are in characteristic zero, then by classification every irreducible
hypergeometric is either a \(\delta\) or is \(\mathcal{F}[1]\) for
\(\mathcal{F}\) a hypergeometric sheaf
\(\mathcal{H}_{\lambda}(\chi\)'s, \(\rho\)'s) of type \((n,n)\) with disjoint
\(\chi\)'s and \(\rho\)'s. Such an \(\mathcal{F}[1]\) is a successive
multiplicative middle convolution (with disjoint \(\chi\)'s and \(\rho\)'s,
each \(*_{!\times}\cong *_{*\times}\)) of hypergeometrics of type \((1,1)\),
\(\mathcal{H}_{\lambda}(\chi,\rho)[1]\) with \(\chi\ne\rho\) (compare
\cite[5.3.1]{Ka-ESDE} for the \(\mathcal{D}\)-module analogue). So we are
reduced to the case when \(K\) is
\(j_*\mathcal{L}_{\chi(x)(\rho/\chi)(\lambda-x)}[1]\). Just as in
characteristic \(p\), we may replace \(K\) by
\(K\otimes\mathcal{L}_{\bar{\chi}}\), and multiplicatively translate
\(\lambda\) to \(1\). This reduces us to treating universally the case when
\(K=j_*\mathcal{L}_{\chi(1-x)}[1]\), with \(\chi\) nontrivial. This will be
proven in \hyperref[cor:2.13.3]{Corollary~2.13.3} below.
\end{proof}

\section[Relations between middle convolution on
\texorpdfstring{\(\mathbb{A}^1\)}{A1} and
\texorpdfstring{\(\mathbb{G}_m\)}{Gm}]{Interlude: surprising relations between
\texorpdfstring{\(*_{\mathrm{mid}}\)}{*mid} on
\texorpdfstring{\(\mathbb{A}^1\)}{A1} and on
\texorpdfstring{\(\mathbb{G}_m\)}{Gm}}\label{sec:2.13}

\sourceheading[lem:2.13.1]{Lemma 2.13.1}
On \(\mathbb{G}_m\) over an algebraically closed field \(k\) of characteristic
\(\ne\ell\), let \(K\) in
\(\mathrm{D}^b_c(\mathbb{G}_m,\overline{\mathbb{Q}}_\ell)\) be perverse, and
let \(\mathcal{L}_\chi\) be any nontrivial Kummer sheaf on \(\mathbb{G}_m\).
Let \(j:\mathbb{G}_m\to\mathbb{A}^1\) be the inclusion, and denote
\(j_*\mathcal{L}_\chi\) on \(\mathbb{A}^1\) simply as \(\mathcal{L}_\chi\). We
have
\[
\begin{aligned}
  K*_{!\times}\mathcal{L}_\chi(x-1)[1]
  &=j^*(j_!(K\otimes\mathcal{L}_{\bar{\chi}})*_{!+}\mathcal{L}_\chi[1]),\\
  K*_{*\times}\mathcal{L}_\chi(x-1)[1]
  &=j^*(\mathrm{R}j_*(K\otimes\mathcal{L}_{\bar{\chi}})
    *_{*+}\mathcal{L}_\chi[1]),\\
  K*_{\mathrm{mid}\times}\mathcal{L}_\chi(x-1)[1]
  &=j^*(j_{!*}(K\otimes\mathcal{L}_{\bar{\chi}})
    *_{\mathrm{mid}+}\mathcal{L}_\chi[1]).
\end{aligned}
\]
\begin{proof}
Let us start with \(K\) on \(\mathbb{G}_m\) any derived category object. For
\(\chi\) nontrivial, we have, for
\(\pi:\mathbb{G}_m\times\mathbb{G}_m\to\mathbb{G}_m\) the map
\((x,t)\mapsto t\), and \(j:\mathbb{G}_m\to\mathbb{A}^1\) the inclusion,
\[
\begin{aligned}
  K*_{!\times}\mathcal{L}_\chi(x-1)
  &:=
  \mathrm{R}\pi_!(K_x\otimes\mathcal{L}_\chi((t/x)-1))\\
  &=
  \mathrm{R}\pi_!((K\otimes\mathcal{L}_{\bar{\chi}})_x
    \otimes\mathcal{L}_\chi(t-x))\\
  &=
  j^*(j_!(K\otimes\mathcal{L}_{\bar{\chi}})*_{!+}\mathcal{L}_\chi).
\end{aligned}
\]
Dualizing this, we get
\[
  \mathbb{D}K*_{*\times}\mathcal{L}_{\bar{\chi}}(x-1)
  =
  j^*(\mathrm{R}j_*(\mathbb{D}K\otimes\mathcal{L}_\chi)
  *_{*+}\mathcal{L}_{\bar{\chi}}).
\]
Replacing \(K\) by \(\mathbb{D}K\) and \(\chi\) by \(\bar{\chi}\), this gives
\[
  K*_{*\times}\mathcal{L}_\chi(x-1)
  =
  j^*(\mathrm{R}j_*(K\otimes\mathcal{L}_{\bar{\chi}})
  *_{*+}\mathcal{L}_\chi).
\]
Suppose now that \(K\) is perverse on \(\mathbb{G}_m\). Taking images, the
right side becomes the pullback of the image of the composite
\[
\begin{aligned}
  j_!(K\otimes\mathcal{L}_{\bar{\chi}})*_{!+}\mathcal{L}_\chi[1]
  &\to
  j_{!*}(K\otimes\mathcal{L}_{\bar{\chi}})
  *_{\mathrm{mid}+}\mathcal{L}_\chi[1]\\
  &\to
  \mathrm{R}j_*(K\otimes\mathcal{L}_{\bar{\chi}})
  *_{*+}\mathcal{L}_\chi[1],
\end{aligned}
\]
namely
\(j_{!*}(K\otimes\mathcal{L}_{\bar{\chi}})
*_{\mathrm{mid}+}\mathcal{L}_\chi[1]\).
\end{proof}

\sourceheading[lem:2.13.2]{Lemma 2.13.2}
On \(\mathbb{G}_m\) over an algebraically closed field \(k\) of characteristic
\(p\ne\ell\), \(p>0\), let \(K\) in
\(\mathrm{D}^b_c(\mathbb{G}_m,\overline{\mathbb{Q}}_\ell)\) be perverse. Then
for \(j:\mathbb{G}_m\to\mathbb{A}^1\) the inclusion, we have
\[
\begin{aligned}
  \operatorname{inv}^*K*_{!\times}j^*\mathcal{L}_\psi[1]
  &=j^*\operatorname{FT}(j_!K),\\
  \operatorname{inv}^*K*_{*\times}j^*\mathcal{L}_\psi[1]
  &=j^*\operatorname{FT}(\mathrm{R}j_*K),\\
  \operatorname{inv}^*K*_{\mathrm{mid}\times}j^*\mathcal{L}_\psi[1]
  &=j^*\operatorname{FT}(j_{!*}K).
\end{aligned}
\]
\begin{proof}
The first assertion results formally from the definitions, cf.
\cite[page~264]{Ka-ESDE}, or \cite[8.6.1]{Ka-GKM}, and the second is the dual of
the first. Since \(j^*\mathcal{L}_\psi[1]\) on \(\mathbb{G}_m\) is in
\(\mathcal{P}\), being perverse irreducible and not an
\(\mathcal{L}_\chi[1]\), all the objects in the first two assertions are
perverse. The third assertion is the image of the first in the second by the
``forget supports'' map, where on the right we think of the source as
\(j^*\operatorname{FT}_!(j_!K)\) and the target as
\(j^*\operatorname{FT}_*(\mathrm{R}j_*K)\).
\end{proof}

As a nice application of these last two results, we can now complete the proof
of \hyperref[thm:2.12.10]{Theorem~2.12.10}, (3) \(\Rightarrow\) (4).

\sourceheading[cor:2.13.3]{Corollary 2.13.3}
On \(\mathbb{G}_m\) over an algebraically closed field \(k\) of characteristic
\(\ne\ell\), let \(\mathcal{L}_\chi\) be any nontrivial Kummer sheaf. Then
\(K=\mathcal{L}_{\chi(x-1)}[1]\) satisfies
\[
  K*_{\mathrm{mid}\times}\mathbb{D}(\operatorname{inv}^*K)=\delta_1.
\]
\begin{proof}
For any perverse \(L\) on \(\mathbb{G}_m\), we have
\[
  L*_{\mathrm{mid}\times}\mathcal{L}_\chi(x-1)[1]
  =
  j^*(j_{!*}(L\otimes\mathcal{L}_{\bar{\chi}})
  *_{\mathrm{mid}+}\mathcal{L}_\chi[1]).
\]
We take
\[
\begin{aligned}
  L
  &:=
  \mathbb{D}(\operatorname{inv}^*\mathcal{L}_\chi(x-1)[1])\\
  &=
  \mathcal{L}_{\bar{\chi}}((1/x)-1)[1]
  =
  \mathcal{L}_{\bar{\chi}}((1-x)/x)[1].
\end{aligned}
\]
Then \(j_{!*}(L\otimes\mathcal{L}_{\bar{\chi}})
=\mathcal{L}_{\bar{\chi}}(1-x)[1]\), and the formula becomes
\[
  \mathbb{D}(\operatorname{inv}^*\mathcal{L}_\chi(x-1)[1])
  *_{\mathrm{mid}\times}\mathcal{L}_\chi(x-1)[1]
  =
  j^*(\mathcal{L}_{\bar{\chi}}(1-x)[1]
  *_{\mathrm{mid}+}\mathcal{L}_\chi[1]).
\]
This last object on \(\mathbb{A}^1\) is (geometrically) the additive
translation by \(1\) of
\(\mathcal{L}_{\bar{\chi}}[1]*_{\mathrm{mid}+}\mathcal{L}_\chi[1]\), which we
have already seen to be \(\delta_0\). So we obtain \(\delta_1\), as required.
\end{proof}

\sourceheading[cor:2.13.4]{Corollary 2.13.4}
On \(\mathbb{G}_m\) over an algebraically closed field \(k\) of characteristic
\(p\ne\ell\), \(p>0\), \(K=\mathcal{L}_\psi[1]\) satisfies
\[
  K*_{\mathrm{mid}\times}\mathbb{D}(\operatorname{inv}^*K)=\delta_1.
\]
\begin{proof}
For any perverse \(L\) on \(\mathbb{G}_m\), we have, by
\hyperref[lem:2.13.2]{Lemma~2.13.2} above,
\[
  \operatorname{inv}^*L*_{\mathrm{mid}\times}j^*\mathcal{L}_\psi[1]
  =
  j^*\operatorname{FT}(j_{!*}L).
\]
For \(L=j^*\mathcal{L}_{\bar{\psi}}[1]\) on \(\mathbb{G}_m\), we have
\(j_{!*}L=\mathcal{L}_{\bar{\psi}}[1]\) on \(\mathbb{A}^1\), and so
\(\operatorname{FT}(j_{!*}L)=\delta_1\). Since \(L=\mathbb{D}K\), the above
formula says precisely
\[
  \operatorname{inv}^*\mathbb{D}K*_{\mathrm{mid}\times}K=\delta_1
  \quad\text{for }K=j^*\mathcal{L}_\psi[1]\text{ on }\mathbb{G}_m.
\]
But \(\operatorname{inv}^*\mathbb{D}K=\mathbb{D}(\operatorname{inv}^*K)\), so
the commutativity of middle convolution gives
\[
  K*_{\mathrm{mid}\times}\mathbb{D}(\operatorname{inv}^*K)=\delta_1,
\]
as required.
\end{proof}

\section{Interpretive remark: Fourier--Bessel Transform}\label{sec:2.14}
We continue to work on \(\mathbb{G}_m\) over an algebraically closed field of
characteristic \(p>0\), \(p\ne\ell\). For any irreducible hypergeometric sheaf
\(\mathcal{H}\), and in particular for any Kloosterman sheaf \(\mathcal{K}\),
\(\mathcal{H}[1]\) lies in \(\mathcal{P}\) and is invertible for middle
multiplicative convolution. Thus \(*_{\mathrm{mid}}(\mathcal{H}[1])\) is an
automorphism of \(\mathcal{P}\), with inverse
\(*_{\mathrm{mid}}\mathbb{D}\operatorname{inv}^*(\mathcal{H}[1])\). In
particular, if we take \(\mathcal{H}\) to be the ``classical'' rank two
Kloosterman sheaf \(\operatorname{Kl}_2\), we find the \(\ell\)-adic analogue
of the classical Fourier--Bessel Transform.

In characteristic \(0\), we can only form such a transform with
hypergeometrics of type \((n,n)\). But already in the case \(n=2\), where we
have the Gauss hypergeometric function, we have a transform which should have
received a classical name, and some classical attention.

\section{Questions about the situation in several variables}\label{sec:2.15}

\sourceheading[lem:2.15.1]{Lemma 2.15.1}
On \(\mathbb{A}^n\) over an algebraically closed field of characteristic
\(p\ne\ell\), \(p>0\), a perverse \(K\) in
\(\mathrm{D}^b_c(\mathbb{A}^n,\overline{\mathbb{Q}}_\ell)\) has
\(\mathcal{P}_!\) if and only if \(\operatorname{FT}(K)\) is of the form
\(\mathcal{F}[n]\) for some sheaf \(\mathcal{F}\).
\begin{proof}
First of all, we already know that for \(M\) and \(K\) perverse, the dual of
\(K*_!M\) is semiperverse, so \(K\) is in \(\mathcal{P}_!\) if and only if
\(K*_!M\) is semiperverse for every perverse \(M\). On the
\(\operatorname{FT}\) side, we are asking when \(L:=\operatorname{FT}(K)\) has
the property that \(N\mapsto N\otimes L[-n]\) maps perverses to
semiperverses. As soon as \(L[-n]\) has some nonzero cohomology sheaf in
strictly positive degree, we can take \(N\) to be \(\delta_\alpha\) for some
point \(\alpha\) where the offending cohomology sheaf lives, and get a
nonsemiperverse answer. So the condition is necessary, and it is trivially
sufficient, since when it holds,
\(N\mapsto N\otimes L[-n]=N\otimes\mathcal{F}\) only decreases the supports of
the \(\mathcal{H}^{-i}\).
\end{proof}

\sourcepara{2.15.2}
As an application of this lemma, \(\delta\)'s on \(\mathbb{A}^n\) and external
products of perverse objects with \(\mathcal{P}_!\) on each \(\mathbb{A}^1\)
factor will satisfy \(\mathcal{P}_!\) on \(\mathbb{A}^n\). What about condition
\(\mathcal{P}\) on \(\mathbb{A}^n\)? On the Fourier Transform side, we are
looking for the sheaves \(\mathcal{F}\) on \(\mathbb{A}^n\) such that
\(\mathcal{F}[n]\) is perverse and such that
\(\mathbb{D}(\mathcal{F}[n])=\mathcal{G}[n]\) for some sheaf \(\mathcal{G}\).
What is the classification of such sheaves?

\sourcepara{2.15.3}
Since \(K\) on \(\mathbb{A}^n\) has \(\mathcal{P}\) iff both \(K\) and
\(\mathbb{D}K\) have \(\mathcal{P}_!\), external products of perverse objects
with \(\mathcal{P}\) on each \(\mathbb{A}^1\) factor will satisfy
\(\mathcal{P}\) on \(\mathbb{A}^n\). What is the intrinsic characterization of
\(\mathcal{P}\) on \(\mathbb{A}^n\)? What if any is the analogue of the
one-variable criterion
\[
  \operatorname{Hom}(\mathcal{L}_\psi[1],K)
  =
  \operatorname{Hom}(K,\mathcal{L}_\psi[1])=0
\]
which we had on \(\mathbb{A}^1\)? There are presumably lots of objects in
\(\mathcal{P}\) on \(\mathbb{A}^n\) other than external products of objects in
\(\mathcal{P}\) on the factors, and \(\delta\)'s. What are they? How can
analyze the situation without \(\operatorname{FT}\)? What happens in
characteristic zero?

\sourcepara{2.15.4}
Already on higher dimensional tori \((\mathbb{G}_m)^n\) we don't seem to know
even that external products of perverse objects with \(\mathcal{P}\) on each
\(\mathbb{G}_m\) factor will satisfy \(\mathcal{P}\) on \((\mathbb{G}_m)^n\).

\section{Questions about the situation on elliptic curves}\label{sec:2.16}

\sourcepara{2.16.1}
What about the situation on an elliptic curve \(E\) over an algebraically
closed field of characteristic \(\ne\ell\), where again we may speak of
\(\mathcal{P}\) (here \(\mathcal{P}_!\) and \(\mathcal{P}_*\) are both
\(\mathcal{P}\))? Any perverse irreducible whose isomorphism class, under
translation, has only a finite stabilizer, is automatically in
\(\mathcal{P}\). We know the stability under \(*\) trivially for
\(\mathcal{P}\) (just as we knew it trivially for \(\mathcal{P}_!\) in the
noncompact case).

\sourcepara{2.16.2}
If we are over \(\overline{\mathbb{F}}_p\), there is a kind of
\(\operatorname{FT}\) defined using the Lang torsors over finite subfields
\(k_0\) over which \(E\) is defined,
\[
  E\to E\text{ by }1-\operatorname{Frob}_{k_0},
\]
with structural group \(E(k_0)\), and pushing out by characters \(\xi\) of the
group \(E(k_0)\) to get rank one lisse sheaves \(\mathcal{L}_\xi\) on \(E\)
against which we can form
\[
  (K,\xi)\mapsto \mathrm{R}\Gamma(E,K\otimes\mathcal{L}_\xi).
\]
The Euler characteristic \(\chi(E,K\otimes\mathcal{L}_\xi)\) is independent of
\(\xi\). In any case, if \(K\) is say perverse irreducible, and not itself any
\(\mathcal{L}_\xi\), then
\(\mathrm{R}\Gamma(E,K\otimes\mathcal{L}_\xi)\) has only its middle cohomology
group possibly nonzero.

Just as in the \(\mathbb{A}^1\) and \(\mathbb{G}_m\) cases, we have

\sourceheading[lem:2.16.3]{Lemma 2.16.3}
On an elliptic curve \(E\) over an algebraically closed field of characteristic
\(\ne\ell\), the necessary and condition that a perverse \(K\) on \(E\) have
\(\mathcal{P}\) is that for each lisse rank one \(\mathcal{L}\) on \(E\)
(automatically translation invariant), we have
\[
  \operatorname{Hom}(\mathcal{L}[1],K)=\operatorname{Hom}(K,\mathcal{L}[1])=0.
\]

\sourcepara{2.16.4}
What are the invertible objects on an elliptic curve \(E\) over an
algebraically closed field of characteristic \(\ne\ell\)? Since
\[
  \chi(E,K*L)=\chi(E,K)\chi(E,L),
\]
only those objects in \(\mathcal{P}\) whose Euler characteristic \(\chi=1\)
can possibly be invertible. Notice that any object in \(\mathcal{P}\) with
\(\chi=1\) is irreducible (since by the Euler--Poincaré formula on \(E\),
\(\chi(E,K)\geq0\) for any perverse \(K\), with equality if and only if \(K\)
is \(\mathcal{F}[1]\) with \(\mathcal{F}\) lisse on \(E\), so a successive
extension of lisse \(\mathcal{L}\)'s of rank one, cf. below). Do we know any
non-\(\delta\) examples of invertible objects? Indeed, do we know any
non-\(\delta\) examples of irreducible perverse sheaves \(K\) on \(E\) with
\(\chi(E,K)=1\)?

\sourcepara{2.16.5}
On \(E\) the Euler--Poincaré formula for \(j_*\mathcal{F}[1]\) gives
\[
  \chi(E,j_*\mathcal{F}[1])
  =
  -\chi(E,j_*\mathcal{F})
  =
  \sum_{x\text{ in }E}(\operatorname{drop}_x(\mathcal{F})
  +\operatorname{swan}_x(\mathcal{F})).
\]
So if \(\chi(E,j_*\mathcal{F}[1])=1\), then there is only one point \(x\) with
a nonzero contribution. At that point, (which by translation we may as well
take to be the origin \(0_E\)) \(\mathcal{F}\)'s local monodromy is tame
(since \(\operatorname{swan}>0\) forces \(\operatorname{drop}>0\)) and a
pseudoreflection. In fact this pseudoreflection must be unipotent, since
\(\det(\mathcal{F})\) is lisse of rank one on \(E-\{0\}\), with tame local
monodromy at \(0\); because \(E\) and \(E-\{0\}\) have the same tame
\((\pi_1)^{\mathrm{ab}}\), \(\det(\mathcal{F})\) is also lisse at zero. In
particular, \(\chi(E,j_*\mathcal{F}[1])=1\) implies
\(\operatorname{rank}(\mathcal{F})\geq2\).

\sourcepara{2.16.6}
In characteristic zero, such an \(\mathcal{F}\) on \(E-\{0\}\) is a pair
\(A,B\) in \(\operatorname{GL}(n,\mathbb{C})\),
\(n:=\operatorname{rank}(\mathcal{F})\), with commutator \(\{A,B\}\) a
unipotent pseudoreflection. If such a \(j_*\mathcal{F}[1]\) is in
\(\mathcal{P}\), then, as noted above, \(\mathcal{F}\) must be irreducible on
\(E-\{0\}\).

\sourcepara{2.16.7}
This brings us to the following question, which we are at present unable to
answer, except by explicit calculation, in rank \(n=2\), where the answer is
negative. Given \(n\geq2\), do there exist elements \(A,B\) in
\(\operatorname{GL}(n,\mathbb{C})\) with commutator \(\{A,B\}\) a unipotent
pseudoreflection, such that the group \(\langle A,B\rangle\) generated by
\(A\) and \(B\) is an irreducible subgroup of
\(\operatorname{GL}(n,\mathbb{C})\)? (As we have seen above, this is
equivalent to the question of whether there exist perverse irreducibles \(K\)
on \(E\) with \(\chi(E,K)=1\) of the form \(\mathcal{F}[1]\) with
\(\mathcal{F}\) of generic rank \(n\).) We may scale both \(A\) and \(B\), and
require that \(A\) and \(B\) lie in \(\operatorname{SL}(n,\mathbb{C})\) if we
like.

\sourcepara{2.16.8}
Let us briefly explain the negative answer to the \(n=2\) case of this
question. Here the ``trick'' is that an element \(X\) in
\(\operatorname{SL}(2,\mathbb{C})\) is unipotent if and only if
\(\operatorname{trace}(X)=2\). So we first want to find two elements \(A\) and
\(B\) in \(\operatorname{SL}(2,\mathbb{C})\) whose commutator has trace \(2\),
then check to see if \(A\) and \(B\) generate an irreducible subgroup of
\(\operatorname{SL}(2,\mathbb{C})\). We may conjugate by
\(\operatorname{SL}(2,\mathbb{C})\) to reduce to the case when \(A\) is upper
triangular. If \(A\) has distinct eigenvalues, then by further conjugation we
may assume it is \(\operatorname{diag}(x,y)\), with \(x\ne y\), \(xy=1\). If
\(A\) has repeated eigenvalues, they must be both \(\pm1\), so either
\(A=\pm1\), or \(A=\pm\) (the standard upper unipotent).

\sourcepara{2.16.9}
If \(A\) is diagonal and nonscalar, we compute \(\{A,B\}\), set its trace
\(=2\), and find that \(B\) is either upper or lower triangular, so
\(\langle A,B\rangle\) can't be irreducible in this case. If \(A\) is scalar,
\(\langle A,B\rangle\) is not irreducible. If \(A\) is \(\pm\)(unipotent), we
find that \(B\) is upper triangular, and again \(\langle A,B\rangle\) is not
irreducible. So it seems that no such irreducibles exist!

\sourcepara{2.16.10}
This seems to put a damper on our elliptic hopes, and makes it likely, or at
least plausible, there are no \(\chi=1\) irreducibles of any rank \(>1\). It
would be interesting to clarify this point.

\sourcepara{2.16.11}
Another point one should clarify is this: the rigid objects
\(\mathcal{F}_{n,\xi}\) on \(E-\{0_E\}\) constructed in the last chapter give
perverse irreducibles \(j_*\mathcal{F}_{n,\xi}[1]\) on \(E\), which have
\(\chi(E,j_*\mathcal{F}_{n,\xi}[1])=n\), hence lie in \(\mathcal{P}\). What
happens when we convolve them with each other?

\section{Appendix 1: the basic lemma on end-exact functors}\label{sec:2.17}

\sourceheading[lem:2.17.1]{Lemma 2.17.1}
We are given two abelian categories \(\mathcal{A}\) and \(\mathcal{B}\), two
exact functors \(S,T\) from \(\mathcal{A}\) to \(\mathcal{B}\), and a morphism
of functors \(\varphi:S\to T\). Then the functor \(\operatorname{Im}(\varphi)\)
from \(\mathcal{A}\) to \(\mathcal{B}\) defined by
\[
  \operatorname{Im}(\varphi)(A):=
  \operatorname{Image}(\varphi_A:S(A)\to T(A))
\]
is end-exact, i.e., it carries injections to injections, and it carries
surjections to surjections.
\begin{proof}
Let us begin with a short exact sequence in \(\mathcal{A}\),
\[
  0\to A\to B\to C\to0.
\]
Applying the functors \(S\) and \(T\), we get a commutative diagram in
\(\mathcal{B}\) with exact rows:
\[
\begin{tikzcd}
0 \arrow[r] & SA \arrow[r] \arrow[d,"\varphi_A"'] &
SB \arrow[r] \arrow[d,"\varphi_B"'] &
SC \arrow[r] \arrow[d,"\varphi_C"'] & 0\\
0 \arrow[r] & TA \arrow[r] & TB \arrow[r] & TC \arrow[r] & 0.
\end{tikzcd}
\]
Applying the snake lemma, we get the six term exact sequence
\[
\begin{aligned}
0&\to\ker(\varphi_A)\to\ker(\varphi_B)\to\ker(\varphi_C)
\xrightarrow{\delta}\operatorname{coker}(\varphi_A)\\
&\to\operatorname{coker}(\varphi_B)\to
\operatorname{coker}(\varphi_C)\to0.
\end{aligned}
\]
We denote by \(\ker(\varphi_C)[\delta]\subset\ker(\varphi_C)\) the kernel of
the coboundary map \(\delta\), and extract the short exact sequence
\[
  0\to\ker(\varphi_A)\to\ker(\varphi_B)\to
  \ker(\varphi_C)[\delta]\to0.
\]
Then we have a commutative diagram in \(\mathcal{B}\)
\[
\begin{tikzcd}
0 \arrow[r] & \ker(\varphi_A) \arrow[r] \arrow[d,hook] &
\ker(\varphi_B) \arrow[r] \arrow[d,hook] &
\ker(\varphi_C)[\delta] \arrow[r] \arrow[d,hook] & 0\\
0 \arrow[r] & SA \arrow[r] & SB \arrow[r] & SC \arrow[r] & 0
\end{tikzcd}
\]
with exact rows and with all vertical arrows injective. Applying the snake to
this diagram gives a short exact sequence
\[
  0\to SA/\ker(\varphi_A)\to SB/\ker(\varphi_B)
  \to SC/\ker(\varphi_C)[\delta]\to0,
\]
which we may rewrite as
\[
  0\to\operatorname{Im}(\varphi_A)\to\operatorname{Im}(\varphi_B)
  \to SC/\ker(\varphi_C)[\delta]\to0.
\]
Thus the map \(\operatorname{Im}(\varphi_A)\to\operatorname{Im}(\varphi_B)\)
is injective. The third term \(SC/\ker(\varphi_C)[\delta]\) of this exact
sequence surjects onto
\(SC/\ker(\varphi_C)\cong\operatorname{Im}(\varphi_C)\), and hence the map
\(\operatorname{Im}(\varphi_B)\to\operatorname{Im}(\varphi_C)\) is surjective,
being the composition of the two surjections
\[
  \operatorname{Im}(\varphi_B)\twoheadrightarrow
  SC/\ker(\varphi_C)[\delta]\twoheadrightarrow
  SC/\ker(\varphi_C)\cong\operatorname{Im}(\varphi_C).
\]
\end{proof}

\section{Appendix 2: twisting representations by characters}\label{sec:2.18}

\sourcepara{2.18.1}
In analyzing the translation-invariant perverse irreducibles on
\(\mathbb{A}^1\) over an algebraically closed field \(k\) of characteristic
\(p>0\), cf. the proof of \hyperref[lem:2.6.13]{Lemma~2.6.13}, we made use of
the following lemma, with \(\Gamma\) the group \(I(\infty)\), \(K\) the field
\(\overline{\mathbb{Q}}_\ell\), and characters all the
\(\mathcal{L}_\psi(\beta x)\), \(\beta\) in \(k\).

\sourceheading[lem:2.18.2]{Lemma 2.18.2}
Let \(\Gamma\) be a group, \(K\) an algebraically closed field of
characteristic zero, and \(M\) a finite-dimensional \(K\)-representation of
\(\Gamma\). Suppose that for an infinity of one-dimensional \(K\)-valued
characters \(\chi\) of \(\Gamma\), there exists an isomorphism
\(M\cong M\otimes\chi\) (as \(K\)-representation of \(\Gamma\)). Then \(M=0\).

We will prove this in the following equivalent form.

\sourceheading[lem:2.18.2bis]{Lemma 2.18.2 bis}
Let \(\Gamma\) be a group, \(K\) an algebraically closed field \(K\) of
characteristic zero, and \(M\) a non-zero finite-dimensional
\(K\)-representation of \(\Gamma\). Denote by \(\operatorname{Stab}(M,\Gamma)\)
the subgroup of the character group \(\operatorname{Hom}(\Gamma,K^\times)\)
consisting of those characters \(\chi\) of \(\Gamma\) for which there exists an
isomorphism \(M\cong M\otimes\chi\) (as \(K\)-representation of \(\Gamma\)).
Then \(\operatorname{Stab}(M,\Gamma)\) is finite.
\begin{proof}
Step 1. We first reduce to the case where \(M\) is a faithful representation
of \(\Gamma\). Denote by \(\ker(M)\subset\Gamma\) the kernel of the
representation \(M\) of \(\Gamma\). Then for any \(\chi\) in
\(\operatorname{Stab}(M,\Gamma)\), \(\chi|\ker(M)\) lies in
\(\operatorname{Stab}(M|\ker(M),\ker(M))\), hence is trivial. Thus we find
\[
  \operatorname{Stab}(M,\Gamma)
  =
  \operatorname{Stab}(M,\Gamma/\ker(M)).
\]

Step 2. We henceforth suppose that \(M\) is a faithful representation, i.e.
that \(\Gamma\subset\operatorname{GL}(M)\). We denote by \(G\) the Zariski
closure of \(\Gamma\) in \(\operatorname{GL}(M)\). We claim that every
character \(\chi\) in \(\operatorname{Stab}(M,\Gamma)\) extends uniquely to an
algebro-geometric character \(\widetilde{\chi}\) in
\(\operatorname{Hom}(G,\mathbb{G}_m)\).

Indeed, given \(\chi\) in \(\operatorname{Stab}(M,\Gamma)\), there exists by
definition an element \(A\) in \(\operatorname{GL}(M)\) such that, for every
\(\gamma\) in \(\Gamma\), we have
\[
\begin{aligned}
  A\gamma A^{-1}&=\gamma\chi(\gamma),\\
  A\gamma A^{-1}\gamma^{-1}&=\chi(\gamma),
\end{aligned}
\]
and remember only
\[
  A\gamma A^{-1}\gamma^{-1}=\text{scalar}.
\]
For fixed \(A\) in \(\operatorname{GL}(M)\), the map of
\(\operatorname{GL}(M)\) to itself defined by
\[
  B\mapsto ABA^{-1}B^{-1}
\]
is an algebro-geometric morphism. In the target \(\operatorname{GL}(M)\), the
set of scalar matrices is Zariski closed. So for fixed \(A\), the set of
\(B\) in \(\operatorname{GL}(M)\) for which
\[
  ABA^{-1}B^{-1}=\text{scalar}
\]
is Zariski closed. It is also a group, as one sees in rewriting this equation
as
\[
  ABA^{-1}=(\text{scalar})B.
\]
Thus the set of such \(B\)'s is a Zariski closed subgroup of
\(\operatorname{GL}(M)\), which contains \(\Gamma\). Therefore it contains
\(G\). Thus for every \(g\) in \(G\), there exists a scalar, say
\(\alpha(g)\), such that
\[
  \alpha(g):=AgA^{-1}g^{-1}.
\]
Clearly the function \(g\mapsto\alpha(g)\) is an algebro-geometric morphism to
\(\mathbb{G}_m\). This function is multiplicative in \(g\), as one sees by
rewriting the defining equation as
\[
  AgA^{-1}=\alpha(g)g.
\]
For if \(g\) and \(h\) are elements of \(G\), we have
\[
  \alpha(hg)hg=AhgA^{-1}=AhA^{-1}AgA^{-1}
  =\alpha(h)h\alpha(g)g=\alpha(h)\alpha(g)hg,
\]
and multiplying both sides by \((hg)^{-1}\) gives
\(\alpha(hg)=\alpha(h)\alpha(g)\). Thus \(g\mapsto\alpha(g)\) is an
algebro-geometric character of \(G\) which agrees with \(\chi\) on
\(\Gamma\). Since \(\Gamma\) is Zariski dense in \(G\), there is at most one
such, which we denote \(\widetilde{\chi}\). Thus
\[
  \widetilde{\chi}(g):=\alpha(g).
\]

Step 3. For every \(\chi\) in \(\operatorname{Stab}(M,\Gamma)\),
\(\widetilde{\chi}\) lies in \(\operatorname{Stab}(M,G)\), as is clear from
the equation
\[
  AgA^{-1}=\alpha(g)g:=\widetilde{\chi}(g)g.
\]
Thus we have
\[
  \operatorname{Stab}(M,G)\cong\operatorname{Stab}(M,\Gamma),
\]
the isomorphism that of restriction to \(\Gamma\).

Step 4. For any \(\widetilde{\chi}\) in \(\operatorname{Stab}(M,G)\), by
taking determinants in the equation
\[
  AgA^{-1}=\widetilde{\chi}(g)g,
\]
we see that \(\widetilde{\chi}(g)^{\dim(M)}=1\). Thus
\(\widetilde{\chi}\) has finite order, and hence is a character of the finite
group \(G/G^0\) of components of \(G\). Therefore
\(\operatorname{Stab}(M,G)\) lies in the finite group
\(\operatorname{Hom}(G/G^0,K^\times)\), hence is finite. Therefore the
isomorphic group \(\operatorname{Stab}(M,\Gamma)\) is finite, as required.
\end{proof}
